\section{Algèbres formellement lisses et anneaux de Cohen}
\label{section:0.19-fr}
\label{0.19-fr}

\setcounter{subsection}{-1}
\subsection{Introduction}
\label{subsection:0.19.0-fr}
\label{0.19.0-fr}

\oldpage[0\textsubscript{IV-1}]{69}

\begin{env}[19.0.1]
\label{0.19.0.1-fr}
Au chapitre IV, nous introduirons et étudierons entre autres une classe
importante de morphismes de préschémas, les morphismes \emph{lisses}\footnote{Aussi
appelés morphismes \emph{simples} dans certains travaux récents (en s'inspirant
de la terminologie classique de «~point simple~»)~; cette terminologie prête
cependant à confusion, notamment en théorie des groupes algébriques.}. Une de
leurs propriétés fondamentales (et qui, jointe à une condition de finitude, les
caractérise) est une propriété de \emph{relèvement des morphismes}~: si
$f:X\to Y$ est un morphisme lisse, $g:Y'\to Y$ un morphisme, alors pour tout
morphisme $h:Y''\to Y'$ faisant de $Y''$ un préschéma «~peu différent~» de
$Y'$, \emph{tout $Y$-morphisme $Y''\to X$ se factorise} en
$Y''\xrightarrow{h}Y'\to X$. De façon précise, bornons-nous au cas où
$Y=\operatorname{Spec}(A)$, $X=\operatorname{Spec}(B)$ sont \emph{affines}~;
$B$ est alors appelée une $A$-algèbre \emph{formellement lisse} si, pour toute
$A$-algèbre $C$ et tout idéal \emph{nilpotent} $\mathfrak J$ de $C$, tout
$A$-homomorphisme $B\to C/\mathfrak J$ \emph{se relève} en
$B\to C\to C/\mathfrak J$. En d'autres termes, l'application
\[
  \operatorname{Hom}_A(B,C)\longrightarrow
  \operatorname{Hom}_A(B,C/\mathfrak J)
\]
est \emph{surjective}. Dans beaucoup d'applications, $X$ apparaîtra comme objet
représentant un \emph{foncteur contravariant représentable} $Y'\mapsto F(Y')$
de la catégorie des $Y$-préschémas dans celle des ensembles, de sorte
($\mathbf{0}_{\mathrm{III}}$, 8.1.8) que l'on aura
$F(Y')=\operatorname{Hom}_Y(Y',X)$. Dans le cas affine, si l'on pose
$F^0(C)=F(\operatorname{Spec}(C))$, la vérification du fait que, sous les
conditions ci-dessus, $F^0(C)\to F^0(C/\mathfrak J)$ est \emph{surjectif}
(qui peut se faire même \emph{sans savoir} que $F$ est représentable) établira
que $f$ est lisse, pourvu que la condition additionnelle de finitude soit
remplie.
\end{env}

\begin{env}[19.0.2]
\label{0.19.0.2-fr}
Pour les \emph{algèbres topologiques} sur un anneau topologique $A$, il y a une
notion analogue d'algèbre \emph{formellement lisse} que nous ne préciserons pas
ici (cf.~19.3.1). L'étude de ces notions est faite d'abord d'un point de vue
élémentaire au §~19, puis, à l'aide des propriétés des différentielles qui
seront développées aux §§~20 et 21, des théorèmes plus délicats seront prouvés
au §~22. Résumons ici les principaux résultats sur les algèbres formellement
lisses~:

\medskip\noindent\textbf{I.} --- Soient $A$, $B$ deux anneaux locaux
noethériens, $\varphi:A\to B$ un homomorphisme local, $k$ le corps résiduel de
$A$, et soit $B_0=B\otimes_A k$~; on munit $A$, $B$ et $B_0$ de leurs
topologies préadiques et $k$ de la topologie discrète. Alors, pour que $B$ soit
une $A$-algèbre formellement lisse, il faut et il suffit que $B$ soit un
$A$-module \emph{plat} et que $B_0$ soit une $k$-algèbre \emph{formellement
lisse} (19.7.1). Ce théorème ramène donc la lissité formelle, pour les anneaux
locaux noethériens, à la même question pour les anneaux locaux noethériens qui
sont des algèbres sur un \emph{corps}.

\medskip\noindent\textbf{II.} --- Soient $k$ un corps, $A$ un anneau local
noethérien qui est une $k$-algèbre. Pour que $A$ soit formellement lisse, il
faut et il suffit que $A$ soit \emph{géométriquement régulier} sur $k$,
c'est-à-dire que pour toute extension \emph{finie} $k'$ de $k$, l'anneau
semi-local $A\otimes_k k'$ soit

\oldpage[0\textsubscript{IV-1}]{70}

\noindent régulier (22.5.8)~; le \emph{complété} $\widehat A$ de $A$ est alors
un anneau isomorphe à un anneau de séries formelles
$K[[T_1,\ldots,T_n]]$ (19.6.5). En outre, la structure de $k$-algèbre de $A$,
lorsque $A$ est supposé être \emph{complet} et formellement lisse sur $k$, est
entièrement déterminée par le corps résiduel $K$ de $A$ et la \emph{dimension}
de $A$~; celle-ci peut d'ailleurs être arbitraire à condition de satisfaire à
l'inégalité
\[
  \dim(A)\geq \operatorname{rg}_k(\Gamma_{K/k}),
\]
où $\Gamma_{K/k}$ est le «~module d'imperfection~» de $K$ (22.2.6).

En particulier, pour qu'une \emph{extension} $K$ de $k$ soit formellement
lisse, il faut et il suffit que $K$ soit une extension \emph{séparable} de $k$
(19.6.1).

\medskip\noindent\textbf{III.} --- Soient $A$ un anneau local noethérien,
$\mathfrak J$ un idéal de $A$ distinct de $A$, $A_0=A/\mathfrak J$, $B_0$ un
anneau local noethérien \emph{complet}, $A_0\to B_0$ un homomorphisme local
faisant de $B_0$ une $A_0$-algèbre formellement lisse. Alors il existe un
anneau local noethérien \emph{complet} $B$, un homomorphisme local $A\to B$
faisant de $B$ un $A$-module plat, et un $A_0$-\emph{isomorphisme}
\[
  B\otimes_A A_0\xrightarrow{\sim}B_0
\]
(de sorte que, par I), $B$ est une $A$-algèbre \emph{formellement lisse})~; en
outre $B$ est déterminé par ces propriétés à isomorphisme près (19.7.2). Ce
théorème contient en particulier les \emph{théorèmes de Cohen} sur la structure
des anneaux locaux noethériens \emph{complets} (19.8), qui joueront un rôle
important dans les §§~6 et 7 du chapitre IV.

\medskip\noindent\textbf{IV.} --- L'intérêt de l'étude des anneaux locaux
noethériens formellement lisses sur un autre provient de la caractérisation
«~ponctuelle~» suivante des morphismes lisses~: si $X$ et $Y$ sont des
préschémas localement noethériens, $f:X\to Y$ un morphisme localement de type
fini, alors, pour que $f$ soit lisse, il faut et il suffit que pour tout
$x\in X$, l'anneau $\mathcal O_x$ soit formellement lisse sur
$\mathcal O_{f(x)}$. En particulier, si $Y=\operatorname{Spec}(k)$, où $k$ est
un corps parfait, dire que $f$ est lisse équivaut à dire que $X$ est un
préschéma \emph{régulier}.

\medskip\noindent\textbf{V.} --- Enfin, nous verrons aux §§~20, 21 et 22 que
la notion d'algèbre formellement lisse s'introduit de façon naturelle dans la
théorie des \emph{différentielles de Kähler}, les deux théories s'éclairant
mutuellement.
\end{env}

\begin{env}[19.0.3]
\label{0.19.0.3-fr}
Dans tout ce paragraphe et les suivants, les anneaux et modules topologiques
seront supposés \emph{linéairement topologisés} ($\mathbf{0}_{\mathrm I}$,
7.1.1)~; les anneaux topologiques considérés seront supposés \emph{commutatifs},
sauf mention expresse du contraire. Rappelons que si $A$ et $B$ sont deux
anneaux topologiques, $\rho:A\to B$ un homomorphisme d'anneaux définissant sur
$B$ une structure de $A$-algèbre, on dit que $B$ est une \emph{$A$-algèbre
topologique} si $\rho$ est \emph{continu} pour les topologies envisagées.

Pour abréger, dans un anneau topologique $A$ (resp. un $A$-module topologique
$M$), nous dirons «~\emph{système fondamental d'idéaux} (resp.
\emph{sous-modules}) \emph{ouverts}~» au lieu de «~système fondamental de
voisinages de 0 formé d'idéaux (resp. sous-modules)~».

Étant donnés un anneau topologique $A$ et un $A$-module $M$, les ensembles
$\mathfrak J M$, où $\mathfrak J$ parcourt un système fondamental d'idéaux
ouverts, forment un système fondamental de sous-modules ouverts d'une topologie
sur $M$ faisant de $M$ un $A$-module topologique, et que l'on dit \emph{déduite}
de la topologie de $A$.

Soit $M$ un $A$-module topologique dont la topologie est \emph{moins fine} que
la topologie déduite de celle de $A$~; alors, si $N$ est un sous-module
\emph{ouvert} de $M$, le $A$-module \emph{discret} $M/N$ est \emph{annulé par
un idéal ouvert} de $A$, car par hypothèse il existe un tel idéal $\mathfrak R$
tel que $\mathfrak R.M\subset N$.

\oldpage[0\textsubscript{IV-1}]{71}

Si $M$ et $N$ sont deux $A$-modules topologiques dont les topologies sont toutes
deux déduites de celle de $A$, alors \emph{tout} $A$-homomorphisme $u:M\to N$
est \emph{continu}, car pour tout voisinage $V$ de 0 dans $N$, il existe par
hypothèse un idéal ouvert $\mathfrak J$ de $A$ tel que $\mathfrak JN\subset V$,
et l'on a donc $u(\mathfrak JM)\subset\mathfrak JN\subset V$.

Lorsque $B$ est une $A$-algèbre topologique, la topologie sur $B$ \emph{déduite}
de celle de $A$ est \emph{plus fine} que la topologie donnée, car pour tout
idéal ouvert $\mathfrak R$ de $B$, il y a par hypothèse un idéal ouvert
$\mathfrak J$ de $A$ tel que $\mathfrak JB\subset\mathfrak R$.
\end{env}

\subsection{Épimorphismes et monomorphismes formels}
\label{subsection:0.19.1-fr}

\begin{proposition}[19.1.1]
\label{0.19.1.1-fr}
Soient $A$ un anneau topologique, $M$, $N$ deux $A$-modules topologiques,
$(W_\lambda)$ un système fondamental de sous-modules ouverts dans $N$,
$u:M\to N$ un $A$-homomorphisme continu, $\widehat M$, $\widehat N$ les
séparés complétés de $M$ et $N$, $\widehat u:\widehat M\to\widehat N$ le
prolongement continu de $u$ aux séparés complétés.
\begin{enumerate}
  \item[(i)] Les conditions suivantes sont équivalentes~:
  \begin{enumerate}
    \item[a)] $u(M)$ est \emph{dense} dans $N$.
    \item[a$'$)] $\widehat u(\widehat M)$ est \emph{dense} dans $\widehat N$.
    \item[a$''$)] Pour tout $\lambda$, l'homomorphisme composé
    $M\xrightarrow{u}N\to N/W_\lambda$ est \emph{surjectif}.
  \end{enumerate}
  \item[(ii)] Les conditions suivantes sont équivalentes~:
  \begin{enumerate}
    \item[b)] L'image réciproque par $u$ de la topologie de $N$ est égale à la
    topologie de $M$.
    \item[b$'$)] $\widehat u$ est un \emph{isomorphisme} du $\widehat A$-module
    topologique $\widehat M$ sur le sous-$\widehat A$-module topologique
    $\widehat u(\widehat M)$ de $\widehat N$ (qui est nécessairement
    \emph{fermé}).
    \item[b$''$)] Les $u^{-1}(W_\lambda)$ forment un système fondamental de
    voisinages de 0 dans $M$.
  \end{enumerate}
\end{enumerate}
Cela résulte immédiatement de la définition des séparés complétés, et du fait
que $N/W_\lambda$ est discret.
\end{proposition}

\begin{env}[19.1.2]
\label{0.19.1.2-fr}
Lorsque les conditions équivalentes de (i) (resp. (ii)) dans (19.1.1) sont
remplies, on dit que $u$ est un \emph{épimorphisme formel} (resp. un
\emph{monomorphisme formel}). On dit que $u$ est un \emph{bimorphisme formel}
s'il est à la fois un monomorphisme formel et un épimorphisme formel~; il
revient au même, en vertu de (19.1.1), de dire que $\widehat u$ est un
\emph{isomorphisme} du $\widehat A$-module topologique $\widehat M$ sur le
$\widehat A$-module topologique $\widehat N$.
\end{env}

\begin{proposition}[19.1.3]
\label{0.19.1.3-fr}
Soient $A$ un anneau topologique, $M$, $N$ deux $A$-modules topologiques,
$u:M\to N$ un $A$-homomorphisme continu. On suppose qu'il existe deux systèmes
fondamentaux de sous-modules ouverts, $(V_\lambda)$, $(W_\lambda)$ dans $M$ et
$N$ respectivement, ayant même ensemble d'indices et tels que
$u(V_\lambda)\subset W_\lambda$ pour tout $\lambda$~; soit
$u_\lambda:M/V_\lambda\to N/W_\lambda$ l'homomorphisme déduit de $u$ par
passage aux quotients. Alors~:
\begin{enumerate}
  \item[(i)] Pour que $u$ soit un épimorphisme formel, il faut et il suffit que,
  pour tout $\lambda$, $u_\lambda$ soit surjectif.
  \item[(ii)] Pour que $u$ soit un monomorphisme formel, il faut et il suffit
  que, pour tout $\lambda$, il existe un $\mu$ tel que
  $V_\mu\subset V_\lambda$ et que $\operatorname{Ker}(u_\mu)$ soit contenu
  dans le noyau $V_\lambda/V_\mu$ de l'application canonique
  $M/V_\mu\to M/V_\lambda$.
  \item[(iii)] Pour que $u$ soit un bimorphisme formel, il faut et il suffit
  que, pour tout $\lambda$, $u_\lambda$ soit surjectif et qu'il existe un
  $\mu$ tel que $V_\mu\subset V_\lambda$ et que l'homomorphisme canonique
  $M/V_\mu\to M/V_\lambda$

  \oldpage[0\textsubscript{IV-1}]{72}

  \noindent se factorise en
  $M/V_\mu\xrightarrow{u_\mu}N/W_\mu\to M/V_\lambda$ (où
  l'homomorphisme $N/W_\mu\to M/V_\lambda$ est nécessairement unique).
\end{enumerate}
\end{proposition}

\begin{proof}
L'assertion (i) est immédiate~; d'autre part,
$\operatorname{Ker}(u_\mu)=u^{-1}(W_\mu)/V_\mu$, et le noyau de l'application
canonique $M/V_\mu\to M/V_\lambda$ est $V_\lambda/V_\mu$~; la condition de
(ii) revient donc à dire que $u^{-1}(W_\mu)\subset V_\lambda$, et l'assertion
(ii) en découle aussitôt~; enfin, lorsque $u_\mu$ est surjectif, il revient au
même de dire que $\operatorname{Ker}(u_\mu)$ est contenu dans
$V_\lambda/V_\mu$, ou de dire que $M/V_\mu\to M/V_\lambda$ se factorise en
$v\circ u_\mu$, où $v$ est un homomorphisme
$N/W_\mu\to M/V_\lambda$, puisque $N/W_\mu$ s'identifie alors à
$(M/V_\mu)/\operatorname{Ker}(u_\mu)$.
\end{proof}

\begin{corollary}[19.1.4]
\label{0.19.1.4-fr}
Soient $A$ un anneau topologique, $M$, $N$ deux $A$-modules topologiques dont
les topologies sont déduites de celle de $A$, $u:M\to N$ un épimorphisme
formel. Soit $(\mathfrak J_\lambda)$ un système fondamental d'idéaux ouverts
dans $A$. Pour que $u$ soit un bimorphisme formel, il faut et il suffit que pour
tout $\lambda$, l'homomorphisme
$u_\lambda:M/\mathfrak J_\lambda M\to N/\mathfrak J_\lambda N$ déduit de $u$
par passage aux quotients soit bijectif.
\end{corollary}

\begin{proof}
On a en effet $u(\mathfrak J_\lambda M)\subset\mathfrak J_\lambda N$ et on
peut appliquer le critère (19.1.3, (iii))~; mais si l'on a une factorisation
\[
  M/\mathfrak J_\mu M\xrightarrow{u_\mu}
  N/\mathfrak J_\mu N\xrightarrow{v}M/\mathfrak J_\lambda M,
\]
l'image de $\mathfrak J_\lambda M/\mathfrak J_\mu M$ par $u_\mu$ est
$\mathfrak J_\lambda N/\mathfrak J_\mu N$, donc l'image de
$\mathfrak J_\lambda N/\mathfrak J_\mu N$ par $v$ est 0 et $v$ se factorise en
\[
  N/\mathfrak J_\mu N\longrightarrow N/\mathfrak J_\lambda N
  \xrightarrow{v'}M/\mathfrak J_\lambda M~;
\]
mais alors l'automorphisme identique de $M/\mathfrak J_\lambda M$ se factorise
en
\[
  M/\mathfrak J_\lambda M\xrightarrow{u_\lambda}
  N/\mathfrak J_\lambda N\xrightarrow{v'}M/\mathfrak J_\lambda M,
\]
ce qui montre que $u_\lambda$ est injectif, et comme on sait déjà qu'il est
surjectif, cela prouve le corollaire.
\end{proof}

\begin{proposition}[19.1.5]
\label{0.19.1.5-fr}
Soient $A$ un anneau topologique, $M$, $N$ deux $A$-modules topologiques,
$u:M\to N$ un $A$-homomorphisme continu. Les conditions suivantes sont
équivalentes~:
\begin{enumerate}
  \item[a)] Pour tout $A$-module topologique discret $E$ et tout
  $A$-homomorphisme continu $v:M\to E$, il existe un $A$-homomorphisme continu
  $w:N\to E$ tel que $v=w\circ u$.
  \item[b)] Pour tout sous-module ouvert $V'$ de $M$, il existe un sous-module
  ouvert $W''$ de $N$, un sous-module ouvert
  $V''\subset V'\cap u^{-1}(W'')$ et un homomorphisme
  $h:N/W''\to M/V'$ tels que l'homomorphisme canonique
  $M/V''\to M/V'$ se factorise en
  \[
    M/V''\xrightarrow{u''}N/W''\xrightarrow{h}M/V',
  \]
  où $u''$ est déduit de $u$ par passage aux quotients.
\end{enumerate}
\end{proposition}

\begin{proof}
Pour voir que a) entraîne b), il suffit d'appliquer a) à $E=M/V'$ et à
l'homomorphisme canonique $v:M\to M/V'$~; alors $W''=w^{-1}(0)$ est un
sous-module ouvert de $N$ et $V''=u^{-1}(W'')$ un sous-module ouvert de $M$
contenu dans $V'$~; ces sous-modules et l'homomorphisme
$h:N/W''\to M/V'$ déduit de $w$ par passage au quotient répondent à la
question. Inversement, pour montrer que b) entraîne a), on peut se limiter au
cas où $v$ est surjectif, en remplaçant $E$ par $v(M)$~; alors
$V'=\operatorname{Ker}(v)$ est un sous-module ouvert de $M$, donc $E$ est
isomorphe au $A$-module discret $M/V'$, et en appliquant b), on obtient

\oldpage[0\textsubscript{IV-1}]{73}

le $A$-homomorphisme continu $w$ cherché en prenant le composé
$N\to N/W''\xrightarrow{h}M/V'$, le diagramme
\[
  \xymatrix{
    M\ar[r]^u\ar[d] & N\ar[d]\\
    M/V''\ar[r]_{u''} & N/W''
  }
\]
étant commutatif.
\end{proof}

Lorsque les conditions équivalentes de (19.1.5) sont remplies, nous dirons que
$u$ est \emph{formellement inversible à gauche}~; comme la condition b) de
(19.1.5) implique que $u^{-1}(W'')\subset V'$, $u$ est alors un monomorphisme
formel. La terminologie est en outre justifiée par les corollaires suivants~:

\begin{corollary}[19.1.6]
\label{0.19.1.6-fr}
S'il existe un $\widehat A$-homomorphisme continu
$s:\widehat N\to\widehat M$ tel que
$s\circ\widehat u=1_{\widehat M}$, alors $u$ est formellement inversible à
gauche.
\end{corollary}

On notera que $\widehat u$ est alors un isomorphisme topologique de
$\widehat M$ sur un \emph{facteur direct topologique} de $\widehat N$. Pour
prouver (19.1.6), il suffit de noter qu'avec les notations de (19.1.5, a)),
$v:M\to E$ se prolonge en un $\widehat A$-homomorphisme continu
$\widehat v:\widehat M\to E$ puisque $E$ est discret~; soient
$j:M\to\widehat M$ et $j':N\to\widehat N$ les homomorphismes canoniques~;
alors, en posant $w=\widehat v\circ s\circ j'$, on a
$w\circ u=\widehat v\circ s\circ\widehat u\circ j=\widehat v\circ j=v$.

\begin{corollary}[19.1.7]
\label{0.19.1.7-fr}
Supposons que les topologies de $M$ et de $N$ soient déduites de celle de $A$,
et soit $(\mathfrak J_\lambda)$ un système fondamental d'idéaux ouverts dans
$A$. Pour que $u$ soit formellement inversible à gauche il faut et il suffit
que, pour tout $\lambda$, l'homomorphisme
$u_\lambda:M/\mathfrak J_\lambda M\to N/\mathfrak J_\lambda N$, déduit de
$u$ par passage aux quotients, soit inversible à gauche (autrement dit, soit
un isomorphisme de $M/\mathfrak J_\lambda M$ sur un facteur direct de
$N/\mathfrak J_\lambda N$).
\end{corollary}

En effet, la condition est suffisante, car, pour
$V'=\mathfrak J_\lambda M$ dans la condition b) de (19.1.5), on répond à la
question en prenant $W''=\mathfrak J_\lambda N$, $V''=V'$ et $h$ tel que
$h\circ u_\lambda$ soit l'identité dans $M/\mathfrak J_\lambda M$.
Inversement, si $u$ est formellement inversible à gauche, alors, pour tout
$\lambda$, il y a, en vertu de (19.1.5, b)), un $\mu$ tel que
$\mathfrak J_\mu\subset\mathfrak J_\lambda$ et un homomorphisme
$h:N/\mathfrak J_\mu N\to M/\mathfrak J_\lambda M$ tels que l'homomorphisme
canonique $M/\mathfrak J_\mu M\to M/\mathfrak J_\lambda M$ se factorise en
\[
  M/\mathfrak J_\mu M\xrightarrow{u_\mu}N/\mathfrak J_\mu N
  \xrightarrow{h}M/\mathfrak J_\lambda M~;
\]
mais comme
$h(\mathfrak J_\lambda(N/\mathfrak J_\mu N))
\subset\mathfrak J_\lambda(M/\mathfrak J_\lambda M)=0$, $h$ se factorise
canoniquement en
\[
  N/\mathfrak J_\mu N\longrightarrow N/\mathfrak J_\lambda N
  \xrightarrow{s}M/\mathfrak J_\lambda M,
\]
et il est immédiat que $s$ est inverse à gauche de $u_\lambda$.

\begin{proposition}[19.1.8]
\label{0.19.1.8-fr}
Soit $A$ un anneau topologique admettant un système fondamental dénombrable
décroissant $(\mathfrak J_n)$ d'idéaux ouverts. Soient $M$, $N$ deux
$A$-modules topologiques dont les topologies sont déduites de celle de $A$,
$u:M\to N$ un $A$-homomorphisme. Posons
$M_n=M/\mathfrak J_nM$, $N_n=N/\mathfrak J_nN$ et soit
$u_n:M_n\to N_n$ le $(A/\mathfrak J_n)$-homomorphisme déduit de $u$ par
passage aux quotients. Supposons que, pour tout $n$,
$P_n=\operatorname{Coker}(u_n)$ soit un $(A/\mathfrak J_n)$-module projectif
et que $M$ soit séparé et complet. Alors, pour que $u$ soit formellement
inversible à gauche, il faut et il suffit que $u$ soit inversible à gauche (et
$u$ est alors un isomorphisme topologique de $M$ sur un facteur direct
topologique de $N$).
\end{proposition}

\oldpage[0\textsubscript{IV-1}]{74}

\begin{proof}
En vertu de (19.1.7), on a des diagrammes commutatifs
\[
  \xymatrix{
    0\ar[r] & M_n\ar[r]^{u_n} & N_n\ar[r]^{p_n} & P_n\ar[r] & 0\\
    0\ar[r] & M_{n+1}\ar[u]^f\ar[r]_{u_{n+1}} &
      N_{n+1}\ar[u]^g\ar[r]_{p_{n+1}} & P_{n+1}\ar[u]^h\ar[r] & 0
  }
\]
où les lignes sont des suites exactes \emph{scindées}, autrement dit, il existe
pour chaque $n$ un homomorphisme $s_n:P_n\to N_n$ tel que
$p_n\circ s_n=1_{P_n}$. Nous allons montrer que l'on peut, par récurrence sur
$n$, définir un homomorphisme $s'_n:P_n\to N_n$ tel que
$p_n\circ s'_n=1_{P_n}$ et que les diagrammes
\[
  \xymatrix{
    P_n\ar[r]^{s'_n} & N_n\\
    P_{n+1}\ar[u]^h\ar[r]_{s'_{n+1}} & N_{n+1}\ar[u]^g
  }
\]
soient commutatifs. En effet, $g\circ s_{n+1}-s'_n\circ h$ est un
homomorphisme de $P_{n+1}$ dans
\[
  u_n(M_n)=u_{n+1}(M_{n+1})/
  \mathfrak J_nu_{n+1}(M_{n+1})~;
\]
comme $P_{n+1}$ est un $(A/\mathfrak J_{n+1})$-module \emph{projectif}, cet
homomorphisme se factorise en
\[
  P_{n+1}\xrightarrow{t_{n+1}}u_{n+1}(M_{n+1})
  \longrightarrow
  u_{n+1}(M_{n+1})/\mathfrak J_nu_{n+1}(M_{n+1}),
\]
d'où l'on conclut aussitôt que $s'_{n+1}=s_{n+1}-t_{n+1}$ répond à la
question. Cela étant, de la décomposition de $N_n$ en somme directe de
$u_n(M_n)$ et de $s'_n(P_n)$, on déduit aussitôt un homomorphisme
$w_n:N_n\to M_n$ inverse à gauche de $u_n$ et tel que les diagrammes
\[
  \xymatrix{
    N_n\ar[r]^{w_n} & M_n\\
    N_{n+1}\ar[u]^g\ar[r]_{w_{n+1}} & M_{n+1}\ar[u]^f
  }
\]
soient commutatifs. Le système projectif $(w_n)$ admet alors une limite
projective $\widehat w:\widehat N\to\widehat M=M$, d'où par composition avec
l'homomorphisme canonique $N\to\widehat N$, un homomorphisme $w:N\to M$ tel
que, pour tout $n$, l'endomorphisme $(w\circ u)_n$ de
$M_n=M/\mathfrak J_nM$ déduit de $w\circ u$ par passage aux quotients soit
l'identité~; cela entraîne que $w\circ u=1_M$ puisque $M$ est séparé et
complet.
\end{proof}

\begin{proposition}[19.1.9]
\label{0.19.1.9-fr}
Soient $A$ un anneau topologique préadmissible ($\mathbf{0}_{\mathrm I}$,
7.1.2), $\mathfrak Q$ un idéal de définition de $A$,
$(\mathfrak J_\lambda)$ un système fondamental d'idéaux ouverts de $A$.
Soient $M$, $N$ deux $A$-modules topologiques dont les topologies sont déduites
de celles de $A$, et tels que pour tout $\lambda$,
$N/\mathfrak J_\lambda N$ soit un $(A/\mathfrak J_\lambda)$-module projectif
(cf. (19.2.3)). Soit $u:M\to N$ un $A$-homomorphisme. Alors les conditions
suivantes sont équivalentes~:
\begin{enumerate}
  \item[a)] $u$ est formellement inversible à gauche.

  \oldpage[0\textsubscript{IV-1}]{75}

  \item[b)] L'homomorphisme
  $u_0:M/\mathfrak QM\to N/\mathfrak QN$ déduit de $u$ par passage aux
  quotients est inversible à gauche.
\end{enumerate}
\end{proposition}

\begin{proof}
On a vu (19.1.7) que la condition a) équivaut à dire que $u_\lambda$ est
inversible à gauche pour tout $\lambda$~; comme $\mathfrak Q$ est un idéal
ouvert, donc contient un $\mathfrak J_\lambda$, on en déduit aussitôt que
$u_0$ est inversible à gauche. Pour montrer inversement que b) entraîne a),
notons que pour tout $\lambda$, $\mathfrak Q/\mathfrak J_\lambda$ est par
hypothèse un idéal nilpotent de $A/\mathfrak J_\lambda$. Notre assertion
résultera de la proposition suivante.
\end{proof}

\begin{proposition}[19.1.10]
\label{0.19.1.10-fr}
Soient $A$ un anneau, $M$, $N$ deux $A$-modules, $N$ étant projectif,
$u:M\to N$ un $A$-homomorphisme. Soit $\mathfrak J$ un idéal de $A$~; on
suppose vérifiée l'une des conditions suivantes~:
\begin{enumerate}
  \item[(i)] $\mathfrak J$ est nilpotent.
  \item[(ii)] $\mathfrak J$ est contenu dans le radical de $A$ et $M$ est de
  type fini.
\end{enumerate}
Alors, pour que $u$ soit inversible à gauche, il faut et il suffit que
l'homomorphisme
$u_0:M/\mathfrak JM\to N/\mathfrak JN$ de $(A/\mathfrak J)$-modules, déduit
de $u$ par passage aux quotients, soit inversible à gauche.
\end{proposition}

\begin{proof}
La condition étant évidemment nécessaire, prouvons qu'elle est suffisante.
Soit $v_0$ un inverse à gauche de $u_0$~; l'homomorphisme composé
$N\to N/\mathfrak JN\xrightarrow{v_0}M/\mathfrak JM$ se factorise en
$N\xrightarrow{v}M\to M/\mathfrak JM$ puisque $N$ est projectif~; alors
$w=v\circ u$ est un endomorphisme de $M$ tel que l'endomorphisme $w_0$ de
$M/\mathfrak JM$ déduit par passage aux quotients soit l'identité, et il suffit
de prouver que $w$ est lui-même bijectif (car alors $w^{-1}\circ v$ sera un
inverse à gauche de $u$). Distinguons maintenant les deux cas.

(i) Pour tout $n$ on a le diagramme commutatif
\[
  \xymatrix{
    (\mathfrak J^n/\mathfrak J^{n+1})\otimes_{A/\mathfrak J}
      (M/\mathfrak JM)\ar[r]\ar[d]_{1\otimes\operatorname{gr}^0(w)} &
      \mathfrak J^nM/\mathfrak J^{n+1}M\ar[d]^{\operatorname{gr}^n(w)}\\
    (\mathfrak J^n/\mathfrak J^{n+1})\otimes_{A/\mathfrak J}
      (M/\mathfrak JM)\ar[r] & \mathfrak J^nM/\mathfrak J^{n+1}M
  }
\]
où les flèches horizontales sont \emph{surjectives}, et comme
$\operatorname{gr}^0(w)=w_0$ est l'identité, il en est de même de
$\operatorname{gr}^n(w)$ qui, \emph{a fortiori}, est bijectif. La filtration
$\mathfrak J$-préadique sur $M$ étant finie puisque $\mathfrak J$ est
nilpotent, on en conclut que $w$ est bijectif (Bourbaki, \emph{Alg. comm.},
chap.~III, §~2, n\textsuperscript{o}~8, cor.~3 du th.~1).

(ii) Il suffit de montrer que pour tout idéal maximal $\mathfrak m$ de $A$,
l'endomorphisme $w_{\mathfrak m}$ de $M_{\mathfrak m}$ est bijectif
(Bourbaki, \emph{Alg. comm.}, chap.~II, §~3, n\textsuperscript{o}~3, th.~1)
et comme $\mathfrak JA_{\mathfrak m}\subset\mathfrak mA_{\mathfrak m}$, et
$A_{\mathfrak m}/\mathfrak JA_{\mathfrak m}=(A/\mathfrak J)_{\mathfrak m}$,
on est ramené à prouver la proposition lorsque $A$ est un anneau local. En
outre, on peut supposer que $\mathfrak J$ est l'idéal maximal
$\mathfrak J_0$ de $A$, car si $u_0$ est inversible

\oldpage[0\textsubscript{IV-1}]{76}

à gauche, il en est de même de
$u_{00}:M/\mathfrak J_0M\to N/\mathfrak J_0M$ obtenu par tensorisation de
$u_0$ avec $1_{A/\mathfrak J_0}$, puisque l'on a
\[
  M/\mathfrak J_0M=(M/\mathfrak JM)\otimes_{A/\mathfrak J}(A/\mathfrak J_0).
\]
Supposons donc $\mathfrak J$ maximal, de sorte que $A/\mathfrak J$ est un
corps. Il suffit évidemment de montrer que $M$ est un $A$-module \emph{libre}
sous les conditions de l'énoncé~: en effet, $\det(w_0)$ est alors l'image
canonique de $\det(w)$ dans $A/\mathfrak J$, donc $\det(w)$ n'appartient pas à
l'idéal $\mathfrak J$ et est par suite inversible. Or le
$(A/\mathfrak J)$-espace vectoriel $M/\mathfrak JM$ étant libre de type fini,
il y a un $A$-module de type fini $L$ et un $A$-homomorphisme $f:L\to M$ tel
que l'homomorphisme $f_0:L/\mathfrak JL\to M/\mathfrak JM$ déduit de $f$ par
passage aux quotients soit bijectif. Comme $M$ est de type fini, on en conclut
tout d'abord que $f$ est surjectif par le lemme de Nakayama (Bourbaki,
\emph{Alg. comm.}, chap.~II, §~3, n\textsuperscript{o}~2, cor.~1 de la
prop.~4)~; en outre, si $g=u\circ f$, l'homomorphisme
$g_0:L/\mathfrak JL\to N/\mathfrak JN$ déduit par passage aux quotients est
inversible à gauche, et puisqu'ici $L$ est libre, la remarque du début prouve
que $g$ est lui-même inversible à gauche~; or cela entraîne évidemment que
$f$ est \emph{injectif}, ce qui achève la démonstration.
\end{proof}

Mentionnons en passant les corollaires suivants de (19.1.10)~:

\begin{corollary}[19.1.11]
\label{0.19.1.11-fr}
Soient $A$ un anneau local, $k$ son corps résiduel, $M$ un $A$-module de type
fini, $N$ un $A$-module projectif, $u:M\to N$ un homomorphisme. Pour que $u$
soit inversible à gauche, il faut et il suffit qu'il existe un système de
générateurs $(x_i)_{1\leq i\leq m}$ de $M$ tel que les images par
$u\otimes 1:M\otimes_Ak\to N\otimes_Ak$ des $x_i\otimes1$ soient linéairement
indépendantes dans $N\otimes_Ak$~; les $x_i$ forment alors une base de $M$.
\end{corollary}

La condition est évidemment nécessaire, car si $u$ est inversible à gauche,
$M$ est un $A$-module projectif de type fini, donc libre. Inversement, si la
condition est satisfaite, il est clair que les $x_i\otimes1$ forment une base
de $M\otimes_Ak$ et que $u\otimes1$ est inversible à gauche~; il suffit donc
d'appliquer (19.1.10) à l'idéal maximal $\mathfrak J$ de $A$.

\begin{corollary}[19.1.12]
\label{0.19.1.12-fr}
Soient $A$ un anneau, $M$ un $A$-module de type fini, $N$ un $A$-module
projectif, $u:M\to N$ un homomorphisme. Pour tout idéal premier
$\mathfrak p\in\operatorname{Spec}(A)$, les conditions suivantes sont
équivalentes~:
\begin{enumerate}
  \item[a)] L'homomorphisme
  $u_{\mathfrak p}:M_{\mathfrak p}\to N_{\mathfrak p}$ est inversible à
  gauche.
  \item[b)] L'homomorphisme
  $u\otimes1:M\otimes_Ak(\mathfrak p)\to N\otimes_Ak(\mathfrak p)$ est
  injectif (on rappelle que $k(\mathfrak p)$ désigne le corps résiduel de $A$
  en $\mathfrak p$).
  \item[c)] Il existe un système fini d'éléments $x_i\in M$
  $(1\leq i\leq m)$ tel que les images des $x_i$ dans $M_{\mathfrak p}$
  engendrent $M_{\mathfrak p}$, et un système de $m$ formes linéaires
  $y_i^*$ sur $N$ $(1\leq i\leq m)$ tel que
  $\det(\langle y_i^*,u(x_j)\rangle)\notin\mathfrak p$.
  \item[d)] Il existe $f\in A-\mathfrak p$ tel que l'homomorphisme
  $u_f:M_f\to N_f$ soit inversible à gauche.
\end{enumerate}
En outre, l'ensemble des $\mathfrak p\in\operatorname{Spec}(A)$ vérifiant ces
conditions est ouvert dans $\operatorname{Spec}(A)$.
\end{corollary}

La dernière condition est conséquence triviale de d). Comme $N$ est
projectif, il est facteur direct d'un $A$-module libre $A^{(I)}$~; en outre,
comme $M$ est de type fini, $u(M)$ est contenu dans un sous-module de $A^{(I)}$
de la forme $A^n$ pour $n$ fini~; comme chacun des énoncés a), b), c), d) est
équivalent à l'énoncé correspondant où l'on remplace $N$ par $N\oplus P$ ($u$
étant toujours supposé appliquer $M$ dans $N$), on est ramené au cas où $N$
est \emph{libre de type fini}. Il est trivial que d) entraîne a), et a) et b)
sont équivalentes en vertu de (19.1.11)~; d'ailleurs a) entraîne que
$M_{\mathfrak p}$ est libre (19.1.11), donc a) entraîne c), en prenant pour les
$x_i$ une famille dont les images dans $M_{\mathfrak p}$ forment une base de ce

\oldpage[0\textsubscript{IV-1}]{77}

$A_{\mathfrak p}$-module et en notant que (puisque $N$ est libre de type fini)
toute forme linéaire sur le $A_{\mathfrak p}$-module $N_{\mathfrak p}$ s'écrit
$u=v/f$, où $f\in A-\mathfrak p$, et où $v$ est une forme linéaire sur le
$A$-module $N$. Il est clair que c) entraîne b), et il reste donc à voir que a)
entraîne d). Or, comme $N$ est de présentation finie,
$(\operatorname{Hom}_A(N,M))_{\mathfrak p}$ s'identifie canoniquement à
$\operatorname{Hom}_{A_{\mathfrak p}}(N_{\mathfrak p},M_{\mathfrak p})$
(Bourbaki, \emph{Alg. comm.}, chap.~II, §~2, n\textsuperscript{o}~7,
prop.~19). Si $w_{\mathfrak p}$ est inverse à gauche de $u_{\mathfrak p}$, il
existe donc un homomorphisme $w:N\to M$ et un élément $f\in A-\mathfrak p$
tels que $w=w_{\mathfrak p}\otimes(1/f)1_{A_{\mathfrak p}}$. La relation
$w_{\mathfrak p}\circ u_{\mathfrak p}=1_{M_{\mathfrak p}}$ s'écrit donc aussi
$(w\circ u)\otimes1_{A_{\mathfrak p}}=f.1_{M_{\mathfrak p}}$. Mais comme $M$
est un $A$-module de type fini, il existe $g\in A-\mathfrak p$ tel que
l'endomorphisme $g((w\circ u)-f.1_M)$ s'annule en tous les générateurs de $M$,
donc soit nul. Si on pose $h=fg$, et
$u_h=u\otimes1_{A_h}$, $w_h=w\otimes1_{A_h}$, on a donc
$w_h(u_h(z))=(f/1)z$ pour tout $z\in M_h$~; mais comme $h/1$ est inversible
dans $A_h$, il en est de même de $f/1=f'$, et $f'^{-1}w_h$ est par suite
inverse à gauche de $u_h$, ce qui achève la démonstration.

\begin{proposition}[19.1.13]
\label{0.19.1.13-fr}
Supposons vérifiées les hypothèses de (19.1.9) et supposons en outre que
$(\mathfrak J_\lambda)$ soit une suite décroissante $(\mathfrak J_n)$ et que
$M$ soit séparé et complet. Alors les conditions a) et b) de (19.1.9) sont
aussi équivalentes à~:
\begin{enumerate}
  \item[c)] $u$ est inversible à gauche.
\end{enumerate}
\end{proposition}

On sait déjà que c) entraîne a) (19.1.6). Inversement, si a) est vérifiée, on
sait (avec les notations de (19.1.8)) que $M_n$ est facteur direct du
$(A/\mathfrak J_n)$-module projectif $N_n$, donc $P_n$ est aussi isomorphe à
un facteur direct de $N_n$, et par suite est projectif~; il suffit donc
d'appliquer (19.1.8).

Notons enfin la proposition suivante~:

\begin{proposition}[19.1.14]
\label{0.19.1.14-fr}
Soient $A$ un anneau, $M$ un $A$-module de type fini, $N$ un $A$-module
projectif, $u:M\to N$ un $A$-homomorphisme.
\begin{enumerate}
  \item[(i)] Pour que $u$ soit inversible à gauche, il faut et il suffit que,
  pour tout idéal maximal $\mathfrak m$ de $A$,
  $u_{\mathfrak m}:M_{\mathfrak m}\to N_{\mathfrak m}$ soit inversible à
  gauche.
  \item[(ii)] Soit $A'$ une $A$-algèbre qui soit un $A$-module fidèlement
  plat. Pour que $u$ soit inversible à gauche, il faut et il suffit que
  $u\otimes1:M\otimes_AA'\to N\otimes_AA'$ soit inversible à gauche.
\end{enumerate}
\end{proposition}

Comme dans (19.1.12), on peut se borner au cas où $N$ est libre de type fini~;
dire que $u$ est inversible à gauche signifie alors que $u$ est injectif et
que le module quotient $P=N/u(M)$ est \emph{projectif}, car $u(M)$ sera alors
facteur direct de $N$. Notons en outre que puisque $M$ est de type fini, $P$
est de \emph{présentation finie}. Cela étant~:

(i) La condition est évidemment nécessaire. Inversement, si elle est
satisfaite, on sait que $u$ est injectif (Bourbaki, \emph{Alg. comm.},
chap.~II, §~3, n\textsuperscript{o}~3, th.~1) et comme
$P_{\mathfrak m}=N_{\mathfrak m}/u_{\mathfrak m}(M_{\mathfrak m})$ est
projectif pour tout $\mathfrak m$, on sait que cela entraîne que $P$ est
projectif (\emph{loc. cit.}, §~5, n\textsuperscript{o}~2, th.~1).

(ii) Ici encore, la condition est trivialement nécessaire. Inversement, si
elle est remplie, on sait que $u$ est injectif ($\mathbf{0}_{\mathrm I}$,
6.4.1) et comme $P\otimes_AA'=\operatorname{Coker}(u\otimes1)$ est projectif,
donc plat, on en déduit que $P$ est un $A$-module plat
($\mathbf{0}_{\mathrm I}$, 6.6.3), donc projectif puisqu'il est de présentation
finie (Bourbaki, \emph{Alg. comm.}, chap.~II, §~5,
n\textsuperscript{o}~2, cor.~2 du th.~1).

\oldpage[0\textsubscript{IV-1}]{78}

\begin{remark}[19.1.15]
\label{0.19.1.15-fr}
La notion «~duale~» de celle d'homomorphisme formellement inversible à gauche
est celle d'homomorphisme \emph{formellement inversible à droite}~; un tel
$A$-homomorphisme continu $u:M\to N$ vérifie, par définition, la condition
suivante~: pour tout sous-module ouvert $W'$ de $N$, il existe un sous-module
ouvert $V'\subset u^{-1}(W')$ de $M$, un sous-module ouvert $W''\subset W'$ de
$N$ et un homomorphisme $h:N/W''\to M/V'$ tels que l'homomorphisme canonique
$N/W''\to N/W'$ se factorise en
\[
  N/W''\xrightarrow{h}M/V'\xrightarrow{u'}N/W',
\]
où $u'$ est déduit de $u$ par passage aux quotients. Cela entraîne que $u$ est
un épimorphisme formel~; s'il existe un $\widehat A$-homomorphisme continu
$r:\widehat N\to\widehat M$ tel que $\widehat u\circ r=1_{\widehat N}$, on
vérifie aussitôt que $u$ est formellement inversible à droite. Si les
conditions de (19.1.7) sont remplies, pour que $u$ soit formellement inversible
à droite, il faut et il suffit que les $u_\lambda$ soient inversibles à droite
(c'est-à-dire que le noyau de $u_\lambda$ est facteur direct de
$M/\mathfrak J_\lambda M$ et que $u_\lambda$ est un isomorphisme d'un
supplémentaire de $\operatorname{Ker}(u_\lambda)$ sur
$N/\mathfrak J_\lambda N$). Nous laissons au lecteur le soin d'énoncer et de
démontrer les propositions correspondant à (19.1.8) et (19.1.9) (dans
l'analogue de (19.1.8), il faut supposer $M$ séparé et complet et que $M_n$ est
un $(A/\mathfrak J_n)$-module projectif~; dans l'analogue de (19.1.9), il faut
supprimer l'hypothèse sur les $N/\mathfrak J_\lambda N$, mais supposer par
contre que, pour tout $\lambda$, $M/\mathfrak J_\lambda M$ est un
$(A/\mathfrak J_\lambda)$-module projectif).
\end{remark}

\subsection{Modules formellement projectifs}
\label{subsection:0.19.2-fr}

\begin{definition}[19.2.1]
\label{0.19.2.1-fr}
Soient $A$ un anneau topologique. On dit qu'un $A$-module topologique $M$ est
\emph{formellement projectif} s'il vérifie la condition suivante~: pour tout
idéal ouvert $\mathfrak J$ de $A$, tout couple de $(A/\mathfrak J)$-modules
(discrets) $P$, $Q$, tout $A$-homomorphisme surjectif $u:P\to Q$ et tout
$A$-homomorphisme continu $v:M\to Q$, il existe un $A$-homomorphisme continu
$w:M\to P$ tel que $v=u\circ w$.
\end{definition}

\begin{env}[19.2.2]
\label{0.19.2.2-fr}
Pour vérifier la condition de (19.2.1), on peut évidemment se borner (en
remplaçant $Q$ par $v(M)$ et $P$ par $u^{-1}(v(M))$) au cas où $v$ est
lui-même surjectif~; alors $Q$ est isomorphe à $M/V$, où $V$ est un
sous-module ouvert de $M$ tel que $\mathfrak JM\subset V$~; la condition de
(19.2.1) équivaut alors à dire que pour tout $(A/\mathfrak J)$-module discret
$P$ et tout $A$-homomorphisme surjectif $u:P\to M/V$, il existe un sous-module
ouvert $V'\subset V$ dans $M$ et un $A$-homomorphisme $w:M/V'\to P$, tels que
l'homomorphisme canonique $M/V'\to M/V$ se factorise en
\[
  M/V'\xrightarrow{w}P\xrightarrow{u}M/V.
\]
On notera qu'il suffit de vérifier cette condition lorsque $\mathfrak J$
parcourt un système fondamental de voisinages de 0 dans $A$ formé d'idéaux.
\end{env}

\begin{env}[19.2.3]
\label{0.19.2.3-fr}
Supposons qu'il existe un système fondamental $(\mathfrak J_\lambda)$ d'idéaux
ouverts dans $A$ et un système fondamental $(V_\lambda)$ de sous-modules
ouverts de $M$, ayant \emph{même} ensemble d'indices que
$(\mathfrak J_\lambda)$ et tel que, pour tout $\lambda$, $M/V_\lambda$ soit un
$(A/\mathfrak J_\lambda)$-module \emph{projectif}. Alors $M$ est formellement
projectif~: il suffit en effet, avec les notations de (19.2.2), de prendre
$\lambda$ tel que $\mathfrak J_\lambda\subset\mathfrak J$ et
$V_\lambda\subset V$~; comme $P$ est aussi un
$(A/\mathfrak J_\lambda)$-module, la factorisation de l'homomorphisme canonique
$M/V_\lambda\to M/V$ en $M/V_\lambda\to P\to M/V$ résulte de ce qu'il s'agit
d'un $(A/\mathfrak J_\lambda)$-homomorphisme et de ce que $M/V_\lambda$ est
supposé être un $(A/\mathfrak J_\lambda)$-module projectif.

Lorsque la condition plus stricte de ce numéro est satisfaite, on dit que $M$
est \emph{strictement formellement projectif}.
\end{env}

\begin{proposition}[19.2.4]
\label{0.19.2.4-fr}
Soient $A$ un anneau topologique, $M$ un $A$-module topologique dont la
topologie est déduite de celle de $A$. Pour que $M$ soit formellement
projectif, il faut et il suffit qu'il soit strictement formellement projectif.
\end{proposition}

\oldpage[0\textsubscript{IV-1}]{79}

\begin{proof}
On vient de voir que la condition est suffisante. Inversement, supposons $M$
formellement projectif et soit $(\mathfrak J_\lambda)$ un système fondamental
d'idéaux ouverts dans $A$. Pour tout $\lambda$, soient $P$ un
$(A/\mathfrak J_\lambda)$-module libre et
$u:P\to M/\mathfrak J_\lambda M$ un
$(A/\mathfrak J_\lambda)$-homomorphisme surjectif. Il existe donc un
$\mathfrak J_\mu\subset\mathfrak J_\lambda$ tel que l'homomorphisme canonique
$M/\mathfrak J_\mu M\to M/\mathfrak J_\lambda M$ se factorise en
\[
  M/\mathfrak J_\mu M\xrightarrow{w}P\xrightarrow{u}M/\mathfrak J_\lambda M;
\]
mais comme
$w(\mathfrak J_\lambda(M/\mathfrak J_\mu M))\subset\mathfrak J_\lambda P=0$,
$w$ se factorise en $M/\mathfrak J_\lambda M\xrightarrow{v}P$, et il est
clair que $u\circ v$ est l'identité dans $M/\mathfrak J_\lambda M$, ce qui
prouve que ce $(A/\mathfrak J_\lambda)$-module est projectif.
\end{proof}

\begin{proposition}[19.2.5]
\label{0.19.2.5-fr}
Soient $A$ un anneau topologique, $M$ un $A$-module topologique.
\begin{enumerate}
  \item[(i)] Pour que $M$ soit formellement projectif (resp. strictement
  formellement projectif), il faut et il suffit que le $\hat A$-module
  topologique $\hat M$ le soit.
  \item[(ii)] Soit $A'$ une $A$-algèbre topologique. Si $M$ est formellement
  projectif (resp. strictement formellement projectif), alors
  $M\otimes_A A'$ (muni de la topologie produit tensoriel) est un $A'$-module
  topologique formellement projectif (resp. strictement formellement
  projectif).
\end{enumerate}
\end{proposition}

\begin{proof}
\begin{enumerate}
  \item[(i)] Il suffit de remarquer que lorsque $\mathfrak J$ (resp. $V$)
  parcourt l'ensemble des idéaux ouverts de $A$ (resp. l'ensemble des
  sous-modules ouverts de $M$), le séparé complété $\hat{\mathfrak J}$
  (resp. $\hat V$) parcourt l'ensemble des idéaux ouverts de $\hat A$
  (resp. l'ensemble des sous-modules ouverts de $\hat M$), et
  $\hat A/\hat{\mathfrak J}=A/\mathfrak J$ (resp.
  $\hat M/\hat V=M/V$) à un isomorphisme canonique près~; comme la notion
  de module formellement projectif (resp. strictement formellement projectif)
  ne fait intervenir que les $A/\mathfrak J$ et les $M/V$, on en déduit
  aussitôt (i).
  \item[(ii)] Supposons d'abord $M$ formellement projectif et posons
  $M'=M\otimes_A A'$~; soient $\mathfrak J'$ un idéal ouvert de $A'$, $P'$,
  $Q'$ deux $(A'/\mathfrak J')$-modules discrets, $u':P'\to Q'$ un
  $A'$-homomorphisme surjectif, $v':M'\to Q'$ un $A'$-homomorphisme continu.
  Il y a un idéal ouvert $\mathfrak J$ de $A$ tel que
  $\mathfrak JA'\subset\mathfrak J'$, donc $P'$ et $Q'$ sont aussi des
  $(A/\mathfrak J)$-modules discrets. Si l'on considère le
  $A$-homomorphisme composé $v:M\xrightarrow{j}M'\xrightarrow{v'}Q'$ qui est
  continu, l'hypothèse entraîne qu'il existe un $A$-homomorphisme continu
  $w:M\to P'$ tel que $v=u'\circ w$~; mais comme $P'$ est un $A'$-module
  topologique, $w$ se factorise en
  $M\xrightarrow{j}M'\xrightarrow{w'}P'$, où $w'$ est un $A'$-homomorphisme
  continu, et comme $v'\circ j=(u'\circ w')\circ j$, on en conclut que
  $v'=u'\circ w'$.

  Supposons ensuite que $M$ soit strictement formellement projectif~; soit
  $(\mathfrak J_\lambda)$ (resp. $(V_\lambda)$) un système fondamental
  d'idéaux ouverts (resp. de sous-modules ouverts) dans $A$ (resp. $M$) tel
  que $(M/V_\lambda)$ soit un $(A/\mathfrak J_\lambda)$-module projectif,
  et soit $(\mathfrak J'_\mu)$ un système fondamental d'idéaux ouverts dans
  $A'$. On a un système fondamental de voisinages de $0$ dans $M'$ en prenant
  les sous-modules
  \[
    \operatorname{Im}(M\otimes_A\mathfrak J'_\mu)
    +\operatorname{Im}(V_\lambda\otimes_A A')=W_{\lambda\mu},
  \]
  où $\lambda$ et $\mu$ sont tels que
  $\mathfrak J_\lambda A'\subset\mathfrak J'_\mu$. Comme
  \[
    (M\otimes_A A')/W_{\lambda\mu}
    =(M/V_\lambda)\otimes_{A/\mathfrak J_\lambda}
      (A'/\mathfrak J'_\mu)
  \]
  et que $(M/V_\lambda)$ est un
  $(A/\mathfrak J_\lambda)$-module projectif, $M'/W_{\lambda\mu}$
  est un $(A'/\mathfrak J'_\mu)$-module projectif.
\end{enumerate}
\end{proof}

\subsection{Algèbres formellement lisses}
\label{subsection:0.19.3-fr}

\begin{definition}[19.3.1]
\label{0.19.3.1-fr}
Soient $A$ un anneau topologique, $B$ une $A$-algèbre topologique. On dit que
$B$ est une $A$-algèbre \emph{formellement lisse} si, pour toute
$A$-algèbre topologique discrète $C$ et tout idéal nilpotent $\mathfrak J$ de
$C$, tout $A$-homomorphisme continu $u:B\to C/\mathfrak J$ se factorise en
\[
  B\xrightarrow{v}C\xrightarrow{\varphi}C/\mathfrak J,
\]
où $v$ est un homomorphisme continu et $\varphi$ l'homomorphisme canonique.
\end{definition}

\oldpage[0\textsubscript{IV-1}]{80}

La définition (19.3.1) revient à dire que la propriété suivante a lieu~:
\begin{enumerate}
  \item[(P)] Pour tout idéal ouvert $\mathfrak R$ de la $A$-algèbre $B$ et
  tout $A$-homomorphisme $u':B/\mathfrak R\to C/\mathfrak J$, il y a un
  idéal ouvert $\mathfrak R'\subset\mathfrak R$ de $B$ tel que
  l'homomorphisme
  $B\to B/\mathfrak R\xrightarrow{u'}C/\mathfrak J$ se factorise en
  \[
    B\to B/\mathfrak R'\xrightarrow{v'}C\to C/\mathfrak J,
  \]
  où $v'$ est un $A$-homomorphisme.
\end{enumerate}
En effet, si $u:B\to C/\mathfrak J$ est un $A$-homomorphisme continu, il a
pour noyau $\mathfrak R$ un idéal ouvert de $B$, donc $u$ se factorise en
$B\to B/\mathfrak R\xrightarrow{u'}C/\mathfrak J$, et si (P) est
vérifiée, il suffit de l'appliquer à $u'$ et de prendre pour $v$ la composée
$B\to B/\mathfrak R'\xrightarrow{v'}C$ pour satisfaire aux conditions de
(19.3.1). Inversement, supposons que $B$ soit une $A$-algèbre formellement
lisse~; donnons-nous un idéal ouvert $\mathfrak R$ de $B$ et un
$A$-homomorphisme $u':B/\mathfrak R\to C/\mathfrak J$, et appliquons la
définition (19.3.1) à
$u:B\to B/\mathfrak R\xrightarrow{u'}C/\mathfrak J$~; si
$v:B\to C$ est un $A$-homomorphisme continu tel que $u$ se factorise en
$B\xrightarrow{v}C\to C/\mathfrak J$, l'idéal
$\mathfrak R'=\operatorname{Ker}(v)\cap\mathfrak R$ de $B$ est ouvert et
contenu dans $\mathfrak R$~; par suite $v$ se factorise en
$B\to B/\mathfrak R'\xrightarrow{v'}C$, et $v'$ vérifie bien la condition
(P).

\begin{proposition}[19.3.2]
\label{0.19.3.2-fr}
Soient $A$ un anneau discret, $V$ un $A$-module projectif~; l'algèbre
symétrique $B=\mathbf S_A^\bullet(V)$, munie de la topologie discrète, est une
$A$-algèbre formellement lisse.
\end{proposition}

\begin{proof}
En effet, les notations $C$ et $\mathfrak J$ ayant le même sens que dans
(19.3.1), soit $u:B\to C/\mathfrak J$ un homomorphisme de $A$-algèbres, qui
par restriction à $V=\mathbf S_A^1(V)$, donne un homomorphisme de $A$-modules
$u_1:V\to C/\mathfrak J$~; comme $V$ est projectif, $u_1$ se factorise en
$V\xrightarrow{v_1}C\xrightarrow{\varphi}C/\mathfrak J$, et $v_1$ se
prolonge en un homomorphisme de $A$-algèbres
$v:\mathbf S_A^\bullet(V)\to C$, tel que le composé
$\varphi\circ v$ coïncide avec $u$ dans $V$, donc
$\varphi\circ v=u$, ce qui démontre la proposition.
\end{proof}

\begin{corollary}[19.3.3]
\label{0.19.3.3-fr}
Si $A$ est un anneau discret, toute algèbre de polynômes
$B=A[T_\alpha]_{\alpha\in I}$, munie de la topologie discrète, est une
$A$-algèbre formellement lisse.
\end{corollary}

\begin{proposition}[19.3.4]
\label{0.19.3.4-fr}
Soit $A$ un anneau topologique, et soit
$B=A[[T_\alpha]]_{\alpha\in I}$ un anneau de séries formelles
$\sum_{\lambda\in\mathbb{N}^{(I)}}c_\lambda T^\lambda$ (algèbre large sur
$A$ du monoïde $\mathbb{N}^{(I)}$, identifié en tant que $A$-module au produit
$A^{\mathbb{N}^{(I)}}$). Si l'on munit $B$ de la topologie produit, $B$ est une
$A$-algèbre formellement lisse.
\end{proposition}

Soit $(\mathfrak L_\mu)_{\mu\in M}$ un système fondamental d'idéaux ouverts
dans $A$. Pour toute partie finie $H$ de $I$, tout $\mu\in M$ et tout entier
$n$, désignons par $\mathfrak K_{H,\mu,n}$ le voisinage de $0$ dans $B$ formé
des $(c_\lambda)$ tels que pour tout
$\lambda=(\lambda_\alpha)_{\alpha\in I}$ tel que
$\lambda_\alpha=0$ pour $\alpha\notin H$ et
$\sum_{\alpha\in H}\lambda_\alpha\leq n$, on ait
$c_\lambda\in\mathfrak L_\mu$. On vérifie immédiatement que les
$\mathfrak K_{H,\mu,n}$ forment un système fondamental de voisinages de $0$
dans $B$ et sont des idéaux de $B$, donc la topologie produit est compatible
avec la structure de $A$-algèbre de $B$.

Notons d'abord, avec les mêmes notations, le
\begin{lemma}[19.3.4.1]
\label{0.19.3.4.1-fr}
Soit $E$ une $A$-algèbre discrète.
\begin{enumerate}
  \item[(i)] Si $f:B\to E$ est un $A$-homomorphisme continu, il existe une
  partie finie $H$ de $I$ telle que $f(T_\alpha)=0$ pour $\alpha\notin H$,
  et $f(T_\alpha)$ est nilpotent dans $E$ pour tout $\alpha\in H$.
  \item[(ii)] Inversement, soient $H$ une partie finie de $I$,
  $(z_\alpha)_{\alpha\in H}$ une famille d'éléments nilpotents de $E$.
  Il existe un $A$-homomorphisme continu $g:B\to E$ et un seul tel que
  $g(T_\alpha)=z_\alpha$ pour $\alpha\in H$ et
  $g(T_\alpha)=0$ pour $\alpha\notin H$.
\end{enumerate}
\end{lemma}

\oldpage[0\textsubscript{IV-1}]{81}

\begin{proof}
(i) résulte aussitôt de ce que $f^{-1}(0)$ est un voisinage de $0$ dans le
$A$-module produit $A^{\mathbb{N}^{(I)}}$, d'où $f(T^\lambda)=0$ sauf pour un
nombre fini de valeurs de $\lambda\in\mathbb{N}^{(I)}$. Pour prouver (ii), il
suffit de remarquer que l'anneau de polynômes
$B'=A[T_\alpha]_{\alpha\in I}$ est dense dans $B$~; l'existence et
l'unicité de la restriction $g|B'$ sont triviales et sa continuité résulte de
l'hypothèse que les $f(T_\alpha)$ pour $\alpha\in H$ sont nilpotents,
car si $(f(T_\alpha))^n=0$ pour tout $\alpha\in H$, on a
$g(T^\lambda)=0$ pour tout $\lambda=(\lambda_\alpha)_{\alpha\in I}$ de
support fini sauf ceux tels que $\lambda_\alpha=0$ pour
$\alpha\notin H$ et $\lambda_\alpha\leq n$ pour
$\alpha\in H$, c'est-à-dire sauf pour un nombre fini de valeurs de
$\lambda$.
\end{proof}

\begin{proof}[Démonstration de la proposition~{\rm(19.3.4)}]
Ce lemme étant établi, et les notations $C$ et $\mathfrak J$ ayant le même
sens que dans (19.3.1), on a $u(T_\alpha)=0$ sauf pour
$\alpha\in H$, $H$ étant une partie finie de $I$, et les
$z_\alpha=u(T_\alpha)$ pour $\alpha\in H$ sont nilpotents dans
$C/\mathfrak J$~; comme $\mathfrak J$ est nilpotent, il existe une famille
$(x_\alpha)_{\alpha\in H}$ d'éléments nilpotents de $C$ dont les images
canoniques dans $C/\mathfrak J$ sont les $z_\alpha$~; si $v$ est le
$A$-homomorphisme continu de $B$ dans $C$ tel que
$v(T_\alpha)=0$ pour $\alpha\notin H$, $v(T_\alpha)=x_\alpha$ pour
$\alpha\in H$, il est clair que $u$ se factorise en
$B\xrightarrow{v}C\to C/\mathfrak J$.
\end{proof}

\begin{proposition}[19.3.5]
\label{0.19.3.5-fr}
Soit $A$ un anneau topologique.
\begin{enumerate}
  \item[(i)] $A$ est une $A$-algèbre formellement lisse.
  \item[(ii)] Si $B$ est une $A$-algèbre formellement lisse et $C$ une
  $B$-algèbre formellement lisse, alors $C$ est une $A$-algèbre formellement
  lisse.
  \item[(iii)] Soient $B$ une $A$-algèbre formellement lisse, $A'$ une
  $A$-algèbre topologique~; alors la $A'$-algèbre topologique
  $B\otimes_A A'$ ($\mathbf 0_{\rm I}$, 7.7.5 et 7.7.6) est formellement
  lisse.
  \item[(iv)] Soient $B$ une $A$-algèbre topologique, $S$ (resp. $T$) une
  partie multiplicative de $A$ (resp. $B$) telle que l'image canonique de $S$
  dans $B$ soit contenue dans $T$. Si $B$ est une $A$-algèbre formellement
  lisse, alors $T^{-1}B$ est une $S^{-1}A$-algèbre formellement lisse.
  \item[(v)] Soient $B_i$ ($1\leq i\leq n$) des $A$-algèbres topologiques.
  Pour que $\prod_{i=1}^n B_i$ soit une $A$-algèbre formellement lisse, il
  faut et il suffit que chacune des $B_i$ le soit.
\end{enumerate}
\end{proposition}

\oldpage[0\textsubscript{IV-1}]{82}

\begin{proof}
\begin{enumerate}
  \item[(i)] Si $C$ est une $A$-algèbre discrète,
  $\varphi:C\to C/\mathfrak J$ l'homomorphisme canonique de $C$ sur une
  $A$-algèbre quotient quelconque de $C$, le \emph{seul} $A$-homomorphisme
  de $A$ dans $C/\mathfrak J$ est l'homomorphisme composé
  $A\xrightarrow{\psi}C\xrightarrow{\varphi}C/\mathfrak J$, où
  $\psi$ est l'homomorphisme définissant la structure de $A$-algèbre de
  $C$~; comme $\psi$ est continu, la condition de (19.3.1) est trivialement
  vérifiée.
  \item[(ii)] Soient $\alpha:A\to B$, $\beta:B\to C$ les homomorphismes
  continus définissant respectivement la structure de $A$-algèbre sur $B$ et
  celle de $B$-algèbre sur $C$, de sorte que $\beta\circ\alpha$ définit la
  structure de $A$-algèbre sur $C$. Soient $E$ une $A$-algèbre discrète,
  $\mathfrak L$ un idéal nilpotent de $E$, $u:C\to E/\mathfrak L$ un
  $A$-homomorphisme continu, de sorte que $u\circ\beta\circ\alpha$ est
  l'homomorphisme définissant la structure de $A$-algèbre de
  $E/\mathfrak L$. Comme $B$ est une $A$-algèbre formellement lisse, le
  $A$-homomorphisme continu $u\circ\beta:B\to E/\mathfrak L$ se
  factorise en $B\xrightarrow{v}E\to E/\mathfrak L$, où $v$ est un
  $A$-homomorphisme continu~; $v$ et $u\circ\beta$ définissent alors sur
  $E$ et $E/\mathfrak L$ respectivement des structures de $B$-algèbre
  topologique, pour lesquelles $E/\mathfrak L$ est encore l'algèbre quotient
  (discrète) de la $B$-algèbre $E$. En outre, $u$ est un $B$-homomorphisme
  continu, donc se factorise en $C\xrightarrow{w}E\to E/\mathfrak L$,
  où $w$ est un $B$-homomorphisme continu~; comme
  $w\circ\alpha$ est le $A$-homomorphisme définissant la structure de
  $A$-algèbre sur $E$, $w$ est bien un $A$-homomorphisme continu, d'où notre
  assertion.
  \item[(iii)] Soient $C$ une $A'$-algèbre topologique discrète,
  $\mathfrak J$ un idéal nilpotent de $C$,
  $u:B\otimes_A A'\to C/\mathfrak J$ un $A'$-homomorphisme continu. Si on
  compose $u$ et l'homomorphisme canonique
  $\rho:B\to B\otimes_A A'$, on obtient (puisque
  $A\to B\xrightarrow{\rho}B\otimes_A A'$ est égal au composé
  $A\to A'\xrightarrow{\sigma}B\otimes_A A'$) un $A$-homomorphisme
  continu qui, par hypothèse, se factorise donc en
  $B\xrightarrow{v}C\to C/\mathfrak J$, où $v$ est un $A$-homomorphisme
  continu. L'égalité des composés $A\to B\xrightarrow{v}C$ et
  $A\to A'\xrightarrow{w}C$ entraîne donc l'existence d'un homomorphisme
  continu d'anneaux $f:B\otimes_A A'\to C$ tel que
  $v=f\circ\rho$ et $w=f\circ\sigma$ ($\mathbf 0_{\rm I}$, 7.7.6)~;
  comme les homomorphismes composés
  $B\xrightarrow{\rho}B\otimes_A A'\xrightarrow{f}C\to C/\mathfrak J$
  et
  $B\xrightarrow{\rho}B\otimes_A A'\xrightarrow{u}C/\mathfrak J$
  (resp.
  $A'\xrightarrow{\sigma}B\otimes_A A'\xrightarrow{f}C\to C/\mathfrak J$
  et
  $A'\xrightarrow{\sigma}B\otimes_A A'\xrightarrow{u}C/\mathfrak J$)
  sont égaux, on a bien la factorisation
  $u:B\otimes_A A'\xrightarrow{f}C\to C/\mathfrak J$, ce qui établit
  (iii).
  \item[(iv)] La topologie considérée sur $S^{-1}A$ (resp. $T^{-1}B$) est
  naturellement celle pour laquelle un système fondamental de voisinages de
  $0$ est formé des $S^{-1}\mathfrak J$ (resp. $T^{-1}\mathfrak R$), où
  $\mathfrak J$ (resp. $\mathfrak R$) parcourt un système fondamental
  d'idéaux ouverts dans $A$ (resp. $B$) ($\mathbf 0_{\rm I}$, 7.6.1). Si
  $\alpha:A\to B$ est l'homomorphisme canonique, il est clair que
  l'homomorphisme canonique $\alpha':S^{-1}A\to T^{-1}B$ déduit de
  $\alpha$ est continu. Soient alors $C$ une $S^{-1}A$-algèbre topologique
  discrète, $\mathfrak J$ un idéal nilpotent de cette algèbre,
  $u:T^{-1}B\to C/\mathfrak J$ un $S^{-1}A$-homomorphisme continu~; alors
  le composé $B\to T^{-1}B\xrightarrow{u}C/\mathfrak J$ est un
  $A$-homomorphisme continu qui, par hypothèse, se factorise en
  $B\xrightarrow{v}C\to C/\mathfrak J$, où $v$ est un $A$-homomorphisme
  continu. Comme pour tout $t\in T$, $t/1$ est inversible dans $T^{-1}B$,
  $u(t/1)$ est inversible dans $C/\mathfrak J$. Comme $\mathfrak J$ est
  nilpotent, tout élément de la classe $u(t/1)$ dans $C$, et en particulier
  $v(t)$, est inversible dans $C$ ($\mathbf 0_{\rm I}$, 7.1.12), et par suite
  $v$ se factorise en $B\to T^{-1}B\xrightarrow{w}C$~; comme $v$ est
  continu, il en est de même de $w$ ($\mathbf 0_{\rm I}$, 7.6.6), et c'est
  un $S^{-1}A$-homomorphisme car le composé
  $A\to B\to T^{-1}B\xrightarrow{w}C$ est égal à
  \[
    A\to S^{-1}A\xrightarrow{\alpha'}T^{-1}B\xrightarrow{w}C,
  \]
  donc $S^{-1}A\xrightarrow{\alpha'}T^{-1}B\xrightarrow{w}C$ est
  l'homomorphisme canonique définissant sur $C$ la structure de
  $S^{-1}A$-algèbre. Enfin, les homomorphismes composés
  $B\to T^{-1}B\xrightarrow{w}C\to C/\mathfrak J$ et
  $B\to T^{-1}B\xrightarrow{u}C/\mathfrak J$ étant égaux, le même
  raisonnement montre que $u$ est bien égal au composé
  $T^{-1}B\xrightarrow{w}C\to C/\mathfrak J$, d'où (iv).
  \item[(v)] Enfin, (v) est immédiat, la donnée d'un $A$-homomorphisme
  continu de $B=\prod_{i=1}^n B_i$ dans $C$ (resp. $C/\mathfrak J$)
  équivalant à celle de $n$ $A$-homomorphismes continus $B_i\to C$
  (resp. $B_i\to C/\mathfrak J$) et tout $A$-homomorphisme continu
  $B_j\to C$ (resp. $B_j\to C/\mathfrak J$) donnant par composition un
  $A$-homomorphisme continu $B\to B_j\to C$
  (resp. $B\to B_j\to C/\mathfrak J$).
\end{enumerate}
\end{proof}

\oldpage[0\textsubscript{IV-1}]{83}

\begin{proposition}[19.3.6]
\label{0.19.3.6-fr}
Soient $A$ un anneau topologique, $B$ une $A$-algèbre topologique, $\hat A$
et $\hat B$ les séparés complétés respectifs de $A$ et $B$, de sorte que
$\hat B$ est une $\hat A$-algèbre topologique. Les conditions suivantes
sont équivalentes~:
\begin{enumerate}
  \item[a)] $B$ est une $A$-algèbre formellement lisse.
  \item[b)] $\hat B$ est une $A$-algèbre formellement lisse.
  \item[c)] $\hat B$ est une $\hat A$-algèbre formellement lisse.
\end{enumerate}
\end{proposition}

\begin{proof}
Bien entendu, la structure de $\hat A$-algèbre sur $\hat B$ est définie
par l'homomorphisme $\hat\varphi$, si $\varphi:A\to B$ est
l'homomorphisme définissant la structure de $A$-algèbre sur $B$. Comme toute
$A$-algèbre discrète $C$ est séparée et complète, elle est aussi une
$\hat A$-algèbre (par prolongement canonique de l'homomorphisme de $A$ dans
$C$), et un $A$-homomorphisme continu de $B$ dans $C$ donne par prolongement
canonique un $\hat A$-homomorphisme (et \emph{a fortiori} un
$A$-homomorphisme) de $\hat B$ dans $C$ (autrement dit, tout
$A$-homomorphisme $B\to C$ se factorise en $B\to\hat B\to C$ de façon
unique). Ces remarques et la définition (19.3.1) entraînent aussitôt la
proposition.
\end{proof}

\begin{corollary}[19.3.7]
\label{0.19.3.7-fr}
Sous les hypothèses de (19.3.5), (iv), la $A\{S^{-1}\}$-algèbre topologique
$B\{T^{-1}\}$ est formellement lisse.
\end{corollary}

Cela résulte des définitions ($\mathbf 0_{\rm I}$, 7.6.1 et 7.6.7) et de
(19.3.5), (iv) et (19.3.6).

\begin{proposition}[19.3.8]
\label{0.19.3.8-fr}
Soient $A$ un anneau topologique, $B$ une $A$-algèbre topologique, et supposons
que pour tout idéal ouvert $\mathfrak R$ de $B$, $\mathfrak R^2$ soit aussi
ouvert. Soient $A'$, $B'$ des anneaux topologiques obtenus en munissant $A$ et
$B$ de topologies plus fines que les topologies initiales, et supposons que
l'homomorphisme canonique $\varphi:A\to B$ soit encore un homomorphisme
continu de $A'$ dans $B'$. Alors, si $B'$ est une $A'$-algèbre formellement
lisse, $B$ est une $A$-algèbre formellement lisse.
\end{proposition}

Il suffit d'appliquer le lemme suivant~:

\begin{lemma}[19.3.8.1]
\label{0.19.3.8.1-fr}
Soient $C$ une $A$-algèbre discrète, $\mathfrak J$ un idéal nilpotent de $C$.
Supposons que le carré de tout idéal ouvert de $B$ soit ouvert. Alors, si un
homomorphisme composé
$v:B\xrightarrow{u}C\to C/\mathfrak J$ est continu, l'homomorphisme $u$
est continu.
\end{lemma}

\begin{proof}
En effet, $\mathfrak R=u^{-1}(\mathfrak J)$ est le noyau de $v$ et il est
ouvert par hypothèse~; comme il existe un entier $n$ tel que
$\mathfrak J^n=0$, $\mathfrak R^n$ est contenu dans
$\operatorname{Ker}(u)$~; mais par récurrence sur $h$, il résulte de
l'hypothèse que $\mathfrak R^{2^h}$ est ouvert, donc aussi
$\mathfrak R^n$ pour tout $n$, et par suite $\operatorname{Ker}(u)$ est lui
aussi ouvert, ce qui prouve notre assertion.
\end{proof}

On notera que l'hypothèse sur $B$ signifie que la topologie de $B$ est borne
supérieure des topologies $\mathfrak R$-préadiques, où $\mathfrak R$ parcourt
l'ensemble des idéaux ouverts de $B$. Si $B$ est un anneau préadique
($\mathbf 0_{\rm I}$, 7.1.9), cette condition est toujours vérifiée.

\begin{env}[19.3.9]
\label{0.19.3.9-fr}
Nous allons maintenant voir que la propriété d'être formellement lisse
implique des propriétés de «~relèvement~» d'homomorphismes sous des conditions
plus générales que celles de la définition (19.3.1).
\end{env}

\oldpage[0\textsubscript{IV-1}]{84}

\begin{proposition}[19.3.10]
\label{0.19.3.10-fr}
Soient $A$ un anneau topologique, $B$ une $A$-algèbre formellement lisse.
Soient $C$ une $A$-algèbre topologique, $\mathfrak J$ un idéal de $C$,
vérifiant les conditions suivantes~:
\begin{enumerate}
  \item[$1^\circ$] $C$ est métrisable et complet.
  \item[$2^\circ$] $\mathfrak J$ est fermé et la suite
  $(\mathfrak J^n)$ tend vers $0$.
\end{enumerate}
Alors, tout $A$-homomorphisme continu $u:B\to C/\mathfrak J$ se factorise en
$B\xrightarrow{v}C\to C/\mathfrak J$, où $v$ est un
$A$-homomorphisme continu.
\end{proposition}

\begin{proof}
Soit $(\mathfrak Q_n)$ une suite décroissante d'idéaux de $C$ formant un
système fondamental de voisinages de $0$. Pour tout $n$, considérons la
$A$-algèbre discrète $C/\mathfrak Q_n$ et l'idéal
$(\mathfrak J+\mathfrak Q_n)/\mathfrak Q_n$ de cette algèbre~; comme il
existe $k$ tel que $\mathfrak J^k\subset\mathfrak Q_n$,
$(\mathfrak J+\mathfrak Q_n)/\mathfrak Q_n$ est nilpotent. Pour tout
$n$, soit $u_n$ l'homomorphisme continu
\[
  B\xrightarrow{u}C/\mathfrak J\to C/(\mathfrak J+\mathfrak Q_n)
  =(C/\mathfrak Q_n)/((\mathfrak J+\mathfrak Q_n)/\mathfrak Q_n)~;
\]
par hypothèse $u_n$ se factorise en
\[
  B\xrightarrow{v_n}C/\mathfrak Q_n
  \xrightarrow{\varphi_n}C/(\mathfrak J+\mathfrak Q_n),
\]
où $v_n$ est un $A$-homomorphisme continu. Montrons qu'on peut choisir de
proche en proche les $v_n$ de sorte que $v_n$ se factorise en
\[
  B\xrightarrow{v_{n+1}}C/\mathfrak Q_{n+1}\to C/\mathfrak Q_n
\]
pour tout $n$ (autrement dit, que les $v_n$ forment un \emph{système
projectif} d'homomorphismes). Cela résultera du lemme suivant.
\end{proof}

\begin{lemma}[19.3.10.1]
\label{0.19.3.10.1-fr}
Soient $B$ une $A$-algèbre formellement lisse, $E$, $E'$ deux $A$-algèbres
topologiques discrètes, $\mathfrak R$ (resp. $\mathfrak R'$) un idéal
nilpotent de $E$ (resp. $E'$), $f:E\to E'$ un $A$-homomorphisme surjectif tel
que $f(\mathfrak R)\subset\mathfrak R'$, $f':E/\mathfrak R\to
E'/\mathfrak R'$ l'homomorphisme déduit de $f$ par passage aux quotients.
Soient $u:B\to E/\mathfrak R$ un $A$-homomorphisme continu, $u'$
l'homomorphisme composé
$B\xrightarrow{u}E/\mathfrak R\xrightarrow{f'}E'/\mathfrak R'$, et soit
$v':B\to E'$ un $A$-homomorphisme continu tel que $u'$ se factorise en
$B\xrightarrow{v'}E'\to E'/\mathfrak R'$. Alors il existe un
$A$-homomorphisme continu $v:B\to E$ tel que $v'$ se factorise en
$B\xrightarrow{v}E\xrightarrow{f}E'$.
\end{lemma}

Dans le diagramme commutatif
\[
\xymatrix{
  E \ar[r]^f \ar[d] & E' \ar[d] \\
  E/\mathfrak R \ar[r]_{f'} & E'/\mathfrak R'
}
\]
tous les homomorphismes sont surjectifs~; si $\mathfrak L$ est le noyau de
$f$, $E'$ s'identifie à $E/\mathfrak L$ et $\mathfrak R'$ à
$(\mathfrak R+\mathfrak L)/\mathfrak L$. Utilisons maintenant le lemme
élémentaire suivant.

\begin{lemma}[19.3.10.2]
\label{0.19.3.10.2-fr}
Soient $F$ un anneau (non nécessairement commutatif), $\mathfrak a$,
$\mathfrak b$ deux idéaux bilatères de $F$~; alors le produit fibré
\[
  (F/\mathfrak a)\times_{F/(\mathfrak a+\mathfrak b)}(F/\mathfrak b)
\]
s'identifie canoniquement à $F/(\mathfrak a\cap\mathfrak b)$.
\end{lemma}

\begin{proof}
Ce n'est autre qu'une application particulière de (18.1.7), où, dans le
diagramme (18.1.7.1), on remplace $G$ par
$F/(\mathfrak a\cap\mathfrak b)$, $\mathfrak J'$ par
$\mathfrak b/(\mathfrak a\cap\mathfrak b)$, $F$ par
$F/\mathfrak b$, $\mathfrak R$ par
$(\mathfrak a+\mathfrak b)/\mathfrak b$, $B$ par
$F/(\mathfrak a+\mathfrak b)$ et $\theta$ par le $F$-homomorphisme
bijectif canonique
$(\mathfrak a+\mathfrak b)/\mathfrak b\to
\mathfrak a/(\mathfrak a\cap\mathfrak b)$.
\end{proof}

\oldpage[0\textsubscript{IV-1}]{85}

\begin{proof}[Suite de la démonstration du lemme~{\rm(19.3.10.1)}]
Appliquant ce lemme à la situation de (19.3.10.1), la commutativité du
diagramme
\[
\xymatrix{
  B \ar[r]^{v'} \ar[d]_u & E'=E/\mathfrak L \ar[d] \\
  E/\mathfrak R \ar[r]_{f'} & E'/\mathfrak R'
}
\]
montre l'existence d'un homomorphisme
$w:B\to E/(\mathfrak R\cap\mathfrak L)$ tel que $v'$ et $u$ se
factorisent en
$B\xrightarrow{w}E/(\mathfrak R\cap\mathfrak L)\to E/\mathfrak L$
et
$B\xrightarrow{w}E/(\mathfrak R\cap\mathfrak L)\to E/\mathfrak R$
respectivement~; comme le noyau de $w$ contient l'intersection de ceux de
$v'$ et de $u$, il est ouvert dans $B$ et $w$ est continu. Enfin, il est clair
que $\mathfrak R\cap\mathfrak L$ est nilpotent, donc $w$ se factorise en
$B\xrightarrow{v}E\to E/(\mathfrak R\cap\mathfrak L)$, où $v$ est
un $A$-homomorphisme continu répondant à la question.
\end{proof}

\begin{proof}[Suite de la démonstration de la proposition~{\rm(19.3.10)}]
L'existence de $v_{n+1}$ résulte du lemme (19.3.10.1) appliqué en remplaçant
$E$, $E'$, $\mathfrak R$, $\mathfrak R'$ par
$C/\mathfrak Q_{n+1}$, $C/\mathfrak Q_n$,
$(\mathfrak J+\mathfrak Q_{n+1})/\mathfrak Q_{n+1}$ et
$(\mathfrak J+\mathfrak Q_n)/\mathfrak Q_n$ respectivement, $u$, $u'$
et $v'$ par $u_{n+1}$, $u_n$ et $v_n$ respectivement. Notons que, comme $C$
est séparé et complet, il est égal à $\varprojlim_n(C/\mathfrak Q_n)$, et
$v=\varprojlim_n v_n$ est un $A$-homomorphisme continu de $B$ dans $C$.
Comme $\mathfrak J$ est fermé dans $C$, on a
$\mathfrak J=\bigcap_n(\mathfrak J+\mathfrak Q_n)$~; comme
$C/\mathfrak J$ est métrisable et complet et que les
$(\mathfrak J+\mathfrak Q_n)/\mathfrak J$ forment un système fondamental
de voisinages de $0$ dans $C/\mathfrak J$, on a
\[
  C/\mathfrak J=\varprojlim_n C/(\mathfrak J+\mathfrak Q_n)
\]
et $\varprojlim_n\varphi_n$ est l'application canonique
$\varphi:C\to C/\mathfrak J$. Comme $\varprojlim_n u_n=u$, on a bien
$u=v\circ\varphi$, C.Q.F.D.
\end{proof}

\begin{corollary}[19.3.11]
\label{0.19.3.11-fr}
Soient $A$ un anneau topologique, $B$ une $A$-algèbre formellement lisse,
$C$ un anneau local noethérien complet, $\mathfrak J$ un idéal distinct de
$C$, $\varphi:A\to C$ un homomorphisme continu, faisant de $C$ une
$A$-algèbre topologique. Alors tout $A$-homomorphisme continu
$u_0:B\to C_0=C/\mathfrak J$ se factorise en
$B\xrightarrow{u}C\to C/\mathfrak J$, où $u$ est un $A$-homomorphisme
continu.
\end{corollary}

Toutes les conditions de (19.3.10) sont en effet remplies
($\mathbf 0_{\rm I}$, 7.3.5).

\begin{proposition}[19.3.12]
\label{0.19.3.12-fr}
Soient $A$ un anneau topologique, $B$ une $A$-algèbre formellement lisse,
$C$ une $A$-algèbre topologique, $\mathfrak J$ un idéal de $C$, vérifiant les
conditions suivantes~:
\begin{enumerate}
  \item[$1^\circ$] Il existe un système fondamental d'idéaux ouverts
  $\mathfrak Q_\lambda$ de $C$, tels que les
  $C_\lambda=C/\mathfrak Q_\lambda$ soient des anneaux artiniens et que
  l'homomorphisme canonique $C\to\varprojlim_\lambda C_\lambda$ soit un
  isomorphisme d'anneaux topologiques.
  \item[$2^\circ$] L'idéal $\mathfrak J$ est fermé dans $C$ et
  topologiquement nilpotent.
  \item[$3^\circ$] Le carré de tout idéal ouvert de $B$ est ouvert.
\end{enumerate}
Dans ces conditions, tout $A$-homomorphisme continu
$u:B\to C/\mathfrak J$ se factorise en
$B\xrightarrow{v}C\to C/\mathfrak J$, où $v$ est un $A$-homomorphisme
continu.
\end{proposition}

\begin{proof}
Soit $\mathfrak J_\lambda=(\mathfrak J+\mathfrak Q_\lambda)/\mathfrak Q_\lambda$
l'image canonique de $\mathfrak J$ dans $C_\lambda$, qui est un idéal
nilpotent, et soit $u_\lambda$ l'homomorphisme composé
$B\xrightarrow{u}C/\mathfrak J\to C_\lambda/\mathfrak J_\lambda
=C/(\mathfrak J+\mathfrak Q_\lambda)$. Montrons que tout
$A$-homomorphisme $w:B\to C$ tel que $u$ se factorise en
$B\xrightarrow{w}C\to C/\mathfrak J$ est nécessairement continu~; en
effet, pour tout $\lambda$, il y a un idéal ouvert $\mathfrak R$ de $B$
tel que $u(\mathfrak R)\subset
(\mathfrak J+\mathfrak Q_\lambda)/\mathfrak J$, d'où
$w(\mathfrak R)\subset\mathfrak J+\mathfrak Q_\lambda$, donc il y a
une puissance $\mathfrak R^{2^m}$ de $\mathfrak R$ telle que
$w(\mathfrak R^{2^m})\subset\mathfrak Q_\lambda$, ce qui établit
notre assertion puisque $\mathfrak R^{2^m}$ est ouvert dans $B$. On peut
donc se borner à trouver un $A$-homomorphisme $w$ ayant la propriété précédente
sans s'inquiéter de ses propriétés de continuité. Or, l'ensemble
$\operatorname{Hom}(B,C)$ de tous les $A$-homomorphismes d'algèbres de $B$
dans $C$ est égal à
$\varprojlim_\lambda\operatorname{Hom}(B,C_\lambda)$, et ce dernier est
\emph{fermé} dans le $C_\lambda$-module $C_\lambda^B$, muni de la topologie
produit, pour laquelle il est \emph{linéairement compact} puisque
  $C_\lambda$ est artinien. Pour tout $\lambda$, soit $W_\lambda$

\oldpage[0\textsubscript{IV-1}]{86}

l'ensemble des $w_\lambda\in\operatorname{Hom}(B,C_\lambda)$ tels que
$u_\lambda$ se factorise en
$B\xrightarrow{w_\lambda}C_\lambda\to
C_\lambda/\mathfrak J_\lambda$~; $W_\lambda$ est une variété
linéaire fermée dans le $C_\lambda$-module
  $\operatorname{Hom}(B,C_\lambda)$, non vide puisque $B$ est formellement
  lisse. Dans le produit
$\prod_{\mu\leq\lambda}\operatorname{Hom}(B,C_\mu)=E_\lambda$,
considérons la sous-variété linéaire $U_\lambda$ formée des
$(w_\mu)_{\mu\leq\lambda}$ tels que
$w_\mu=\varphi_{\mu\lambda}\circ w_\lambda$ pour
$\mu\leq\lambda$ et $w_\lambda\in W_\lambda$, où
$\varphi_{\mu\lambda}:C_\lambda\to C_\mu$ est
l'homomorphisme canonique, qui est aussi fermée. Enfin, soit $V_\lambda$
la variété linéaire dans le produit
$\prod_\mu\operatorname{Hom}(B,C_\mu)$, produit de $U_\lambda$
et des $\operatorname{Hom}(B,C_\mu)$ pour
$\mu\not\leq\lambda$, qui est encore fermée. Tout revient à voir que
l'intersection des $V_\lambda$ est \emph{non vide}, car un élément $w$ de
cette intersection appartiendra à
$\varprojlim_\lambda\operatorname{Hom}(B,C_\lambda)$ par définition.
D'ailleurs, comme $\mathfrak J$ est fermé, et
$C=\varprojlim_\lambda C_\lambda$ linéairement compact,
$C/\mathfrak J$ s'identifie à
$\varprojlim_\lambda C_\lambda/\mathfrak J_\lambda$ et
l'application canonique $\psi:C\to C/\mathfrak J$ à
$\varprojlim_\lambda\psi_\lambda$, où $\psi_\lambda$ est
l'application canonique
$C_\lambda\to C_\lambda/\mathfrak J_\lambda$. On conclut donc
comme dans (19.3.10) que $\psi\circ w=u$.
\end{proof}

\subsection{Premiers critères de lissité formelle}
\label{subsection:0.19.4-fr}

\begin{proposition}[19.4.1]
\label{0.19.4.1-fr}
Soient $A$ un anneau topologique, $B$ une $A$-algèbre topologique~; supposons
qu'il existe deux familles filtrantes décroissantes
$(\mathfrak J_\alpha)_{\alpha\in I}$,
$(\mathfrak R_\alpha)_{\alpha\in I}$ d'idéaux de $A$ et $B$
respectivement, telles que~:
\begin{enumerate}
  \item[$1^\circ$] $(\mathfrak J_\alpha)$ tend vers $0$ dans $A$ et
  $(\mathfrak R_\alpha)$ tend vers $0$ dans $B$~;
  \item[$2^\circ$] pour tout $\alpha\in I$ on a
  $\mathfrak J_\alpha B\subset\mathfrak R_\alpha$ (de sorte que
  $B/\mathfrak R_\alpha$ est une
  $(A/\mathfrak J_\alpha)$-algèbre topologique)~;
  \item[$3^\circ$] pour tout $\alpha\in I$,
  $B/\mathfrak R_\alpha$ est une
  $(A/\mathfrak J_\alpha)$-algèbre formellement lisse.
\end{enumerate}
Alors $B$ est une $A$-algèbre formellement lisse.
\end{proposition}

\begin{proof}
En effet, soient $C$ une $A$-algèbre discrète, $\mathfrak L$ un idéal
nilpotent de $C$, $\mathfrak R$ un idéal ouvert de $B$~; par hypothèse il y a
un $\alpha\in I$ tel que
$\mathfrak R_\alpha\subset\mathfrak R$, donc $B/\mathfrak R$ est
quotient de $B/\mathfrak R_\alpha$ par un idéal ouvert. Tout
$A$-homomorphisme $u':B/\mathfrak R\to C/\mathfrak L$ est aussi un
$(A/\mathfrak J_\alpha)$-homomorphisme, donc il existe un idéal ouvert
$\mathfrak R'$ de $B$ tel que
$\mathfrak R_\alpha\subset\mathfrak R'\subset\mathfrak R$ et que
$B/\mathfrak R_\alpha\to B/\mathfrak R\xrightarrow{u'}C/\mathfrak L$
se factorise en
$B/\mathfrak R_\alpha\to B/\mathfrak R'\to C\to C/\mathfrak L$,
où le second homomorphisme est un
$(A/\mathfrak J_\alpha)$-homomorphisme~; d'où la conclusion, en vertu de
la remarque suivant (19.3.1).
\end{proof}

\begin{corollary}[19.4.2]
\label{0.19.4.2-fr}
Soient $A$ un anneau topologique, $B$ une $A$-algèbre topologique,
$(\mathfrak J_\alpha)_{\alpha\in I}$ une famille filtrante décroissante
d'idéaux de $A$ tendant vers $0$. Pour que $B$ soit une $A$-algèbre
formellement lisse, il faut et il suffit que pour tout $\alpha\in I$, si
l'on pose $A_\alpha=A/\mathfrak J_\alpha$,
$B_\alpha=B/\mathfrak J_\alpha B$, $B_\alpha$ soit une
$A_\alpha$-algèbre formellement lisse.
\end{corollary}

La condition est suffisante par (19.4.1), et elle est aussi nécessaire par
(19.3.5), (iii).

\begin{proposition}[19.4.3]
\label{0.19.4.3-fr}
Soient $A$ un anneau topologique, $B$ une $A$-algèbre topologique. Supposons
que pour toute $A$-algèbre discrète $C$ et tout idéal $\mathfrak J$ de $C$
tel que $\mathfrak J^2=0$, tout $A$-homomorphisme continu
$u:B\to C/\mathfrak J$ se factorise en
$B\xrightarrow{v}C\to C/\mathfrak J$, où $v$ est un $A$-homomorphisme
continu. Alors $B$ est une $A$-algèbre formellement lisse.
\end{proposition}

\begin{proof}
En effet, avec les mêmes notations, soit $\mathfrak R$ un idéal nilpotent
quelconque de $C$, et considérons un $A$-homomorphisme continu
$u':B\to C/\mathfrak R$. Supposons que $\mathfrak R^m=0$, et posons
$C_i=C/\mathfrak R^i$ pour $1\leq i\leq m$, de sorte que $C_m=C$, que
$\mathfrak R^{i-1}/\mathfrak R^i$ est un idéal de carré nul
$\mathfrak J_i$ dans $C_i$, et $C_i=C_{i+1}/\mathfrak J_{i+1}$~;
l'hypothèse entraîne alors par récurrence sur $i$ l'existence de
$A$-homomorphismes continus $u_i:B\to C_i$ tels que $u_1=u$ et que $u_i$
se factorise en
$B\xrightarrow{u_{i+1}}C_{i+1}\to C_{i+1}/\mathfrak J_{i+1}=C_i$~;
d'où la conclusion.
\end{proof}

\begin{proposition}[19.4.4]
\label{0.19.4.4-fr}
Soient $A$ un anneau topologique, $B$ une $A$-algèbre topologique
(commutative). Pour que $B$ soit une $A$-algèbre formellement lisse, il faut et
il suffit que pour

\oldpage[0\textsubscript{IV-1}]{87}

tout $B$-module topologique discret $L$, annulé par un idéal ouvert de $B$, on
ait (cf. (18.5.1))
\[
  \operatorname{Exalcotop}_A(B,L)=0.
\]
\end{proposition}

\begin{proof}
Soit $(\mathfrak R_\lambda)$ un système fondamental décroissant d'idéaux
ouverts de $B$ et posons $B_\lambda=B/\mathfrak R_\lambda$.
Considérons une $A$-algèbre topologique discrète $C$ et un idéal
$\mathfrak J$ de $C$ tel que $\mathfrak J^2=0$, de sorte que $C$ est une
$A$-extension de $C/\mathfrak J$ par $\mathfrak J$~; supposons donné un
$A$-homomorphisme $u:B_\lambda\to C/\mathfrak J$ et formons la
$A$-extension $E_\lambda$ de $B_\lambda$ par $\mathfrak J$, image
réciproque de $C$ par l'homomorphisme $u$ (18.2.5)~; c'est une $A$-algèbre
topologique pour la topologie discrète. Si
$\operatorname{Exalcotop}_A(B,L)=0$, il existe par définition (18.5.1) un
$\mu$ tel que $\mathfrak R_\mu\subset\mathfrak R_\lambda$ et que
l'extension image réciproque $E_\mu$ soit $A$-triviale~; mais cela signifie
(18.1.6) qu'il existe un $A$-homomorphisme continu
$v:B_\mu\to E_\lambda$ tel que l'homomorphisme canonique
$B_\mu\to B_\lambda$ se factorise en
$B_\mu\xrightarrow{v}E_\lambda\to B_\lambda$~; on en conclut
aussitôt que $B_\mu\to B_\lambda\xrightarrow{u}C/\mathfrak J$ se
factorise en
$B_\mu\xrightarrow{v}E_\lambda\to C\to C/\mathfrak J$, et cela
prouve, en vertu de (19.3.1) et (19.4.3), que $B$ est une $A$-algèbre
formellement lisse. La réciproque est immédiate, en appliquant (19.3.1) au cas
où $C$ est une $A$-algèbre topologique qui est une $A$-extension de
$B_\lambda$ par $L$, et à l'homomorphisme identique
$B_\lambda\to B_\lambda=C/L$.
\end{proof}

Lorsque $A$ et $B$ sont des anneaux discrets, le critère de lissité formelle
(19.4.4) se réduit à
\begin{equation}
\tag{19.4.4.1}
\label{0.19.4.4.1-fr}
\operatorname{Exalcom}_A(B,L)=0\qquad\text{pour tout $B$-module $L$};
\end{equation}
autrement dit, \emph{toute $A$-extension commutative de $B$ par un
$B$-module doit être $A$-triviale}.

\begin{corollary}[19.4.5]
\label{0.19.4.5-fr}
Soient $A$ un anneau topologique, $B$ une $A$-algèbre topologique
(commutative).
\begin{enumerate}
  \item[(i)] Supposons que $B$ soit une $A$-algèbre formellement lisse.
  Alors, pour tout idéal ouvert $\mathfrak R$ de $B$, tout
  $(B/\mathfrak R)$-module $L$ et tout $2$-cocycle de Hochschild
  $A$-bilinéaire et symétrique $f$ de
  $(B/\mathfrak R)\times(B/\mathfrak R)$ dans $L$ (18.4.3), il existe un
  idéal ouvert $\mathfrak R'\subset\mathfrak R$ tel que, si
  $\varphi:B/\mathfrak R'\to B/\mathfrak R$ est l'homomorphisme
  canonique, $f\circ(\varphi\times\varphi)$ soit un $2$-cobord de
  Hochschild $A$-bilinéaire de
  $(B/\mathfrak R')\times(B/\mathfrak R')$ dans $L$.
  \item[(ii)] Si $B$ est un $A$-module formellement projectif (19.2.1) et si
  la condition de (i) est vérifiée, $B$ est une $A$-algèbre formellement
  lisse.
\end{enumerate}
\end{corollary}

\begin{proof}
\begin{enumerate}
  \item[(i)] Le $2$-cocycle $f$ définit une $A$-extension de Hochschild $C$
  de $B/\mathfrak R$ par $L$ (18.4.3). Appliquant (19.3.1) à $C$, à l'idéal
  de carré nul $L$ de $C$ et à l'homomorphisme identique
  $B/\mathfrak R\to B/\mathfrak R=C/L$, on en déduit la condition (i) en
  vertu de (18.4.3).
  \item[(ii)] Appliquons le critère (19.4.3), en considérant un idéal ouvert
  $\mathfrak R$ de $B$, un idéal ouvert $\mathfrak J$ de $A$ tel que
  $\mathfrak JB\subset\mathfrak R$, et une $(A/\mathfrak J)$-extension $E$
  de $B/\mathfrak R$ par $L$. Comme $B$ est un $A$-module formellement
  projectif, l'application $A$-linéaire continue canonique
  $\rho:B\to B/\mathfrak R$ se factorise en
  $B\xrightarrow{w}E\to B/\mathfrak R$, où $w$ est une application
  $A$-linéaire continue (19.2.1), qui se factorise donc elle-même en
  $B\to B/\mathfrak R'\xrightarrow{w'}E$ où $\mathfrak R'$ est un idéal
  ouvert de $B$ contenu dans $\mathfrak R$~; remplaçant $\mathfrak J$ par
  un idéal plus petit, on peut supposer que
  $\mathfrak JB\subset\mathfrak R'$. Alors l'image réciproque par
  l'homomorphisme canonique $B/\mathfrak R'\to B/\mathfrak R$ de
  l'extension $E$ de $B/\mathfrak R$ par $L$, est une
  $(A/\mathfrak J)$-extension \emph{de Hochschild} $E'$ de
  $B/\mathfrak R'$ par $L$~;

  \oldpage[0\textsubscript{IV-1}]{88}

  appliquant (i) à un cocycle définissant cette extension (18.4.3), on en
  conclut qu'il y a un idéal ouvert
  $\mathfrak R''\subset\mathfrak R'$ de $B$ tel que l'image réciproque
  $E''$ de $E$ par $B/\mathfrak R''\to B/\mathfrak R$ soit
  $A$-triviale, C.Q.F.D.
\end{enumerate}
\end{proof}

\begin{corollary}[19.4.6]
\label{0.19.4.6-fr}
Soient $A$ un anneau topologique, $B$ une $A$-algèbre topologique qui soit un
$A$-module formellement projectif. Soit $A'$ une $A$-algèbre \emph{munie de
la topologie déduite de celle de $A$}. Supposons en outre que $A'$ soit un
$A$-module fidèlement plat, et que l'une des conditions suivantes soit
remplie~:
\begin{enumerate}
  \item[$1^\circ$] Il existe un système fondamental
  $(\mathfrak J_\lambda)$ d'idéaux ouverts de $A$ et un système fondamental
  $(\mathfrak M_\lambda)$ d'idéaux ouverts de $B$, ayant même ensemble
  d'indices et tels que, pour tout $\lambda$, on ait
  $\mathfrak J_\lambda B\subset\mathfrak M_\lambda$ et que
  $B/\mathfrak M_\lambda$ soit un
  $(A/\mathfrak J_\lambda)$-module projectif de type fini.
  \item[$2^\circ$] $A'$ est un $A$-module projectif de type fini.
\end{enumerate}
Alors, pour que $B'=B\otimes_A A'$ (muni de la topologie produit tensoriel)
soit une $A'$-algèbre formellement lisse, il faut et il suffit que $B$ soit une
$A$-algèbre formellement lisse.
\end{corollary}

\begin{proof}
La suffisance de la condition est contenue dans (19.3.5), (iii), sans aucune
hypothèse supplémentaire sur $B$ ni $A'$. Pour démontrer la réciproque, nous
allons appliquer le critère (19.4.5)~; sous l'hypothèse $2^\circ$, nous
désignerons encore par $(\mathfrak M_\lambda)$ un système fondamental
d'idéaux ouverts de $B$, et, pour tout $\lambda$, par
$\mathfrak J_\lambda$ un idéal ouvert de $A$ tel que
$\mathfrak J_\lambda B\subset\mathfrak M_\lambda$~; dans les deux
cas, nous poserons
$A_\lambda=A/\mathfrak J_\lambda$,
$B_\lambda=B/\mathfrak M_\lambda$,
$A'_\lambda=A_\lambda\otimes_A A'$,
$B'_\lambda=B_\lambda\otimes_A A'$. Soit $f$ un $2$-cocycle de
Hochschild $A_\lambda$-bilinéaire symétrique de
$B_\lambda\times B_\lambda$ dans un $B_\lambda$-module $L$~; par
extension des scalaires, on en déduit un $2$-cocycle de Hochschild
$f'=f\otimes1$, $A'_\lambda$-bilinéaire symétrique, de
$B'_\lambda\times B'_\lambda$ dans $L'=L\otimes_A A'$. Comme par
hypothèse $B'$ est une $A'$-algèbre formellement lisse, il existe, par (19.4.5),
(i), un indice $\mu$ tel que
$\mathfrak M_\mu\subset\mathfrak M_\lambda$, et tel que, si
$\varphi':B'_\mu\to B'_\lambda$ est l'application canonique,
$f'\circ(\varphi'\times\varphi')$ soit un $2$-cobord de Hochschild de
$B'_\mu\times B'_\mu$ dans $L'$~; autrement dit, son image $c'$ dans le
groupe de Hochschild
$\mathrm H^2_{A'_\mu}(B'_\mu,L')^s$ est nulle~; il est clair que si
$\varphi:B_\mu\to B_\lambda$ est l'homomorphisme canonique, et $c$ la
classe de $f\circ(\varphi\times\varphi)$ dans le groupe de Hochschild
$\mathrm H^2_{A_\mu}(B_\mu,L)^s$, $c'$ est l'image canonique de $c$.
Or, si $P_\bullet$ est le complexe relatif aux anneaux $A_\mu$ et
$B_\mu$ défini dans (18.4.5), servant au calcul de
$\mathrm H^2_{A_\mu}(B_\mu,L)^s$, le complexe analogue relatif aux
anneaux $A'_\mu$ et $B'_\mu$ est évidemment
$P_\bullet\otimes_A A'$~; dans l'hypothèse $1^\circ$, la construction de
$P_\bullet$ montre que c'est un $A_\mu$-module projectif de type fini.
On conclut donc de Bourbaki, \emph{Alg.}, chap.~II, 3\textsuperscript{e} éd., §~5,
n\textsuperscript{o}~3, prop.~7 que, sous les deux hypothèses, on a
\[
  \operatorname{Hom}_{A'_\mu}(P_\bullet\otimes_A A',L\otimes_A A')
  =\operatorname{Hom}_{A_\mu}(P_\bullet,L)\otimes_A A'
\]
à un isomorphisme canonique près~; comme $A'$ est un $A$-module plat, on a par
suite (18.4.5)
\[
  \mathrm H^2_{A'_\mu}(B'_\mu,L')^s
  =\mathrm H^2_{A_\mu}(B_\mu,L)^s\otimes_A A'
\]
et l'on peut donc écrire $c'=c\otimes1$~; mais comme $A'$ est un $A$-module
fidèlement plat, l'hypothèse $c'=0$ entraîne $c=0$, ce qui achève la
démonstration.
\end{proof}

\begin{proposition}[19.4.7]
\label{0.19.4.7-fr}
Soient $A$ un anneau topologique préadmissible, $\mathfrak J$ un idéal de
définition de $A$ ($\mathbf 0_{\rm I}$, 7.1.2), $B$ une $A$-algèbre
topologique qui soit un $A$-module formellement projectif. Considérons les
anneaux topologiques quotients $A_0=A/\mathfrak J$,
$B_0=B/\mathfrak JB$~; alors, pour que $B$ soit une $A$-algèbre formellement
lisse, il faut et il suffit que $B_0$ soit une $A_0$-algèbre formellement
lisse.
\end{proposition}

\oldpage[0\textsubscript{IV-1}]{89}

\begin{proof}
La nécessité de la condition résulte de (19.4.2). Pour voir qu'elle est
suffisante, notons d'abord qu'en considérant un système fondamental de
voisinages ouverts de $0$ dans $A$ formés d'idéaux
$\mathfrak J_\alpha$ contenus dans $\mathfrak J$, on peut, en vertu de
(19.4.2), se ramener au cas où $A$ est discret puisque
$B/\mathfrak J_\alpha B$ est un
$(A/\mathfrak J_\alpha)$-module formellement projectif (19.2.5)~; par
définition d'un anneau préadmissible, $\mathfrak J$ est alors
\emph{nilpotent}. Il suffit en outre de démontrer la proposition lorsque
$\mathfrak J^2=0$, car si $\mathfrak J^m=0$, on l'appliquera successivement
aux anneaux $A_k=A/\mathfrak J^k$ et $B_k=B/\mathfrak J^kB$ et à l'idéal
$\mathfrak J^{k-1}/\mathfrak J^k$ de $A_k$ pour $k=2,3,\ldots,m$, en
notant (19.2.5) que $B_k=B\otimes_A A_k$ est un $A_k$-module formellement
projectif. Appliquons le critère (19.4.5), (ii) en considérant un idéal ouvert
$\mathfrak R$ de $B$ et un $2$-cocycle de Hochschild $A$-bilinéaire
symétrique $f$ de $(B/\mathfrak R)\times(B/\mathfrak R)$ dans un
$(B/\mathfrak R)$-module $L$. Considérons d'abord le cas particulier où
$\mathfrak JL=0$, de sorte que $L$ peut aussi être considéré comme un
$(B_0/\mathfrak RB_0)$-module, et $f$ se factorise en
\[
  (B/\mathfrak R)\times(B/\mathfrak R)\to
  (B_0/\mathfrak RB_0)\times(B_0/\mathfrak RB_0)
  \xrightarrow{f_0}L,
\]
où $f_0$ est un $2$-cocycle de Hochschild bilinéaire symétrique. En vertu de
l'hypothèse, il y a donc un idéal ouvert
$\mathfrak R'\subset\mathfrak R$ dans $B$ et une application
$A_0$-linéaire $g_0:B_0/\mathfrak R'B_0\to L$ telle que
\[
  f_0(\varphi_0(x),\varphi_0(y))
  =xg_0(y)-g_0(xy)+g_0(x)y
\]
pour $x,y$ dans $B_0/\mathfrak R'B_0$, $\varphi_0$ étant
l'homomorphisme canonique
$B_0/\mathfrak R'B_0\to B_0/\mathfrak RB_0$. On en conclut aussitôt que
l'application $A$-linéaire composée
$g:B/\mathfrak R'\to B_0/\mathfrak R'B_0\xrightarrow{g_0}L$ est telle
que $dg=f\circ(\varphi\times\varphi)$, où
$\varphi:B/\mathfrak R'\to B/\mathfrak R$ est l'homomorphisme
canonique.

Passons maintenant au cas général, et considérons d'abord le
$(B/\mathfrak R)$-module $L/\mathfrak JL=L'$, pour lequel on a
$\mathfrak JL'=0$~; si $f'$ est l'application $A$-bilinéaire composée
$(B/\mathfrak R)\times(B/\mathfrak R)\xrightarrow{f}L\to L'$, $f'$ est
encore un $2$-cocycle de Hochschild symétrique, et en vertu de ce qui précède,
il existe un idéal ouvert $\mathfrak R'\subset\mathfrak R$ dans $B$ et
une application $A$-linéaire $g':B/\mathfrak R'\to L'$ vérifiant
$dg'=f'\circ(\varphi'\times\varphi')$ pour l'application canonique
$\varphi':B/\mathfrak R'\to B/\mathfrak R$. Comme $B$ est un
$A$-module formellement projectif, il existe un idéal ouvert
$\mathfrak R'_1\subset\mathfrak R'$ et une application $A$-linéaire
$g_1:B/\mathfrak R'_1\to L$ telle que l'homomorphisme
$B/\mathfrak R'_1\to B/\mathfrak R'\xrightarrow{g'}L'$ se factorise en
$B/\mathfrak R'_1\xrightarrow{g_1}L\to L'$. Considérons alors le
$2$-cocycle de Hochschild
\[
  f_1(x,y)=f(\varphi_1(x),\varphi_1(y))-xg_1(y)+g_1(xy)-g_1(x)y,
\]
application $A$-bilinéaire symétrique de
$(B/\mathfrak R'_1)\times(B/\mathfrak R'_1)$ dans $L$. Le choix de
$g_1$ entraîne que $f_1$ prend ses valeurs dans $\mathfrak JL$. Comme
$\mathfrak J(\mathfrak JL)=0$, on peut de nouveau appliquer le premier cas,
et il y a donc un idéal ouvert
$\mathfrak R''\subset\mathfrak R'_1$ et une application $A$-linéaire
$g_2:B/\mathfrak R''\to\mathfrak JL$ telle que
\[
  f_1(\varphi_2(x),\varphi_2(y))=xg_2(y)-g_2(xy)+g_2(x)y
\]
dans $(B/\mathfrak R'')\times(B/\mathfrak R'')$,
$\varphi_2:B/\mathfrak R''\to B/\mathfrak R'_1$ étant l'application
canonique~; l'application $A$-linéaire
$g=g_2+g_1\circ\varphi_2:B/\mathfrak R''\to L$ vérifie donc
$dg=f\circ(\varphi\times\varphi)$ pour l'application canonique
$\varphi:B/\mathfrak R''\to B/\mathfrak R$. C.Q.F.D.
\end{proof}

\subsection{Lissité formelle et anneaux gradués associés}
\label{subsection:0.19.5-fr}

\begin{env}[19.5.1]
\label{0.19.5.1-fr}
\oldpage[0\textsubscript{IV-1}]{90}
Soient $C$ un anneau topologique (commutatif), $V$ un $C$-module topologique,
et considérons l'algèbre symétrique
$\mathbf S_C^*(V)=\bigoplus_{n\geq0}\mathbf S_C^n(V)$, que nous allons munir
canoniquement d'une \emph{topologie linéaire}, compatible avec sa structure de
$C$-algèbre. Pour cela, considérons un sous-module ouvert $U$ de $V$, et
l'\emph{idéal gradué} $U.\mathbf S_C^*(V)$ qu'il engendre dans
$\mathbf S_C^*(V)$~; nous prendrons comme système fondamental de voisinages de
$0$ dans $\mathbf S_C^*(V)$ l'ensemble des sommes
$\mathfrak a.\mathbf S_C^*(V)+U.\mathbf S_C^*(V)$, où $\mathfrak a$ (resp.
$U$) parcourt un système fondamental d'idéaux ouverts (resp. de sous-modules
ouverts) de $C$ (resp. $V$). On notera que si la topologie de $V$ est moins
fine que la topologie déduite de celle de $C$, on peut se borner aux couples
$(\mathfrak a,U)$ tels que $\mathfrak aV\subset U$, donc
$\mathfrak a.\mathbf S_C^*(V)+U.\mathbf S_C^*(V)
=\mathfrak a.1_C+U.\mathbf S_C^*(V)$~; dans ce cas la topologie induite sur
chacun des $\mathbf S_C^n(V)$ pour $n\geq1$ a pour système fondamental de
voisinages de $0$ les $U.\mathbf S_C^{n-1}(V)$, où $U$ parcourt les
sous-modules ouverts de $V$~; en particulier, sur $V=\mathbf S_C^1(V)$, elle
coïncide avec la topologie donnée (en général elle est moins fine que cette
dernière). En outre, dans tous les cas, l'\emph{algèbre topologique}
$\mathbf S_C^*(V)$ ainsi définie est, pour les catégories des $C$-modules
topologiques et des $C$-algèbres topologiques, la solution du même problème
universel que pour les catégories de $C$-modules et de $C$-algèbres. En effet,
soit $E$ une $C$-algèbre topologique, $u:V\to E$ un homomorphisme de
$C$-modules, $\mathbf S(u)=\widetilde u$ son extension canonique à
$\mathbf S_C^*(V)$. Supposons $u$ continu~; alors, si $\mathfrak b$ est un
idéal ouvert de $E$, son image réciproque $U=u^{-1}(\mathfrak b)$ est un
sous-$C$-module ouvert de $V$, et l'image par $\widetilde u$ de
$U.\mathbf S_C^*(V)$ est donc contenue dans $\mathfrak b$~; comme par ailleurs
il existe un idéal ouvert $\mathfrak a$ de $C$ tel que
$\mathfrak a.E\subset\mathfrak b$, donc
$\widetilde u(\mathfrak a.\mathbf S_C^*(V))\subset\mathfrak b$, cela prouve
que $\widetilde u$ est continu. Réciproquement, si $\widetilde u$ est continu,
et si $\mathfrak b$ est un idéal ouvert de $E$, il existe un sous-module ouvert
$U$ de $V$ tel que $\widetilde u(U.\mathbf S_C^*(V))\subset\mathfrak b$, et en
particulier $\widetilde u(U.1_C)\subset\mathfrak b$, c'est-à-dire
$u(U)\subset\mathfrak b$, donc $u$ est continu. Rappelons en outre que l'on a
un isomorphisme canonique de $C$-modules topologiques (discrets)
\begin{equation}
\label{0.19.5.1.1-fr}
  \mathbf S_C^n(V)/U.\mathbf S_C^{n-1}(V)
  \xrightarrow{\sim}\mathbf S_C^n(V/U).
  \tag{19.5.1.1}
\end{equation}
\end{env}

\begin{env}[19.5.2]
\label{0.19.5.2-fr}
Soient $A$ un anneau topologique, $B$ une $A$-algèbre topologique
(commutative), $\mathfrak J$ un idéal de $B$ (non nécessairement ouvert ni
fermé)~; dans tout ce qui suit, on munit $\mathfrak J$ et les
$\mathfrak J^n$ ($n\geq2$) de la topologie \emph{induite} par celle de $B$,
les quotients $C=B/\mathfrak J$,
$\mathfrak J^n/\mathfrak J^{n+1}=\gr_{\mathfrak J}^n(B)$ de la topologie
\emph{quotient}, de sorte que les $\mathfrak J^n/\mathfrak J^{n+1}$ sont des
$C$-modules topologiques~; l'injection canonique
$\mathfrak J/\mathfrak J^2\to\gr_{\mathfrak J}^*(B)$ se prolonge en un
homomorphisme (d'abord non topologique) de $C$-algèbres
\[
  \varphi:\mathbf S_C^*(\mathfrak J/\mathfrak J^2)
  \longrightarrow\gr_{\mathfrak J}^*(B)
\]
qui pour chaque $n$ donne un homomorphisme \emph{surjectif} de $C$-modules
\begin{equation}
\label{0.19.5.2.1-fr}
  \varphi_n:\mathbf S_C^n(\mathfrak J/\mathfrak J^2)
  \longrightarrow\mathfrak J^n/\mathfrak J^{n+1}
  =\gr_{\mathfrak J}^n(B).
  \tag{19.5.2.1}
\end{equation}

\oldpage[0\textsubscript{IV-1}]{91}

Lorsque $\mathbf S_C^n(\mathfrak J/\mathfrak J^2)$ est muni de la topologie
définie dans (19.5.1), les homomorphismes $\varphi_n$ sont continus, en vertu
de la propriété universelle (19.5.1) de
$\mathbf S_C^*(\mathfrak J/\mathfrak J^2)$ appliquée à la $C$-algèbre
topologique $E=\gr_{\mathfrak J/\mathfrak J^n}(B/\mathfrak J^n)$ munie de la
topologie produit de celles des $\mathfrak J^k/\mathfrak J^{k+1}$, et à
l'injection canonique $\mathfrak J/\mathfrak J^2\to E$. Notons qu'ici la
topologie de $\mathfrak J/\mathfrak J^2$ est moins fine que la topologie
déduite de celle de $C$ (cette dernière topologie sur
$\mathfrak J/\mathfrak J^2$ étant aussi la topologie déduite de celle de $B$).
\end{env}

\begin{theorem}[19.5.3]
\label{0.19.5.3-fr}
Soient $A$ un anneau topologique, $B$ une $A$-algèbre topologique,
$\mathfrak J$ un idéal de $B$, $C=B/\mathfrak J$ la $A$-algèbre topologique
quotient. On suppose que la $A$-algèbre discrète $C$ est formellement lisse.
Alors~:
\begin{enumerate}
  \item[\rm(i)] Si $B$ est une $A$-algèbre formellement lisse,
  $\mathfrak J/\mathfrak J^2$ est un $C$-module topologique formellement
  projectif (19.2.1).
  \item[\rm(ii)] Si $B$ est une $A$-algèbre formellement lisse et un anneau
  préadmissible ($\mathbf 0_{\rm I}$, 7.1.2), les homomorphismes $\varphi_n$
  (19.5.2) sont des bimorphismes formels (19.1.2).
  \item[\rm(iii)] Réciproquement, supposons que $B$ soit préadmissible, que la
  suite $(\mathfrak J^n)$ tende vers $0$ dans $B$, que
  $\mathfrak J/\mathfrak J^2$ soit un $C$-module formellement projectif et que
  les $\varphi_n$ soient des bimorphismes formels. Alors $B$ est une
  $A$-algèbre formellement lisse.
\end{enumerate}
\end{theorem}

La démonstration de ce théorème est longue et encombrée de détails techniques~;
aussi commencerons-nous par en démontrer un corollaire plus simple, où
apparaissent plus nettement les idées directrices~; ce corollaire est d'ailleurs
le cas particulier du théorème (19.5.3) qui sera le plus fréquemment employé
par la suite.

\begin{corollary}[19.5.4]
\label{0.19.5.4-fr}
Soient $A$ un anneau topologique, $B$ une $A$-algèbre topologique,
$\mathfrak J$ un idéal de $B$ tel que la topologie de $B$ soit la topologie
$\mathfrak J$-préadique. On suppose que la $A$-algèbre discrète
$C=B/\mathfrak J$ est formellement lisse. Alors les trois conditions suivantes
sont équivalentes~:
\begin{enumerate}
  \item[a)] $B$ est une $A$-algèbre formellement lisse.
  \item[b)] $\mathfrak J/\mathfrak J^2$ est un $C$-module projectif et
  l'homomorphisme canonique
  \[
    \varphi:\mathbf S_C^*(\mathfrak J/\mathfrak J^2)
    \longrightarrow\gr_{\mathfrak J}^*(B)
  \]
  est bijectif.
  \item[c)] Le séparé complété $\widehat B$ de $B$ est une
  $\widehat A$-algèbre topologique isomorphe à une $\widehat A$-algèbre de la
  forme $\widehat{B'}$, où $B'=\mathbf S_C^*(V)$, $V$ étant un $C$-module
  projectif et $B'$ étant muni de la topologie ${B'}^+$-préadique, où
  ${B'}^+$ est l'idéal d'augmentation.
\end{enumerate}
\end{corollary}

\begin{env}[19.5.4.1]
\label{0.19.5.4.1-fr}
Prouvons d'abord que $a)$ entraîne que $\mathfrak J/\mathfrak J^2$ est un
$C$-module projectif. Soient donc $P$ et $Q$ deux $C$-modules,
$u:P\to Q$ un $C$-homomorphisme surjectif, et
$v:\mathfrak J/\mathfrak J^2\to Q$ un $C$-homomorphisme. L'anneau
$B/\mathfrak J^2$ étant considéré comme une $B$-extension de $C$ par
$\mathfrak J/\mathfrak J^2$, on en déduit par $v$ une $B$-extension
$E=(B/\mathfrak J^2)\oplus_{(\mathfrak J/\mathfrak J^2)}Q$ de $C$ par $Q$
(18.2.8). Comme $C$ est une $A$-algèbre discrète formellement lisse,
l'extension $E$ est $A$-\emph{triviale} (19.4.4) et peut donc s'identifier à
$\mathrm D_C(Q)$ (18.2.3). L'homomorphisme surjectif $u:P\to Q$ définit alors
canoniquement un $A$-homomorphisme surjectif
$\psi:\mathrm D_C(P)\to\mathrm D_C(Q)$ (18.2.9) dont le noyau est un idéal
$\mathfrak L$ de l'extension $E'=\mathrm D_C(P)$, contenu dans $P$, et
\emph{a fortiori} de carré nul. Soit $f:B\to E=E'/\mathfrak L$
l'homomorphisme définissant la structure de $B$-algèbre de $E$~: comme
$\mathfrak L$ est de carré nul, et que $B$ est une $A$-algèbre formellement
lisse, $f$ se factorise en $B\xrightarrow{g}E'\to E'/\mathfrak L$, où $g$ est
un $A$-homomorphisme continu.

\oldpage[0\textsubscript{IV-1}]{92}

Le diagramme
\[
  \xymatrix{
    E' \ar[r]^{\psi} & E \\
    B \ar[u]^{g} \ar[r]_{\theta} & B/\mathfrak J^2 \ar[u]_{v'}
  }
\]
où $v'$ est déduit de $v$ (18.2.8), est commutatif. D'ailleurs l'image de
$\mathfrak J$ par $\psi\circ g$ est contenue dans $Q$, donc l'image de
$\mathfrak J$ par $g$ est contenue dans $P+\mathfrak L=P$, et l'image de
$\mathfrak J^2$ par $g$ est nulle. Restreignant $g$ et $\theta$ à
$\mathfrak J$, on obtient un diagramme commutatif
\[
  \xymatrix{
    P \ar[r]^{u} & Q \\
    \mathfrak J/\mathfrak J^2 \ar[u]^{w} \ar[ur]_{v}
  }
\]
ce qui prouve que le $C$-module $\mathfrak J/\mathfrak J^2$ est projectif.
\end{env}

\begin{env}[19.5.4.2]
\label{0.19.5.4.2-fr}
Prouvons en second lieu que $a)$ entraîne que $\varphi$ est bijectif, ce qui
achèvera de montrer que $a)$ entraîne $b)$. Posons
$E_n=B/\mathfrak J^{n+1}$, et désignons par $F_n$ le quotient de
$\mathbf S_C^*(\mathfrak J/\mathfrak J^2)$ par la puissance $(n+1)$-ième de
son idéal d'augmentation. Comme $\mathfrak J/\mathfrak J^{n+1}$ est nilpotent
dans $E_n$ et que $C$ est une $A$-algèbre discrète formellement lisse,
l'automorphisme identique de $C$ se factorise en
$C\xrightarrow{f}E_n\to C=E_n/(\mathfrak J/\mathfrak J^{n+1})$ où $f$ est un
$A$-homomorphisme~; d'autre part, $\mathfrak J/\mathfrak J^2$ étant un
$C$-module projectif d'après (19.5.4.1), l'automorphisme identique de
$\mathfrak J/\mathfrak J^2$ se factorise en
$\mathfrak J/\mathfrak J^2\xrightarrow{g}\mathfrak J/\mathfrak J^{n+1}
\to\mathfrak J/\mathfrak J^2$, où $g$ est une application $C$-linéaire~; à
partir de $f$ et $g$ on obtient canoniquement un homomorphisme de $C$-algèbres
$\mathbf S_C^*(\mathfrak J/\mathfrak J^2)\to E_n$, qui d'ailleurs (par
définition de $g$) s'annule sur la puissance $(n+1)$-ième de l'idéal
d'augmentation de $\mathbf S_C^*(\mathfrak J/\mathfrak J^2)$, d'où, en
passant au quotient, un $A$-homomorphisme surjectif d'algèbres
$v:F_n\to E_n$ tel que $\gr^0(v)$ et $\gr^1(v)$ soient les automorphismes
identiques de $C$ et de $\mathfrak J/\mathfrak J^2$. Par définition de
l'homomorphisme canonique $\varphi$, on voit que $\gr^j(v)=\varphi_j$ pour tout
$j\leq n$. Notons maintenant que le noyau $\mathfrak R$ de $v$ est un idéal
nilpotent de $F_n$, de sorte que $E_n$ s'identifie à $F_n/\mathfrak R$. Comme
$B$ est une $A$-algèbre formellement lisse, le $A$-homomorphisme canonique
$p_n:B\to E_n=B/\mathfrak J^{n+1}$, qui est continu, se factorise en
$B\xrightarrow{w}F_n\xrightarrow{v}E_n$, où $w$ est un $A$-homomorphisme
continu~; comme $\gr^0(v)$ est l'identité, $w(\mathfrak J)$ est contenu dans
l'idéal d'augmentation de $F_n$, d'où $w(\mathfrak J^{n+1})=0$, si bien que
$w$ se factorise en
$B\xrightarrow{p_n}B/\mathfrak J^n=E_n\xrightarrow{w'}F_n$, et le composé
$E_n\xrightarrow{w'}F_n\xrightarrow{v}E_n$ est l'identité. En outre, comme
$\gr^0(v)$ et $\gr^1(v)$ sont les automorphismes identiques de $C$ et de
$\mathfrak J/\mathfrak J^2$, il en est de même de $\gr^0(w')$ et de
$\gr^1(w')$. Comme $\gr^*(E_n)$ est engendré par $\gr^1(E_n)$,
l'homomorphisme composé
\[
  \gr^j(F_n)\xrightarrow{\varphi_j}\gr^j(E_n)
  \xrightarrow{\gr^j(w')}\gr^j(F_n)
\]
est l'identité pour tout $j\leq n$, puisque c'est vrai pour $j=0$ et $j=1$~;
faisant $j=n$, on prouve ainsi que $\varphi_n$ est injectif, ce qui achève de
montrer que $a)$ entraîne $b)$.
\end{env}

\begin{env}[19.5.4.3]
\label{0.19.5.4.3-fr}
\oldpage[0\textsubscript{IV-1}]{93}
Prouvons ensuite que $b)$ entraîne $a)$. Le même raisonnement qu'au début de
(19.5.4.2) prouve l'existence d'un $A$-homomorphisme surjectif d'algèbres
$v:F_n\to E_n$ tel que $\gr^j(v)=\varphi_j$ pour tout $j$~; comme $\varphi_j$
est bijectif pour tout $j$ et que les filtrations de $F_n$ et $E_n$ sont
finies, on en conclut que $v$ est \emph{bijectif} (Bourbaki,
\emph{Alg. comm.}, chap.~III, §~2, n\textsuperscript{o}~8, cor.~3 du th.~1).
Soient maintenant $G$ une $A$-algèbre topologique discrète, $\mathfrak R$ un
idéal de carré nul dans $G$, $f:B\to G/\mathfrak R$ un $A$-homomorphisme
continu d'algèbres. Comme $G$ est discrète, il existe un entier $m$ tel que
$f$ s'annule dans $\mathfrak J^m$, donc $f$ se factorise en
$B\to E_n\xrightarrow{f_n}G/\mathfrak R$, où l'on prend $n=2m$. On a donc par
composition un $A$-homomorphisme continu d'algèbres
$r:C\to F_n\xrightarrow{v}E_n\xrightarrow{f_n}G/\mathfrak R$, et comme $C$
est une $A$-algèbre formellement lisse discrète, $r$ se factorise en
$C\xrightarrow{r'}G\to G/\mathfrak R$, de sorte que $G$ est muni par $r'$ d'une
structure de $C$-algèbre. D'autre part, la restriction à
$\mathfrak J/\mathfrak J^2$ de l'homomorphisme
$f'_n:F_n\xrightarrow{v}E_n\xrightarrow{f_n}G/\mathfrak R$ est une application
$C$-linéaire $g:\mathfrak J/\mathfrak J^2\to G/\mathfrak R$. Puisque
$\mathfrak J/\mathfrak J^2$ est un $C$-module projectif, $g$ se factorise en
$\mathfrak J/\mathfrak J^2\xrightarrow{h}G\to G/\mathfrak R$, où $h$ est une
application $C$-linéaire~; par prolongement, on en déduit un homomorphisme de
$C$-algèbres $w:\mathbf S_C^*(\mathfrak J/\mathfrak J^2)\to G$, et par
construction tout élément de degré $m$ a pour image par $w$ un élément de
$\mathfrak R$, donc tout élément de degré $n=2m$ a une image nulle puisque
$\mathfrak R^2=0$~; autrement dit, $w$ se factorise en
$\mathbf S_C^*(\mathfrak J/\mathfrak J^2)\to F_n\xrightarrow{w'}G$. Par
construction, le composé $F_n\xrightarrow{w'}G\xrightarrow{\theta}G/\mathfrak R$
coïncide avec $f'_n$ dans $C$ et dans $\mathfrak J/\mathfrak J^2$, donc est
égal à $f'_n$. Finalement, le composé
$B\to E_n\xrightarrow{v^{-1}}F_n\xrightarrow{w'}G\to G/\mathfrak R$ étant
égal à $f$, on voit que $B$ est une $A$-algèbre formellement lisse.
\end{env}

\begin{env}[19.5.4.4]
\label{0.19.5.4.4-fr}
Il reste à prouver l'équivalence de $a)$ et $c)$. En premier lieu, $c)$ entraîne
$a)$~: en effet, $B'$ est une $C$-algèbre formellement lisse pour la topologie
discrète (19.3.2), donc aussi pour la topologie ${B'}^+$-préadique (19.3.8)~;
comme $C$ est une $A$-algèbre formellement lisse, $B'$ est aussi une
$A$-algèbre formellement lisse (19.3.5, (ii)) et finalement $\widehat{B'}$ est
une $\widehat A$-algèbre formellement lisse (19.3.6), donc $B$ est une
$A$-algèbre formellement lisse (19.3.6). Reste à voir que $a)$ entraîne $c)$.
Notons d'abord que puisque $C$ est une $A$-algèbre formellement lisse,
l'homomorphisme identique $C\to B/\mathfrak J$ se factorise en
$C\to\widehat B\to B/\mathfrak J$, où $C\to\widehat B$ est un
$A$-homomorphisme (19.3.10)~; $\widehat B$, et par suite tous les
$B/\mathfrak J^{n+1}$, sont donc munis de structures de $C$-algèbres. D'autre
part, comme $\mathfrak J/\mathfrak J^2$ est un $C$-module projectif par $b)$,
l'injection canonique $\mathfrak J/\mathfrak J^2\to B/\mathfrak J^2$ permet de
former un système projectif de $C$-homomorphismes
$u_n:\mathfrak J/\mathfrak J^2\to B/\mathfrak J^{n+1}$ pour $n\geq1$, donc
aussi (par la propriété universelle de l'algèbre symétrique) un système
projectif d'homomorphismes de $C$-algèbres
$v_n:\mathbf S_C^*(\mathfrak J/\mathfrak J^2)=B'\to B/\mathfrak J^{n+1}$~;
d'ailleurs il est clair que $v_n$ s'annule dans $({B'}^+)^{n+1}$ et comme
$\gr^0(v_n)=\varphi_0$ et $\gr^1(v_n)=\varphi_1$, on a
$\gr^j(v_n)=\varphi_j$ pour $0\leq j\leq n$. Comme les $\varphi_j$ sont des
isomorphismes par $b)$, et que les filtrations de
$B'/({B'}^+)^{n+1}$ et de $B/\mathfrak J^{n+1}$ sont finies, on en conclut que
$v_n$ est \emph{bijectif} pour

\oldpage[0\textsubscript{IV-1}]{94}

tout $n$ (Bourbaki, \emph{Alg. comm.}, chap.~III, §~2,
n\textsuperscript{o}~8, cor.~3 du th.~1)~; d'où $c)$ par passage à la limite
projective.
\end{env}

\begin{remark}[19.5.5]
\label{0.19.5.5-fr}
Sous les hypothèses générales de (19.5.4), supposons en outre que
$\mathfrak J$ soit un idéal \emph{maximal} de $B$, de sorte que
$C=B/\mathfrak J$ soit un corps, et que $\mathfrak J/\mathfrak J^2$ soit un
$C$-espace vectoriel de dimension \emph{finie}. Alors les conditions $a)$,
$b)$ et $c)$ de (19.5.4) sont encore équivalentes à la suivante~:
\begin{enumerate}
  \item[d)] \emph{Étant donnés une algèbre de polynômes}
  $E=C[T_1,\ldots,T_n]$ \emph{sur le corps} $C$, \emph{et en désignant par}
  $\mathfrak n$ \emph{l'idéal maximal de} $E$ \emph{engendré par les} $T_i$
  $(1\leq i\leq n)$, \emph{alors, pour tout idéal} $\mathfrak b$ \emph{de}
  $E$ \emph{contenant une puissance} $\mathfrak n^k$, \emph{tout
  homomorphisme} $v:B\to E/\mathfrak b$ \emph{de} $A$-\emph{algèbres}
  $C$-\emph{augmentées se factorise en}
  $B\xrightarrow{w}E/\mathfrak n^k\to E/\mathfrak b$, \emph{où} $w$ \emph{est
  un} $A$-\emph{homomorphisme}.
\end{enumerate}
En effet, $\mathfrak J/\mathfrak J^2$ étant ici un $C$-module libre, il suffit
de vérifier que la condition $d)$ entraîne que l'homomorphisme canonique
$\varphi:\mathbf S_C^*(\mathfrak J/\mathfrak J^2)\to\gr_{\mathfrak J}^*(B)$
est bijectif. Or ici $\mathbf S_C^*(\mathfrak J/\mathfrak J^2)$ s'identifie à
l'algèbre de polynômes $E=C[T_1,\ldots,T_n]$, où
$n=\operatorname{rg}_C(\mathfrak J/\mathfrak J^2)$, et l'idéal d'augmentation
de $E$ est l'idéal $\mathfrak n$ engendré par les $T_i$. Pour tout entier
$k\geq0$, l'hypothèse que $C$ est une $A$-algèbre discrète formellement lisse
entraîne, comme dans (19.5.4.2), l'existence d'un $A$-homomorphisme surjectif
$f:E/\mathfrak n^{k+1}\to B/\mathfrak J^{k+1}$, tel que
$\gr^j(f)=\varphi_j$ pour tout $j\leq k$. Si
$\mathfrak b/\mathfrak n^{k+1}=\operatorname{Ker}(f)$,
$B/\mathfrak J^{k+1}$ s'identifie donc à $E/\mathfrak b$, et l'hypothèse $d)$
permet de factoriser l'homomorphisme canonique $B\to B/\mathfrak J^{k+1}$ en
$B\xrightarrow{w}E/\mathfrak n^{k+1}\xrightarrow{f}B/\mathfrak J^{k+1}$~;
comme $\gr^0(f)$ est l'identité, on a
$w(\mathfrak J)\subset\mathfrak n/\mathfrak n^{k+1}$, donc
$w(\mathfrak J^{k+1})=0$~; on conclut comme dans (19.5.4.2) que $\varphi_k$
est injectif pour tout $k$, ce qui achève de prouver notre assertion.
\end{remark}

\begin{env}[19.5.6]
\label{0.19.5.6-fr}
\emph{Démonstration du théorème} (19.5.3). --- Soit
$(\mathfrak b_\alpha)$ un système fondamental décroissant d'idéaux ouverts de
$B$. Nous poserons
\[
  B_\alpha=B/\mathfrak b_\alpha,\qquad
  C_\alpha=B/(\mathfrak b_\alpha+\mathfrak J),\qquad
  \mathfrak J_\alpha=(\mathfrak b_\alpha+\mathfrak J)/\mathfrak b_\alpha,
\]
de sorte que
\[
  C_\alpha=B_\alpha/\mathfrak J_\alpha\qquad\text{et}\qquad
  \mathfrak J_\alpha^n/\mathfrak J_\alpha^{n+1}
  =(\mathfrak b_\alpha+\mathfrak J^n)/(\mathfrak b_\alpha+\mathfrak J^{n+1}).
\]
Les $C$-modules
$((\mathfrak b_\alpha\cap\mathfrak J^n)+\mathfrak J^{n+1})/
\mathfrak J^{n+1}$ de $\mathfrak J^n/\mathfrak J^{n+1}$ forment un système
fondamental de sous-modules ouverts du $C$-module topologique
$\mathfrak J^n/\mathfrak J^{n+1}$~; comme
$(\mathfrak b_\alpha\cap\mathfrak J^n)+\mathfrak J^{n+1}
=\mathfrak J^n\cap(\mathfrak b_\alpha+\mathfrak J^{n+1})$, on a
\[
  \mathfrak J^n/((\mathfrak b_\alpha\cap\mathfrak J^n)+\mathfrak J^{n+1})
  =\mathfrak J^n/(\mathfrak J^n\cap(\mathfrak b_\alpha+\mathfrak J^{n+1}))
  =(\mathfrak b_\alpha+\mathfrak J^n)/(\mathfrak b_\alpha+\mathfrak J^{n+1})
  =\mathfrak J_\alpha^n/\mathfrak J_\alpha^{n+1}.
\]
\end{env}

\begin{env}[19.5.6.1]
\label{0.19.5.6.1-fr}
Soient $P$, $Q$ deux $C_\alpha$-modules discrets,
$u:P\to Q$ un $C_\alpha$-homomorphisme surjectif. L'anneau discret
$B/(\mathfrak b_\alpha+\mathfrak J^2)$ est une $B$-extension de $C_\alpha$ par
l'idéal de carré nul $\mathfrak J_\alpha/\mathfrak J_\alpha^2$. Soit
$v:\mathfrak J/\mathfrak J^2\to Q$ un $C$-homomorphisme continu~; remplaçant
au besoin $\mathfrak b_\alpha$ par un idéal ouvert plus petit de $B$, on peut
supposer que le noyau de $v$ contient le sous-$C$-module ouvert
$((\mathfrak b_\alpha\cap\mathfrak J)+\mathfrak J^2)/\mathfrak J^2$~;
passant au quotient, on en déduit un $C_\alpha$-homomorphisme de modules
discrets $v':\mathfrak J_\alpha/\mathfrak J_\alpha^2\to Q$~; soit
$E_\alpha$ la $B_\alpha$-extension de $C_\alpha$ par $Q$ déduite de
$B/(\mathfrak b_\alpha+\mathfrak J^2)$ au moyen de $v'$, et soit
$v'_*:B/(\mathfrak b_\alpha+\mathfrak J^2)\to E_\alpha$ le
$B_\alpha$-homomorphisme correspondant (18.2.8)~; $E_\alpha$ est une
$A$-algèbre topologique discrète, et l'isomorphisme canonique
$C_\alpha\xrightarrow{\sim}E_\alpha/Q$ donne par composition un
$A$-homomorphisme continu $f:C\to C_\alpha\to E_\alpha/Q$. Mais comme $Q$ est
de carré nul dans $E_\alpha$ et que $C$ est une $A$-algèbre formellement
lisse, $f$ se factorise en $C\xrightarrow{g}E_\alpha\to E_\alpha/Q$ où $g$
est un $A$-homomorphisme continu. Comme $g$ est continu et $E_\alpha$ discret,
il existe un $\beta\geq\alpha$ tel que $g$ se factorise en
$C\to C_\beta\xrightarrow{g'}E_\alpha$. D'autre part, soit $E_\beta$ la
$A$-extension de $C_\beta$ par $Q$, image réciproque de $E_\alpha$ par
l'homomorphisme canonique $C_\beta\to C_\alpha$~; l'existence de $g'$ (tel que
le composé $C_\beta\to E_\alpha\to C_\alpha$ soit l'homomorphisme canonique)
équivaut au fait que $E_\beta$ est une extension $A$-\emph{triviale}, et peut
donc s'identifier à $\mathrm D_{C_\beta}(Q)$. Cela étant, l'homomorphisme
surjectif $u:P\to Q$ définit canoniquement un homomorphisme surjectif
$\mathrm D_{C_\beta}(P)\to\mathrm D_{C_\beta}(Q)$ (18.2.9), dont le noyau est
un idéal $\mathfrak L$ de l'extension $E'_\beta=\mathrm D_{C_\beta}(P)$
contenu dans $P$, et \emph{a fortiori} de carré nul. Soit
$f':B\to E_\beta=E'_\beta/\mathfrak L$ le $A$-homomorphisme continu définissant
la structure de $B$-algèbre topologique de $E_\beta$~; comme $\mathfrak L$ est
de carré nul et que $B$ est une $A$-algèbre formellement lisse, $f'$ se
factorise en $B\xrightarrow{h'}E'_\beta\to E'_\beta/\mathfrak L$, où $h'$ est
un $A$-homomorphisme continu. La construction de $E'_\beta$ donne par
composition un $A$-homomorphisme $\psi:E'_\beta\to E_\beta\to E_\alpha$, et
il est clair que le diagramme
\[
  \xymatrix{
    E'_\beta \ar[r]^{\psi} & E_\alpha \\
    B \ar[u]^{h'} \ar[r]_{\theta} &
    B/(\mathfrak b_\alpha+\mathfrak J^2) \ar[u]_{v'_*}
  }
\]
est commutatif. D'ailleurs, l'image par $\psi\circ h'$ de $\mathfrak J$ est
contenue dans $Q$, donc l'image par $h'$ de $\mathfrak J$ est contenue dans
$P$, et l'image de $\mathfrak J^2$ par $h'$ est nulle. Restreignant $h'$ et
$\theta$ à $\mathfrak J$, on obtient un diagramme commutatif
\[
  \xymatrix{
    P \ar[r]^{u} & Q \\
    \mathfrak J/\mathfrak J^2 \ar[u]^{w} \ar[ur]_{v}
  }
\]
où $w$ est continu, ce qui prouve que
$\mathfrak J/\mathfrak J^2$ est un $C$-module formellement projectif.
\end{env}

\oldpage[0\textsubscript{IV-1}]{95}

\begin{env}[19.5.6.2]
\label{0.19.5.6.2-fr}
Pour tout entier $n\geq0$, nous poserons
\[
  E_n=B/\mathfrak J^{n+1},
  \qquad\text{de sorte que}\qquad E_0=C~;
\]
les idéaux $(\mathfrak b_\alpha+\mathfrak J^{n+1})/\mathfrak J^{n+1}$ forment
un système fondamental d'idéaux ouverts dans $E_n$, et nous poserons
\[
  E_{\alpha,n}=B/(\mathfrak b_\alpha+\mathfrak J^{n+1})
  =B_\alpha/\mathfrak J_\alpha^{n+1},
\]
anneau quotient discret. De même, dans
$\gr_{\mathfrak J}^n(B)=\mathfrak J^n/\mathfrak J^{n+1}$, nous avons vu que les
$((\mathfrak b_\alpha\cap\mathfrak J^n)+\mathfrak J^{n+1})/
\mathfrak J^{n+1}$ forment un système fondamental de voisinages de $0$, le
quotient de $\mathfrak J^n/\mathfrak J^{n+1}$ par ce sous-module s'identifiant
canoniquement à $\mathfrak J_\alpha^n/\mathfrak J_\alpha^{n+1}$. Considérons
l'algèbre symétrique
$\mathbf S_{C_\alpha}^*(\mathfrak J_\alpha/\mathfrak J_\alpha^2)$~; nous
noterons $F_{\alpha,n}$ le quotient de
$\mathbf S_{C_\alpha}^*(\mathfrak J_\alpha/\mathfrak J_\alpha^2)$ par la
puissance $(n+1)$-ième de son idéal d'augmentation. Pour un $n\geq1$ fixé, il
résulte de (19.5.1.1) que les
$\mathbf S_{C_\alpha}^n(\mathfrak J_\alpha/\mathfrak J_\alpha^2)$ sont les
quotients de $\mathbf S_C^n(\mathfrak J/\mathfrak J^2)$ par un système
fondamental de sous-modules ouverts dans ce $C$-module topologique.

Pour abréger le langage, nous dirons que pour $\alpha\leq\beta$, les
homomorphismes canoniques $B_\beta\to B_\alpha$, $C_\beta\to C_\alpha$,
$\mathfrak J_\beta/\mathfrak J_\beta^2\to
\mathfrak J_\alpha/\mathfrak J_\alpha^2$, $E_{\beta,n}\to E_{\alpha,n}$,
$F_{\beta,n}\to F_{\alpha,n}$, etc., sont les \emph{homomorphismes de
transition}.
\end{env}

\begin{lemma}[19.5.6.3]
\label{0.19.5.6.3-fr}
Supposons que la $A$-algèbre $C$ soit formellement lisse, l'anneau $B$
préadmissible et le $C$-module $\mathfrak J/\mathfrak J^2$ formellement
projectif. Alors~:
\begin{enumerate}
  \item[\rm(i)] Pour tout $\alpha$, il existe $\beta\geq\alpha$ et un
  $A$-homomorphisme surjectif d'algèbres
  \[
    v_{\alpha\beta}:F_{\beta,n}\longrightarrow E_{\alpha,n}
  \]
  tel que $\gr^0(v_{\alpha\beta}):C_\beta\to C_\alpha$ et
  $\gr^1(v_{\alpha\beta}):\mathfrak J_\beta/\mathfrak J_\beta^2
  \to\mathfrak J_\alpha/\mathfrak J_\alpha^2$ soient les homomorphismes de
  transition.
  \item[\rm(ii)] Si $\beta$ et $v_{\alpha\beta}$ vérifient les conditions de
  \emph{(i)}, alors, pour tout $\gamma\geq\beta$, il existe
  $\delta\geq\gamma$ et un $A$-homomorphisme surjectif d'algèbres
  $v_{\gamma\delta}:F_{\delta,n}\to E_{\gamma,n}$ vérifiant les conditions de
  \emph{(i)} (pour $\gamma$ et $\delta$) et rendant commutatif le diagramme
  \begin{equation}
  \label{0.19.5.6.4-fr}
  \xymatrix@C=3em{
    F_{\beta,n} \ar[r]^{v_{\alpha\beta}} & E_{\alpha,n} \\
    F_{\delta,n} \ar[u] \ar[r]_{v_{\gamma\delta}} & E_{\gamma,n} \ar[u]
  }
  \tag{19.5.6.4}
  \end{equation}
  où les flèches verticales sont les homomorphismes de transition.
\end{enumerate}
\end{lemma}

\begin{env}[19.5.6.4]
\label{0.19.5.6.4-proof-fr}
\begin{enumerate}
  \item[\rm(i)] Dans la $A$-algèbre topologique discrète $E_{\alpha,n}$,
  l'idéal $\mathfrak J_\alpha/\mathfrak J_\alpha^{n+1}$ est nilpotent, et
  l'isomorphisme identique
  $C_\alpha\xrightarrow{\sim}
  E_{\alpha,n}/(\mathfrak J_\alpha/\mathfrak J_\alpha^{n+1})$ donne par
  composition un $A$-homomorphisme continu
  $C\to C_\alpha\to
  E_{\alpha,n}/(\mathfrak J_\alpha/\mathfrak J_\alpha^{n+1})$~;

\oldpage[0\textsubscript{IV-1}]{96}

  comme $C$ est une $A$-algèbre formellement lisse, cet homomorphisme se
  factorise en
  $C\xrightarrow{f_\alpha}E_{\alpha,n}\to
  E_{\alpha,n}/(\mathfrak J_\alpha/\mathfrak J_\alpha^{n+1})$, où $f_\alpha$
  est un $A$-homomorphisme continu~; par suite,
  $\mathfrak J_\alpha/\mathfrak J_\alpha^{n+1}$ devient, au moyen de
  $f_\alpha$, un $C$-module topologique discret annulé par un idéal ouvert de
  $C$. L'hypothèse que $\mathfrak J/\mathfrak J^2$ est un $C$-module
  formellement projectif entraîne alors l'existence d'une application
  $C$-linéaire continue
  $g_\alpha:\mathfrak J/\mathfrak J^2\to
  \mathfrak J_\alpha/\mathfrak J_\alpha^{n+1}$ rendant commutatif le diagramme
  \[
    \xymatrix{
      \mathfrak J_\alpha/\mathfrak J_\alpha^{n+1} \ar[r] &
      \mathfrak J_\alpha/\mathfrak J_\alpha^2 \\
      & \mathfrak J/\mathfrak J^2 \ar[ul]^{g_\alpha} \ar[u]
    }
  \]
  Comme $f_\alpha$ et $g_\alpha$ sont continus, il y a un
  $\beta\geq\alpha$ tel que ces homomorphismes se factorisent respectivement en
  \[
    C\to C_\beta\xrightarrow{f'_{\alpha\beta}}E_{\alpha,n},
  \]
  \[
    \mathfrak J/\mathfrak J^2\to
    \mathfrak J_\beta/\mathfrak J_\beta^2
    \xrightarrow{g'_{\alpha\beta}}
    \mathfrak J_\alpha/\mathfrak J_\alpha^{n+1}
  \]
  et à partir de $f'_{\alpha\beta}$ et $g'_{\alpha\beta}$, on obtient donc
  canoniquement un homomorphisme de $C_\beta$-algèbres
  \[
    \mathbf S_{C_\beta}^*(\mathfrak J_\beta/\mathfrak J_\beta^2)
    \longrightarrow E_{\alpha,n}
  \]
  qui d'ailleurs (par définition de $g'_{\alpha\beta}$) s'annule dans la
  puissance $(n+1)$-ième de l'idéal d'augmentation de
  $\mathbf S_{C_\beta}^*(\mathfrak J_\beta/\mathfrak J_\beta^2)$~; passant au
  quotient par cette puissance $(n+1)$-ième, on obtient l'homomorphisme
  $v_{\alpha\beta}$ cherché, compte tenu des définitions de $f_\alpha$ et
  $g_\alpha$~; la surjectivité de $v_{\alpha\beta}$ résulte en effet de celle
  des deux homomorphismes $\gr^0(v_{\alpha\beta})$ et
  $\gr^1(v_{\alpha\beta})$, puisque cela entraîne que
  $\gr(v_{\alpha\beta})$ est surjectif (l'algèbre
  $\gr^*(E_{\alpha,n})$ étant engendrée par $\gr^0(E_{\alpha,n})$ et
  $\gr^1(E_{\alpha,n})$), et comme les filtrations considérées sont finies, on
  peut appliquer Bourbaki, \emph{Alg. comm.}, chap.~III, §~2,
  n\textsuperscript{o}~8, cor.~3 du th.~1.

  \item[\rm(ii)] L'hypothèse que $B$ est \emph{préadmissible} signifie que
  l'on peut supposer tous les $\mathfrak b_\alpha$ contenus dans un même
  $\mathfrak b_{\alpha_0}$ dont les puissances \emph{tendent vers} $0$. Cela
  entraîne en particulier que le noyau de l'homomorphisme de transition
  $E_{\gamma,n}\to E_{\alpha,n}$, égal à
  $(\mathfrak b_\alpha+\mathfrak J^{n+1})/
  (\mathfrak b_\gamma+\mathfrak J^{n+1})$, est \emph{nilpotent}~; appliquant le
  lemme (19.3.10.1), on voit qu'on peut supposer $f_\gamma$ choisi de sorte que
  le diagramme
  \[
    \xymatrix{
      & E_{\alpha,n} \\
      C \ar[ur]^{f_\alpha} \ar[r]_{f_\gamma} & E_{\gamma,n} \ar[u]
    }
  \]
  (où la flèche verticale est l'homomorphisme de transition) soit commutatif.
  L'hypothèse que $\mathfrak J/\mathfrak J^2$ est un $C$-module formellement
  projectif permet d'autre part de choisir
  $g_\gamma:\mathfrak J/\mathfrak J^2\to
  \mathfrak J_\gamma/\mathfrak J_\gamma^{n+1}$ de sorte que le diagramme
  \[
    \xymatrix@C=5em@R=3em{
      \mathfrak J_\alpha/\mathfrak J_\alpha^{n+1} \ar[r] &
      \mathfrak J_\alpha/\mathfrak J_\alpha^2 \\
      \mathfrak J_\gamma/\mathfrak J_\gamma^{n+1} \ar[u] \ar[r] &
      \mathfrak J_\gamma/\mathfrak J_\gamma^2 \ar[u] \\
      & \mathfrak J/\mathfrak J^2
        \ar[ul]_{g_\gamma} \ar[u]_{p_\gamma} \ar[uul]^{g_\alpha}
    }
  \]
  soit commutatif (il suffit de remarquer que l'image par $(g_\alpha,p_\gamma)$
  de $\mathfrak J/\mathfrak J^2$ dans le module produit
  $(\mathfrak J_\alpha/\mathfrak J_\alpha^{n+1})
  \times(\mathfrak J_\gamma/\mathfrak J_\gamma^2)$ est contenue (vu la
  définition de $g_\alpha$ et de la relation $\beta\leq\gamma$) dans l'image
  canonique $Q$ dans ce module produit du module
  $P=\mathfrak J_\gamma/\mathfrak J_\gamma^{n+1}$, et d'appliquer la définition
  (19.2.1) à l'homomorphisme \emph{surjectif} $P\to Q$ et à l'homomorphisme
  $(g_\alpha,p_\gamma)$ de $\mathfrak J/\mathfrak J^2$ dans $Q$). Ce choix de
  $f_\gamma$ et de $g_\gamma$ permet alors, en construisant
  $v_{\gamma\delta}$ comme dans \emph{(i)}, d'obtenir en outre la commutativité
  du diagramme (19.5.6.4).
\end{enumerate}
\end{env}

\oldpage[0\textsubscript{IV-1}]{97}

\begin{lemma}[19.5.6.5]
\label{0.19.5.6.5-fr}
Supposons que les $A$-algèbres $B$ et $C$ soient formellement lisses et que
$B$ soit préadmissible (de sorte qu'en vertu de (19.5.6.1), les conditions de
(19.5.6.3) sont satisfaites). Pour tout système de deux indices
$\alpha\leq\beta$ et d'un homomorphisme $v_{\alpha\beta}$ vérifiant les
conditions de (19.5.6.3, (i)), il existe un indice $\lambda\geq\beta$ et un
$A$-homomorphisme d'algèbres
\[
  w_{\beta\lambda}:E_{\lambda,n}\longrightarrow F_{\beta,n}
\]
tel que~: $1^\circ$ $\gr^0(w_{\beta\lambda}):C_\lambda\to C_\beta$ et
$\gr^1(w_{\beta\lambda}):
\mathfrak J_\lambda/\mathfrak J_\lambda^2\to
\mathfrak J_\beta/\mathfrak J_\beta^2$ soient les homomorphismes de
transition~; $2^\circ$ le composé
$E_{\lambda,n}\xrightarrow{w_{\beta\lambda}}F_{\beta,n}
\xrightarrow{v_{\alpha\beta}}E_{\alpha,n}$ soit l'homomorphisme de transition.
\end{lemma}

\begin{proof}
Appliquons le lemme (19.5.6.3, (ii)) avec $\gamma=\beta$, ce qui donne un
$\delta\geq\beta$ et un
$v_{\beta\delta}:F_{\delta,n}\to E_{\beta,n}$. Rappelons que
$v_{\beta\delta}$ est surjectif~; d'autre part, son noyau $\mathfrak R$ est
\emph{nilpotent}~: en effet, l'idéal d'augmentation $F_{\delta,n}^+$ de
$F_{\delta,n}$ est nilpotent par définition, et le noyau de
$\gr^0(v_{\beta\delta}):C_\delta\to C_\beta$ est aussi nilpotent en vertu de
l'hypothèse que $B$ est préadmissible. On voit donc que $E_{\beta,n}$
s'identifie à $F_{\delta,n}/\mathfrak R$. Comme $B$ est une $A$-algèbre
formellement lisse, le $A$-homomorphisme canonique
$p_{\beta,n}:B\to E_{\beta,n}$, qui est continu, se factorise en
$B\xrightarrow{w}F_{\delta,n}\xrightarrow{v_{\beta\delta}}E_{\beta,n}$, où
$w$ est un $A$-homomorphisme continu. Comme $F_{\delta,n}$ est discret, il
existe $\lambda\geq\delta$ tel que $w$ soit nul dans $\mathfrak b_\lambda$,
donc $w$ se factorise en
$B\to B/\mathfrak b_\lambda\xrightarrow{w'_\lambda}F_{\delta,n}$. Considérons
l'homomorphisme composé
$B/\mathfrak b_\lambda\xrightarrow{w'_\lambda}F_{\delta,n}
\xrightarrow{q_{\beta\delta}}F_{\beta,n}$, où $q_{\beta\delta}$ est
l'homomorphisme de transition~; notons que
$\gr^0(q_{\beta\delta})=\gr^0(v_{\beta\delta}):C_\delta\to C_\beta$ est
l'homomorphisme de transition~; comme le composé
$v_{\alpha\beta}\circ q_{\beta\delta}\circ w'_\lambda$ est l'homomorphisme
canonique $B/\mathfrak b_\lambda\to
B/(\mathfrak b_\alpha+\mathfrak J^{n+1})$, cela montre que l'image de
$\mathfrak J_\lambda=(\mathfrak b_\lambda+\mathfrak J)/\mathfrak b_\lambda$
par $q_{\beta\delta}\circ w'_\lambda$ est contenue dans l'idéal d'augmentation
$F_{\beta,n}^+$, et par suite l'image de
$\mathfrak J_\lambda^{n+1}=(\mathfrak b_\lambda+\mathfrak J^{n+1})/
\mathfrak b_\lambda$ est nulle. Autrement dit,
$q_{\beta\delta}\circ w'_\lambda$ se factorise en
\[
  B/\mathfrak b_\lambda\to
  B/(\mathfrak b_\lambda+\mathfrak J^{n+1})=E_{\lambda,n}
  \xrightarrow{w_{\beta\lambda}}F_{\beta,n}
\]
tel que le composé
$E_{\lambda,n}\xrightarrow{w_{\beta\lambda}}F_{\beta,n}
\xrightarrow{v_{\alpha\beta}}E_{\alpha,n}$ soit l'homomorphisme de transition.
Le raisonnement précédent montre aussi que $\gr^0(w_{\beta\lambda})$, qui est
le composé
$B/(\mathfrak b_\lambda+\mathfrak J)\to\gr^0(F_{\delta,n})
\xrightarrow{\gr^0(q_{\beta\delta})}\gr^0(F_{\beta,n})$, est l'homomorphisme
de transition $C_\lambda\to C_\beta$, puisque
$v_{\beta\delta}\circ w$ est l'homomorphisme canonique. En outre, on a aussi
$\gr^1(q_{\beta\delta})=\gr^1(v_{\beta\delta})$ (19.5.6.3), donc le même
raisonnement prouve que $\gr^1(w_{\beta\lambda})$ est l'homomorphisme de
transition.
\end{proof}

\begin{env}[19.5.6.6]
\label{0.19.5.6.6-fr}
\emph{Démonstration de} (19.5.3, (ii)). --- Pour montrer que $\varphi_n$ est un
bimorphisme formel, nous allons appliquer le critère de (19.1.3, (iii)). Les
conditions du lemme (19.5.6.5) étant satisfaites par hypothèse, déterminons,
pour tout indice $\alpha$, $v_{\alpha\beta}$ et $w_{\beta\lambda}$ vérifiant
les conclusions de ce lemme. L'homomorphisme
\[
  \gr^n(v_{\alpha\beta}):\gr^n(F_{\beta,n})
  \longrightarrow\gr^n(E_{\alpha,n})
\]
n'est autre que l'homomorphisme
\[
  \varphi_{\alpha\beta,n}:
  \mathbf S_{C_\beta}^n(\mathfrak J_\beta/\mathfrak J_\beta^2)
  \longrightarrow\mathfrak J_\alpha^n/\mathfrak J_\alpha^{n+1}
\]
déduit de l'homomorphisme canonique $\varphi_n$ (19.5.2.1) par passage aux
quotients~; en effet, il résulte de (19.5.6.3) que $\gr^0(v_{\alpha\beta})$ et
$\gr^1(v_{\alpha\beta})$ coïncident respectivement avec
$\varphi_{\alpha\beta,0}$ et $\varphi_{\alpha\beta,1}$, et la définition de
l'homomorphisme canonique $\varphi$ montre alors, par récurrence sur $j\leq n$,
que $\gr^j(v_{\alpha\beta})$ et $\varphi_{\alpha\beta,j}$ sont égaux pour tout
$j$. Cela étant, comme $\varphi_n$ est surjectif, c'est \emph{a fortiori} un
épimorphisme formel~; en outre, pour $j\leq n$, l'homomorphisme composé
\[
  \gr^j(F_{\lambda,n})\xrightarrow{\varphi_{\lambda\lambda,j}}
  \gr^j(E_{\lambda,n})\xrightarrow{\gr^j(w_{\beta\lambda})}
  \gr^j(F_{\beta,n})
\]
est l'homomorphisme de transition, car c'est vrai pour $j=0$ et $j=1$ en
vertu de (19.5.6.5), et comme $\gr^*(E_{\lambda,n})$ est engendré par
$\gr^1(E_{\lambda,n})$ cela démontre l'assertion par récurrence sur $j$.
Composant avec l'homomorphisme de transition
$\gr^j(F_{\beta,n})\to\gr^j(F_{\alpha,n})$, on voit donc, pour $j=n$, que
l'on a factorisé l'homomorphisme de transition
$\gr^n(F_{\lambda,n})\to\gr^n(F_{\alpha,n})$ en
\[
  \gr^n(F_{\lambda,n})\xrightarrow{\varphi_{\lambda\lambda,n}}
  \gr^n(E_{\lambda,n})\longrightarrow\gr^n(F_{\alpha,n}),
\]
ce qui est la condition du critère (19.1.3) pour que $\varphi_n$ soit un
bimorphisme formel.
\end{env}

\oldpage[0\textsubscript{IV-1}]{98}

\begin{lemma}[19.5.6.7]
\label{0.19.5.6.7-fr}
Supposons que $B$ soit préadmissible, que $\mathfrak J/\mathfrak J^2$ soit un
$C$-module formellement projectif, que $C$ soit une $A$-algèbre formellement
lisse et que les $\varphi_n$ soient des bimorphismes formels. Alors, pour tout
couple d'indices $\alpha$, $\beta$ tel que $\alpha\leq\beta$ et tout
homomorphisme $v_{\alpha\beta}:F_{\beta,n}\to E_{\alpha,n}$ vérifiant les
conditions de (19.5.6.3, (i)), il existe un indice $\gamma\geq\beta$ tel que,
pour tout indice $\mu\geq\gamma$, il y ait un diagramme commutatif de
$A$-homomorphismes
\[
\xymatrix@C=4em@R=3em{
  F_{\beta,n}\ar[r]^{v_{\alpha\beta}} & E_{\alpha,n} \\
  & E_{\mu,n}\ar[ul]_{u_{\beta\mu}}\ar[u]_{p_{\alpha\mu}}
}
\]
où $p_{\alpha\mu}$ est l'homomorphisme de transition.
\end{lemma}

\begin{proof}
Appliquant le critère (19.1.3) à chacun des $\varphi_j$ pour $0\leq j\leq n$,
on voit qu'il existe un indice $\gamma\geq\beta$ et des homomorphismes
uniquement déterminés (et surjectifs) de $C_\gamma$-modules
\[
  w_{\beta\gamma}^j:\gr^j(E_{\gamma,n})\longrightarrow
  \gr^j(F_{\beta,n})
\]
tels que les composés
\[
  \gr^j(F_{\gamma,n})
  \xrightarrow{\varphi_{\gamma\gamma,j}}\gr^j(E_{\gamma,n})
  \xrightarrow{w_{\beta\gamma}^j}\gr^j(F_{\beta,n})
\]
soient les homomorphismes de transition (le fait qu'on puisse choisir le même
indice $\gamma$ pour tous les $\varphi_j$ résulte de ce qu'ils sont en nombre
fini). En outre, l'unicité des $w_{\beta\gamma}^j$ prouve (puisque $\varphi$ est
un homomorphisme d'algèbres graduées) que
$w_{\beta\gamma}=(w_{\beta\gamma}^j)_{0\leq j\leq n}$ est un homomorphisme de
$C_\gamma$-algèbres graduées
\[
  \gr^*(E_{\gamma,n})\longrightarrow
  \gr^*(F_{\beta,n})=F_{\beta,n}.
\]
En outre, comme $\varphi_{\gamma\gamma,0}$ et $\varphi_{\gamma\gamma,1}$ sont
les homomorphismes identiques,
$w_{\beta\gamma}^0:C_\gamma\to C_\beta$ et
$w_{\beta\gamma}^1:\mathfrak J_\gamma/\mathfrak J_\gamma^2
\to\mathfrak J_\beta/\mathfrak J_\beta^2$ sont les homomorphismes de
transition, et il en est donc de même de
$\gr^0(v_{\alpha\beta}\circ w_{\beta\gamma})$ et de
$\gr^1(v_{\alpha\beta}\circ w_{\beta\gamma})$~; comme
$\gr^*(E_{\gamma,n})$ est engendré par $\gr^1(E_{\gamma,n})$, on en conclut que
$\gr^j(v_{\alpha\beta}\circ w_{\beta\gamma})$ est aussi l'homomorphisme de
transition pour $0\leq j\leq n$.

Appliquons maintenant à $\alpha$, $\beta$ et $\gamma$ le lemme (19.5.6.3,
(ii)), ce qui donne le diagramme (19.5.6.4) avec $\delta\geq\gamma$~; puis
répétons le raisonnement du début en remplaçant $\alpha$ et $\beta$ par
$\gamma$ et $\delta$. On obtient ainsi un indice $\lambda\geq\delta$ et un
diagramme commutatif d'homomorphismes
\[
\xymatrix@C=4em@R=4em{
  \gr^*(E_{\gamma,n})\ar[r]^{w_{\beta\gamma}} &
  F_{\beta,n}\ar[r]^{v_{\alpha\beta}} & E_{\alpha,n} \\
  \gr^*(E_{\lambda,n})\ar[u]_{\gr^*(p_{\gamma\lambda})}
    \ar[r]_{w_{\delta\lambda}} &
  F_{\delta,n}\ar[u]_{q_{\beta\delta}}\ar[r]_{v_{\gamma\delta}} &
  E_{\gamma,n}\ar[u]_{p_{\alpha\gamma}}
    \ar@{-->}[ul]_{u_{\beta\gamma}}
}
\]
où les flèches verticales sont les homomorphismes de transition. Tout revient
à prouver l'existence de l'homomorphisme $u_{\beta\gamma}$ laissant le
diagramme commutatif, et pour cela il suffit évidemment de montrer que l'on a
\[
  \operatorname{Ker}(v_{\gamma\delta})
  \subset\operatorname{Ker}(q_{\beta\delta}),
\]
$v_{\gamma\delta}$ et $q_{\beta\delta}$ étant surjectifs. Comme
$w_{\delta\lambda}$ est surjectif, cette dernière relation équivaut à
\[
\begin{split}
  \operatorname{Ker}(v_{\gamma\delta}\circ w_{\delta\lambda})
  &\subset\operatorname{Ker}(q_{\beta\delta}\circ w_{\delta\lambda})\\
  &=\operatorname{Ker}\bigl(w_{\beta\gamma}\circ
  (\gr(p_{\gamma\lambda}))\bigr).
\end{split}
\]
Mais on a vu ci-dessus que
$\gr(p_{\gamma\lambda})=\gr(v_{\gamma\delta}\circ w_{\delta\lambda})$ et il
est clair que
\[
\begin{split}
  \operatorname{Ker}(v_{\gamma\delta}\circ w_{\delta\lambda})
  &\subset\operatorname{Ker}\bigl(\gr(v_{\gamma\delta}\circ
  w_{\delta\lambda})\bigr)\\
  &=\operatorname{Ker}(\gr(p_{\gamma\lambda}))
  \subset\operatorname{Ker}\bigl(w_{\beta\gamma}\circ
  (\gr(p_{\gamma\lambda}))\bigr),
\end{split}
\]
ce qui achève la démonstration du lemme, car pour tout $\mu\geq\gamma$, il
suffira de définir $u_{\beta\mu}$ comme le composé
\[
  E_{\mu,n}\xrightarrow{p_{\gamma\mu}}E_{\gamma,n}
  \xrightarrow{u_{\beta\gamma}}F_{\beta,n}.
\]
\end{proof}

\begin{env}[19.5.6.8]
\label{0.19.5.6.8-fr}
\emph{Démonstration de} (19.5.3, (iii)). --- Soient $G$ une $A$-algèbre
topologique discrète, $\mathfrak R$ un idéal de carré nul dans $G$,
$f:B\to G/\mathfrak R$ un $A$-homomorphisme continu d'algèbres. Comme
$G/\mathfrak R$ est discrète, il y a un indice $\alpha$ tel que $f$ s'annule
dans $\mathfrak b_\alpha$~; par hypothèse, il existe un entier $m$ tel que
$\mathfrak J^m\subset\mathfrak b_\alpha$, donc $f$ se factorise en
\[
  B\xrightarrow{p_\alpha}E_{\alpha,n}
  \xrightarrow{f_{\alpha,n}}G/\mathfrak R,
\]
où l'on prend $n=2m$. Les hypothèses du lemme (19.5.6.3) étant satisfaites, on
a tout d'abord un $\beta\geq\alpha$ et un $A$-homomorphisme composé
\begin{equation}
\label{0.19.5.6.9-fr}
  f_{\alpha,n}\circ v_{\alpha\beta}:F_{\beta,n}
  \xrightarrow{v_{\alpha\beta}}E_{\alpha,n}
  \xrightarrow{f_{\alpha,n}}G/\mathfrak R.
  \tag{19.5.6.9}
\end{equation}
Comme $F_{\beta,n}$ est une $C_\beta$-algèbre, donc \emph{a fortiori} une
$C$-algèbre topologique (discrète), cela donne par composition un
$A$-homomorphisme continu
$r:C\to F_{\beta,n}\to G/\mathfrak R$. Comme d'autre part $C$ est une
$A$-algèbre formellement lisse,

\oldpage[0\textsubscript{IV-1}]{99}

cet homomorphisme se factorise en
$r:C\xrightarrow{r'}G\to G/\mathfrak R$, où $r'$ est un
$A$-homomorphisme continu, si bien que $G$ est ainsi muni d'une structure de
$C$-algèbre topologique (discrète). D'autre part, par composition avec
l'homomorphisme canonique
\[
  \mathfrak J/\mathfrak J^2\longrightarrow
  \mathfrak J_\beta/\mathfrak J_\beta^2\longrightarrow F_{\beta,n},
\]
on déduit de (19.5.6.9) un $C$-homomorphisme continu
$g:\mathfrak J/\mathfrak J^2\to G/\mathfrak R$. Comme
$\mathfrak J/\mathfrak J^2$ est un $C$-module formellement projectif, $g$ se
factorise en
\[
  \mathfrak J/\mathfrak J^2\xrightarrow{h}G\longrightarrow G/\mathfrak R,
\]
où $h$ est un $C$-homomorphisme continu. Comme $G$ est discret, il existe un
indice $\gamma\geq\beta$ tel que l'image par $h$ de
$((\mathfrak b_\gamma\cap\mathfrak J)+\mathfrak J^2)/\mathfrak J^2$ soit
nulle, ainsi que l'image par $r'$ de
$(\mathfrak b_\gamma+\mathfrak J)/\mathfrak J$, de sorte que $h$ se factorise
en
\[
  \mathfrak J/\mathfrak J^2\longrightarrow
  \mathfrak J_\gamma/\mathfrak J_\gamma^2
  \xrightarrow{h'}G,
\]
où $h'$ est un $C_\gamma$-homomorphisme. Par prolongement, on déduit donc de
$h'$ un homomorphisme de $C_\gamma$-algèbres
\[
  w:\mathbf S_{C_\gamma}^*(\mathfrak J_\gamma/\mathfrak J_\gamma^2)
  \longrightarrow G
\]
et par construction, tout élément de degré $m$ a pour image par $w$ un élément
de $\mathfrak R$, donc tout élément de degré $n=2m$ a une image nulle, puisque
$\mathfrak R^2=0$~; autrement dit, $w$ se factorise en
\[
  \mathbf S_{C_\gamma}^*(\mathfrak J_\gamma/\mathfrak J_\gamma^2)
  \longrightarrow F_{\gamma,n}\xrightarrow{w'_\gamma}G.
\]
Appliquons maintenant à
\[
  v_{\alpha\gamma}:F_{\gamma,n}
  \xrightarrow{q_{\beta\gamma}}F_{\beta,n}
  \xrightarrow{v_{\alpha\beta}}E_{\alpha,n}
\]
le lemme (19.5.6.7), dont les hypothèses sont vérifiées~; il existe donc un
$\delta\geq\gamma$ et un homomorphisme
$u_{\gamma\delta}:E_{\delta,n}\to F_{\gamma,n}$ tel que le composé
$v_{\alpha\gamma}\circ u_{\gamma\delta}$ soit l'homomorphisme de transition
$p_{\alpha\delta}:E_{\delta,n}\to E_{\alpha,n}$. On obtient donc finalement un
diagramme commutatif de $A$-homomorphismes continus
\[
\xymatrix@C=3em@R=2.5em{
  && E_{\alpha,n}\ar[r]^{f_{\alpha,n}} & G/\mathfrak R \\
  && F_{\beta,n}\ar[u]_{v_{\alpha\beta}} & \\
  B\ar[r]_{p_\delta} & E_{\delta,n}\ar[uur]^{p_{\alpha\delta}}
    \ar[r]_{u_{\gamma\delta}} &
  F_{\gamma,n}\ar[u]_{q_{\beta\gamma}}\ar[r]^{w'_\gamma} & G\ar[uu]
}
\]
et le composé $f':B\to G$ des homomorphismes de la ligne inférieure est donc
tel que $f$ se factorise en $B\xrightarrow{f'}G\to G/\mathfrak R$, ce qui
prouve que $B$ est une $A$-algèbre formellement lisse.

La démonstration du th.~(19.5.3) est ainsi achevée.
\end{env}

\begin{corollary}[19.5.7]
\label{0.19.5.7-fr}
Soient $A$ un anneau topologique, $B$ une $A$-algèbre topologique,
$(\mathfrak b_\lambda)$ un système fondamental d'idéaux ouverts dans $B$,
$\mathfrak J$ un idéal de $B$, $C=B/\mathfrak J$ la $A$-algèbre topologique
quotient. On pose $C_\lambda=B/(\mathfrak b_\lambda+\mathfrak J)$. On suppose
que~: $1^\circ$ pour tout $n$, la topologie induite sur $\mathfrak J^n$ par
celle de $B$ est aussi la topologie du $B$-module $\mathfrak J^n$ déduite de
la topologie de $B$ (19.0.2) (cette condition sera remplie en particulier si
$B$ est noethérien et sa topologie préadique ($\mathbf 0_{\mathrm I}$,
7.3.2))~; $2^\circ$ $C$ est une $A$-algèbre formellement lisse. Sous ces
conditions~:
\begin{enumerate}
  \item[\rm(i)] Si $B$ est une $A$-algèbre formellement lisse, alors, pour tout
  $\lambda$, $(\mathfrak J/\mathfrak J^2)\otimes_C C_\lambda$ est un
  $C_\lambda$-module projectif.
  \item[\rm(ii)] Si $B$ est une $A$-algèbre formellement lisse et un anneau
  préadmissible, alors pour tout $\lambda$, l'homomorphisme canonique
  \begin{equation}
  \label{0.19.5.7.1-fr}
    \varphi_\lambda=\varphi\otimes 1_{C_\lambda}:
    \mathbf S_{C_\lambda}^*((\mathfrak J/\mathfrak J^2)\otimes_C C_\lambda)
    \longrightarrow\gr_{\mathfrak J}^*(B)\otimes_C C_\lambda
    \tag{19.5.7.1}
  \end{equation}
  est bijectif.

  \oldpage[0\textsubscript{IV-1}]{100}

  \item[\rm(iii)] Réciproquement, supposons que $B$ soit préadmissible, que la
  suite $(\mathfrak J^n)$ tende vers $0$, et que, pour tout $\lambda$,
  $(\mathfrak J/\mathfrak J^2)\otimes_C C_\lambda$ soit un $C_\lambda$-module
  projectif et l'homomorphisme (19.5.4.1) soit bijectif. Alors $B$ est une
  $A$-algèbre formellement lisse.
\end{enumerate}
\end{corollary}

\begin{proof}
En effet, la topologie de $\mathfrak J/\mathfrak J^2$ et celle des
$\mathbf S_C^n(\mathfrak J/\mathfrak J^2)$ sont alors aussi déduites de celle
de $B$ (19.5.1)~; les conclusions sont alors des conséquences immédiates de
(19.5.3) et des critères de (19.1.4) et (19.2.4) spéciaux à ce type de
topologies.
\end{proof}

\begin{remark}[19.5.8]
\label{0.19.5.8-fr}
Supposons que $\mathfrak J/\mathfrak J^2$ soit un $C$-module de type fini et
que, pour la topologie quotient de celle de $B$, $C$ soit un anneau de Zariski~;
soit $\mathfrak r$ un idéal de définition de $C$, de sorte que les
$\mathfrak r^n$ forment un système fondamental de voisinages de $0$ dans $C$.
Alors les conclusions de (i) et (ii) dans (19.5.7) peuvent être remplacées par
les suivantes~:
\begin{enumerate}
  \item[\rm(i')] $\mathfrak J/\mathfrak J^2$ est un $C$-module projectif.
  \item[\rm(ii')] L'homomorphisme canonique
  $\varphi:\mathbf S_C^*(\mathfrak J/\mathfrak J^2)
  \to\gr_{\mathfrak J}^*(B)$ est bijectif.
\end{enumerate}
En effet, il est clair que (i') entraîne la conclusion de (i) dans (19.5.7).
Inversement, si $(\mathfrak J/\mathfrak J^2)\otimes_C C_\lambda$ est un
$C_\lambda$-module projectif pour tout $\lambda$, alors
$(\mathfrak J/\mathfrak J^2)\otimes_C(C/\mathfrak r^n)$ est un
$(C/\mathfrak r^n)$-module projectif (donc plat) pour tout $n$~; on en conclut
que $\mathfrak J/\mathfrak J^2$ est un $C$-module plat
($\mathbf 0_{\mathrm{III}}$, 10.2.2), donc projectif puisqu'il est de
présentation finie (Bourbaki, \emph{Alg. comm.}, chap.~II, \textsection5,
n\textsuperscript{o}~2, cor.~2 du th.~1). D'autre part, les $C$-modules
$\mathbf S_C^n(\mathfrak J/\mathfrak J^2)$ et $\gr_{\mathfrak J}^n(B)$ sont
de type fini, et on sait que lorsque $C$ est un anneau de Zariski, il revient
au même de dire alors que $\varphi_n$ est bijectif ou que
$\widehat\varphi_n$ est bijectif (Bourbaki, \emph{Alg. comm.}, chap.~III,
\textsection3, n\textsuperscript{o}~5, prop.~9), donc (ii) est équivalent à
(ii').
\end{remark}

\subsection{Cas des algèbres sur un corps}
\label{subsection:0.19.6-fr}

\begin{theorem}[19.6.1) (Cohen]
\label{0.19.6.1-fr}
Soient $k$ un corps, $K$ une extension de $k$, $k$ et $K$
étant munis des topologies discrètes. Pour que $K$ soit une $k$-algèbre
formellement lisse, il faut et il suffit que $K$ soit une extension séparable
de $k$.
\end{theorem}

\begin{proof}
La nécessité de la condition sera établie en (19.6.5.1) (et naturellement ne
sera pas utilisée jusque-là)~; nous nous bornerons ici à prouver que la
condition est \emph{suffisante}. Distinguons deux cas~:

\emph{I. --- $K$ est une extension séparable de type fini de $k$.} On sait
alors (Bourbaki, \emph{Alg.}, chap.~V, \textsection9,
n\textsuperscript{o}~3, th.~2) qu'il existe une sous-extension pure
$K'=k(T_1,\ldots,T_n)$ de $K$ telle que $K$ soit une extension algébrique
finie séparable de $K'$. Compte tenu de (19.3.5, (iii)), on peut donc se
borner, soit au cas où $K=K'$, soit au cas où $K$ est algébrique fini sur $k$.
Dans le premier cas, on sait que $A=k[T_1,\ldots,T_n]$ est une $k$-algèbre
formellement lisse (19.3.3), donc il en est de même de
$k(T_1,\ldots,T_n)=S^{-1}A$, où $S=A-\{0\}$ (19.3.5, (iv)). Dans le second
cas, \emph{tous} les groupes de cohomologie de Hochschild
$\mathrm H_k^n(K,L)$, pour un $(K,K)$-bimodule quelconque $L$, sont nuls~: en
effet, si l'on considère la $k$-algèbre produit tensoriel $C=K\otimes_k K$,
$L$ est un $C$-module à gauche et le groupe de cohomologie
$\mathrm H_k^n(K,L)$ est égal à $\operatorname{Ext}_C^n(K,L)$, où $K$ est
aussi considéré comme $(K,K)$-bimodule,

\oldpage[0\textsubscript{IV-1}]{101}

donc comme $C$-module à gauche (M, IX, 4). Or, comme $K$ est une extension
finie séparable de $k$, on sait que $K\otimes_k K$ est composée directe
d'extensions de $k$, dont l'une est $K$ lui-même (Bourbaki, \emph{Alg.},
chap.~VIII, \textsection8, prop.~3)~; cela entraîne donc que $K$ est un
$C$-module \emph{projectif}, d'où notre assertion. Toute $k$-extension de $K$
de noyau $L$ est donc $k$-triviale, et \emph{a fortiori} les $k$-extensions
commutatives le sont, d'où le théorème dans ce cas (19.4.4).

\emph{II. --- Cas général.} --- Avec les notations de (18.4.5), il s'agit de
prouver que $\operatorname{Hom}_K(\mathrm H_2(P'_\bullet),L)=0$ pour
\emph{tout} $K$-espace vectoriel $L$, ce qui signifie évidemment que
$\mathrm H_2(P'_\bullet)=0$. Si $K$ est réunion d'une famille filtrante de
sous-extensions $K^{(\alpha)}$ de $k$, $P'_\bullet$ est \emph{limite
inductive} des complexes correspondants $P'^{(\alpha)}_\bullet$, puisque le
foncteur $\varinjlim$ commute au produit tensoriel dans la catégorie des
$k$-modules~; par l'exactitude du foncteur $\varinjlim$ dans cette catégorie,
on a donc
$\mathrm H_2(P'_\bullet)=\varinjlim\mathrm H_2(P'^{(\alpha)}_\bullet)$.
Comme l'hypothèse que $K$ est séparable entraîne qu'il en est de même de toute
sous-extension de $K$ et que $K$ est réunion des sous-extensions de type fini,
la première partie de la démonstration entraîne bien que $K$ est une
$k$-algèbre formellement lisse. C.Q.F.D.
\end{proof}

\begin{corollary}[19.6.2]
\label{0.19.6.2-fr}
Soient $A$ un anneau local séparé et complet, $K$ son corps résiduel, $k$ un
sous-corps de $K$ tel que $K$ soit une extension séparable de $k$. Alors il
existe un sous-corps $K'$ de $A$ contenant $k$, tel que la restriction à $K'$
de l'homomorphisme canonique $A\to K$ soit un isomorphisme de $K'$ sur $K$.
\end{corollary}

\begin{proof}
Soit $\mathfrak m$ l'idéal maximal de $A$. En vertu de l'hypothèse et de
(19.6.1), $K$ est une $k$-algèbre formellement lisse~; appliquant (19.3.10) en
remplaçant $C$ par $A$ et $\mathfrak J$ par $\mathfrak m$, on voit que
l'automorphisme identique de $K=A/\mathfrak m$ se factorise en
$K\xrightarrow{u}A\to A/\mathfrak m$, où $u$ est un $k$-homomorphisme, donc
nécessairement un $k$-isomorphisme de $K$ sur un sous-corps $K'$ contenant
$k$.
\end{proof}

\begin{corollary}[19.6.3]
\label{0.19.6.3-fr}
Soient $A$ un anneau local noethérien complet, $\mathfrak m$ son idéal maximal,
$k$ son corps résiduel. Les conditions suivantes sont équivalentes~:
\begin{enumerate}
  \item[a)] $A$ contient un sous-corps.
  \item[b)] Si $p$ est la caractéristique de $k$, on a $pA=0$.
  \item[c)] Il existe un corps $K$ tel que $A$ soit isomorphe à un anneau
  quotient d'un anneau de séries formelles $B=K[[T_1,\ldots,T_n]]$.
\end{enumerate}
Lorsqu'il en est ainsi, on peut supposer que $A$ est isomorphe à $B/\mathfrak b$,
où $\mathfrak b$ est contenu dans le carré de l'idéal maximal de $B$.
\end{corollary}

\begin{proof}
Il est immédiat que $c)$ implique $a)$, car $A$ est $\neq0$ puisque c'est un
anneau local~; comme $K$ est un corps et que l'homomorphisme composé
$K\xrightarrow{\varphi}B\to A$ (où $\varphi$ est l'injection canonique) n'est
pas nul, il est injectif. Pour voir que $a)$ entraîne $b)$, il suffit de
remarquer que si $K$ est un sous-corps de $A$, $K$ est isomorphe à un
sous-corps de $k$ et a par suite même caractéristique $p$~; comme $p.1=0$ dans
$K$, donc dans $A$, on a $pA=0$. Réciproquement, $b)$ entraîne $a)$, car le
composé des homomorphismes canoniques
$\mathbf Z\xrightarrow{j}A\to k$ a pour noyau $p\mathbf Z$, et l'on a
$j(p\mathbf Z)=0$, donc $\operatorname{Ker}(j)=p\mathbf Z$ et $A$ contient un
anneau $A'$ isomorphe à $\mathbf Z/p\mathbf Z$ et ne rencontrant pas
$\mathfrak m$~; si $p>0$, $A'$ est un corps~; sinon, tout

\oldpage[0\textsubscript{IV-1}]{102}

élément de $A'$ étant inversible dans $A$, le corps des fractions de $A'$
(isomorphe à $\mathbf Q$) est aussi contenu dans $A$.

Prouvons enfin que $a)$ entraîne $c)$~; la condition $a)$ entraîne que $A$
contient un sous-corps premier $k_0$, isomorphe au sous-corps premier de $k$,
et dont $k$ est par suite extension séparable. Appliquant (19.6.2), on voit
d'abord qu'il existe un isomorphisme $f$ de $k$ sur un sous-corps $K$ de $A$.
D'autre part, soit $(x_i)_{1\leq i\leq n}$ un système d'éléments de
$\mathfrak m$ tel que les classes mod.~$\mathfrak m^2$ des $x_i$ forment une
base (sur $k$) de $\mathfrak m/\mathfrak m^2$. Puisque $A$ est complet, il y a
alors un homomorphisme continu d'anneaux
$u:k[[T_1,\ldots,T_n]]\to A$ tel que $u$ soit égal à $f$ dans $k$ et que
$u(T_i)=x_i$ pour tout $i$, et cet homomorphisme est surjectif en vertu du
choix des $x_i$ (Bourbaki, \emph{Alg. comm.}, chap.~III, \textsection2,
n\textsuperscript{o}~9, prop.~11).
\end{proof}

\begin{theorem}[19.6.4]
\label{0.19.6.4-fr}
Soient $k$ un corps, $A$ une $k$-algèbre locale noethérienne, $\mathfrak m$ son
idéal maximal, $K=A/\mathfrak m$ son corps résiduel, $A$ étant muni de la
topologie $\mathfrak m$-préadique. On suppose que $K$ soit une extension
séparable de $k$. Alors les conditions suivantes sont équivalentes~:
\begin{enumerate}
  \item[a)] $A$ est une $k$-algèbre formellement lisse.
  \item[b)] Le complété $\widehat A$ de $A$ est $k$-isomorphe à un anneau de
  séries formelles $K[[T_1,\ldots,T_n]]$ (dont la structure de $k$-algèbre est
  définie par l'homomorphisme composé
  \[
    k\xrightarrow{\varphi}K\xrightarrow{\psi}K[[T_1,\ldots,T_n]],
  \]
  où $\psi$ est l'injection canonique et $\varphi$ l'homomorphisme définissant
  la structure d'extension de $K$).
  \item[b')] $\widehat A$ est un anneau isomorphe à un anneau de séries
  formelles $K[[T_1,\ldots,T_n]]$.
  \item[c)] $A$ est un anneau local régulier.
\end{enumerate}
\end{theorem}

\begin{proof}
Le fait que $a)$ entraîne $b)$ est le cas particulier de ((19.5.4),
équivalence de $a)$ et $c)$), appliqué en y remplaçant $A$ par $k$, $B$ par
$A$ et $\mathfrak J$ par $\mathfrak m$, compte tenu de (19.6.1). Il est clair
que $b)$ entraîne $b')$ et $b')$ entraîne $c)$ par (17.1.1). Enfin si $c)$ est
vérifiée, il résulte de (17.1.1, $a)$) et de ((19.5.4), équivalence de $b)$ et
$a)$), appliqué à $K=A/\mathfrak m$ remplaçant $C$, que $A$ est une
$k$-algèbre formellement lisse.
\end{proof}

\begin{corollary}[19.6.5]
\label{0.19.6.5-fr}
Soient $k$ un corps, $A$ une $k$-algèbre locale noethérienne, $\mathfrak m$ son
idéal maximal, $A$ étant muni de la topologie $\mathfrak m$-préadique.
Supposons que $A$ soit une $k$-algèbre formellement lisse~; alors $A$ est
géométriquement régulier sur $k$, autrement dit (cf.~$\mathbf{IV}$, 6.7.6),
pour toute extension finie $k'$ de $k$, l'anneau semi-local
$A'=A\otimes_k k'$ est régulier (17.3.6). En outre, si $K$ est le corps
résiduel de $A$, $\widehat A$ est isomorphe à un anneau de séries formelles
$K[[T_1,\ldots,T_n]]$.
\end{corollary}

\begin{proof}
Comme $\widehat A'=\widehat A\otimes_k k'$, il résulte de (19.3.6) et de
(17.1.5) (appliqué aux anneaux locaux en les idéaux maximaux de $A'$) qu'on
peut se borner au cas où $A$ est complet~; alors $A'$ est une $k'$-algèbre
formellement lisse (19.3.5, (iii)) et par ailleurs est composée directe de
$k'$-algèbres locales, qui sont aussi formellement lisses (19.3.5, (v)). On est
donc ramené à prouver que $A$ est \emph{régulier}. Soit $k_0$ le sous-corps
premier de $k$~; comme $k$ est une extension séparable de $k_0$, c'est une
$k_0$-algèbre formellement lisse (19.6.1) et en vertu de l'hypothèse, il en est
de même de $A$ (19.3.5, (ii)). Comme le corps résiduel $K$ de $A$ est une
extension séparable de $k_0$, on peut alors appliquer (19.6.4) à $A$ et $k_0$,
d'où la conclusion.
\end{proof}

\oldpage[0\textsubscript{IV-1}]{103}

\begin{corollary}[19.6.5.1]
\label{0.19.6.5.1-fr}
Soient $A$ un anneau local noethérien, $k$ un sous-corps de $A$ tel que $A$
soit une $k$-algèbre formellement lisse (pour sa topologie préadique). Alors
tout corps $K$ tel que $k\subset K\subset A$ est une extension séparable de
$k$.
\end{corollary}

\begin{proof}
En effet, pour toute extension finie $k'$ de $k$, l'anneau $K\otimes_k k'$ 
s'identifie à un sous-anneau de $A'=A\otimes_k k'$~; comme $A'$ est un anneau
régulier, il est réduit, donc il en est de même de $K$, ce qui prouve que $K$
est une extension séparable de $k$ (Bourbaki, \emph{Alg.}, chap.~VIII,
\textsection7, n\textsuperscript{o}~3, th.~1).
\end{proof}

On notera que cela démontre que la condition de l'énoncé de (19.6.1) est
\emph{nécessaire}.

\begin{remark}[19.6.5.2]
\label{0.19.6.5.2-fr}
Si $K$ est une extension \emph{non séparable} de $k$, l'anneau de séries
formelles $K[[T_1,\ldots,T_n]]$ \emph{n'est donc pas} une $k$-algèbre
formellement lisse (pour la structure de $k$-algèbre usuelle rappelée dans
(19.6.4)). D'autre part, il y a des $k$-algèbres formellement lisses $A$ qui
sont des anneaux locaux noethériens complets dont le corps résiduel $K$ est une
extension \emph{non séparable} de $k$ (par exemple les complétés des localisés
de $B=k[T_1,\ldots,T_n]$ en des idéaux maximaux $\mathfrak n$ tels que
$B/\mathfrak n$ soit une extension algébrique finie non séparable de $k$)~; un
tel anneau ne peut donc être $k$-isomorphe à $K[[T_1,\ldots,T_n]]$ bien qu'il
lui soit isomorphe en vertu de (19.6.5).

On peut dans certains cas préciser le corollaire (19.6.5)~:
\end{remark}

\begin{proposition}[19.6.6]
\label{0.19.6.6-fr}
Soient $k$ un corps, $A$ une $k$-algèbre locale noethérienne, $\mathfrak m$ son
idéal maximal, $A$ étant muni de la topologie $\mathfrak m$-préadique.
Supposons que le corps résiduel $K=A/\mathfrak m$ de $A$ soit tel qu'il existe
une extension radicielle finie $k_0$ de $k$ pour laquelle le corps résiduel de
l'anneau local $K\otimes_k k_0$ soit une extension séparable de $k_0$ (nous
exprimerons plus tard cette propriété en disant que $K$ est de
\emph{multiplicité radicielle finie} sur $k$ ($\mathbf{IV}$, 4.7.4) et nous
verrons que cette condition est toujours satisfaite lorsque $K$ est une
\emph{extension de type fini} de $k$). Alors les conditions suivantes sont
équivalentes~:
\begin{enumerate}
  \item[a)] $A$ est une $k$-algèbre formellement lisse.
  \item[b)] Il existe une extension radicielle finie $k'$ de $k$ telle que
  $\widehat A\otimes_k k'$ soit $k'$-isomorphe à une algèbre de séries
  formelles $K'[[T_1,\ldots,T_n]]$, où $K'$ est une extension séparable de
  $k'$.
  \item[b')] Il existe une extension finie $k'$ de $k$ telle que l'anneau
  semi-local $\widehat A\otimes_k k'$ ait au moins un anneau local composant
  qui soit $k'$-isomorphe à une algèbre de séries formelles
  $K'[[T_1,\ldots,T_n]]$, où $K'$ est une extension séparable de $k'$.
  \item[c)] $A$ est géométriquement régulier sur $k$ (19.6.5).
\end{enumerate}
\end{proposition}

\begin{proof}
Notons d'abord que si $k'$ est une extension radicielle de $k$, il n'y a qu'un
seul idéal de $A'=A\otimes_k k'$ au-dessus de $\mathfrak m$, formé des éléments
dont une puissance $p^h$-ième ($p$ exposant caractéristique de $k$) est dans
$\mathfrak m$ pour un $h$ convenable (Bourbaki, \emph{Alg. comm.}, chap.~V,
\textsection2, n\textsuperscript{o}~3, lemme~4)~; $A'$ est donc un anneau
local, et il en est de même de
$K\otimes_k k'=(A\otimes_k k')/(\mathfrak m\otimes_k k')$~; en outre les corps
résiduels de ces deux anneaux sont identiques. Rappelons d'autre part que si
$K$ est une extension séparable de $k$, alors, pour toute extension finie
$k''$ de $k$, $K\otimes_k k''$ est composé direct de corps (Bourbaki,
\emph{Alg.}, chap.~VIII, \textsection7, n\textsuperscript{o}~3, cor.~1 du
th.~1), et par suite $\mathfrak m\otimes_k k''$ est le radical de
$A\otimes_k k''$, et les corps composants de $K\otimes_k k''$ sont les corps
résiduels aux idéaux maximaux de $A\otimes_k k''$~; en

\oldpage[0\textsubscript{IV-1}]{104}

outre, ces corps sont des extensions séparables de $k''$ puisque pour toute
extension $k_1$ de $k''$,
$K\otimes_k k_1=(K\otimes_k k'')\otimes_{k''}k_1$ est sans radical
(\emph{loc. cit.}, th.~1). Si l'on applique ces remarques au corps $K$ de
multiplicité radicielle finie sur $k$, on voit que pour \emph{toute} extension
finie $k''$ de $k_0$, $(K\otimes_k k'')_{\mathrm{red}}$ est séparable sur
$k''$.

Passons à la démonstration de (19.6.6). Il est trivial que $b)$ implique
$b')$. Montrons que $b')$ entraîne $a)$. Si l'on pose
$B'=K'[[T_1,\ldots,T_n]]$, l'hypothèse que $K'$ est séparable sur $k'$
entraîne, par les remarques du début, que pour toute extension finie $k''$ de
$k'$, les composants de l'anneau semi-local complet $B'\otimes_{k'}k''$
(égal à l'anneau de séries formelles
$(K'\otimes_{k'}k'')[[T_1,\ldots,T_n]]$) sont les anneaux
$K_i''[[T_1,\ldots,T_n]]$, où les $K_i''$ sont les corps composants de
$K'\otimes_{k'}k''$. Comme les composants locaux de $B'\otimes_{k'}k''$ sont
aussi des composants locaux de
\[
  \widehat A\otimes_k k''
  =(\widehat A\otimes_k k')\otimes_{k'}k'',
\]
on voit que l'hypothèse $b')$ est encore vérifiée lorsqu'on remplace $k'$ par
une quelconque de ses extensions finies. En particulier, on peut supposer que
$k'$ est \emph{quasi-galoisienne}, donc extension galoisienne d'une extension
radicielle $k'_0$ de $k$~; $A'=\widehat A\otimes_k k'$ peut alors s'écrire
$A'_0\otimes_{k'_0}k'$, où $A'_0=\widehat A\otimes_k k'_0$ est un anneau
local complet. On sait alors (Bourbaki, \emph{Alg.}, chap.~VIII,
\textsection8, prop.~4) que $A'$ est composé direct d'anneaux locaux
\emph{isomorphes} à $A'_0$ en tant que $k'_0$-algèbres. Or, un de ces anneaux
est par hypothèse $k'$-isomorphe à un anneau de séries formelles
$K'[[T_1,\ldots,T_n]]$, où $K'$ est une extension séparable de $k'$, donc
\emph{aussi} une extension séparable de $k'_0$~; compte tenu de (19.6.4),
$A'_0$ est donc une $k'_0$-algèbre formellement lisse. Mais comme $k'_0$ est
un $k$-module fidèlement plat, projectif et de type fini, on conclut de
(19.4.7) que $\widehat A$ est une $k$-algèbre formellement lisse.

On a déjà vu (19.6.5) que $a)$ implique $c)$~; il reste donc à vérifier que
$c)$ implique $b)$. Or, si $c)$ est vérifiée, alors, pour toute extension
radicielle finie $k'$ de $k$ contenant $k_0$, $A\otimes_k k'$ est un anneau
local régulier, dont le corps résiduel $K'$ est une extension séparable de
$k'$~; il résulte donc de (19.6.4) que son complété
$\widehat A\otimes_k k'$ est $k'$-isomorphe à un anneau de séries formelles
$K'[[T_1,\ldots,T_n]]$.
\end{proof}

\begin{remarks}[19.6.7]
\label{0.19.6.7-fr}
\begin{enumerate}
  \item[\rm(i)] On notera que l'hypothèse que $K$ est de multiplicité
  radicielle finie sur $k$ n'a été utilisée que dans la démonstration de
  l'implication $b')\Rightarrow a)$.
  \item[\rm(ii)] Nous prouverons plus tard (22.5.8) que $a)$ et $c)$ sont
  équivalentes, \emph{sans hypothèse} sur l'extension $K$ de $k$.
\end{enumerate}
\end{remarks}

\subsection{Cas des homomorphismes locaux; théorèmes d'existence et
d'unicité}
\label{subsection:0.19.7-fr}

Dans ce numéro, lorsqu'un anneau semi-local est considéré comme un anneau
topologique, il est toujours sous-entendu qu'il s'agit de sa topologie
$\mathfrak r$-préadique, où $\mathfrak r$ est son radical. Tout homomorphisme
local d'anneaux locaux est donc automatiquement \emph{continu}.

\begin{theorem}[19.7.1]
\label{0.19.7.1-fr}
Soient $A$, $B$ deux anneaux locaux noethériens, $\mathfrak m$, $\mathfrak n$
leurs idéaux maximaux respectifs, $k=A/\mathfrak m$ le corps résiduel de $A$~;
on suppose $A$ et $B$ munis respectivement des topologies
$\mathfrak m$-préadique et $\mathfrak n$-préadique. Soit
$\varphi:A\to B$ un homomorphisme local, et posons $B_0=B\otimes_A k$.
Les propriétés suivantes sont équivalentes~:
\begin{enumerate}
  \item[a)] $B$ est une $A$-algèbre \emph{formellement lisse}.

\oldpage[0\textsubscript{IV-1}]{105}

  \item[b)] $B$ est un $A$-module \emph{plat} et
  $B_0=B/\mathfrak mB$ (muni de la topologie quotient) est une $k$-algèbre
  \emph{formellement lisse}.
\end{enumerate}
\end{theorem}

\begin{proof}
La démonstration se fait en plusieurs étapes.

\medskip
\noindent\phantomsection\label{0.19.7.1.1-fr}\textbf{(19.7.1.1)}\enspace
Démontrons d'abord que $b)$ entraîne $a)$~; nous allons appliquer le critère
(19.4.7), avec $\mathfrak J=\mathfrak m$~; en vertu de la seconde hypothèse
dans $b)$, tout revient à montrer que $B$ est un $A$-module
\emph{formellement projectif}. L'hypothèse entraîne que pour tout $h>0$,
$B/\mathfrak m^hB$ est un $(A/\mathfrak m^h)$-module plat
($\mathbf 0_{\mathrm{III}}$, 10.2.1)~; comme les $\mathfrak m^h$ forment un
système fondamental de voisinages de 0 dans $A$ et que
$(B/\mathfrak m^hB)\otimes_{A/\mathfrak m^h}k=B_0$, on peut remplacer $A$ et
$B$ par $A/\mathfrak m^h$ et $B/\mathfrak m^hB$ respectivement, et par suite
supposer $A$ \emph{artinien} (donc discret). Comme $B_0$ est une $k$-algèbre
formellement lisse, c'est un anneau \emph{régulier} (19.6.5)~; soit
$(x_i^0)_{1\leq i\leq n}$ un système régulier de paramètres pour $B_0$
(17.1.6), et pour tout $i$, soit $x_i\in B$ tel que $x_i^0$ soit son image
dans $B_0=B/\mathfrak mB$~; comme les $x_i^0$ engendrent l'idéal maximal
$\mathfrak n_0=\mathfrak n/\mathfrak mB$ de $B_0$, les idéaux
\[
  \mathfrak Q_m=\sum_{i=1}^n(x_i^0)^mB_0 \qquad (m>0)
\]
forment un système fondamental de voisinages de 0 dans $B_0$, car
$\mathfrak Q_m$ est évidemment contenu dans $\mathfrak n_0^m$, et d'autre
part contient $\mathfrak n_0^{mn}$. Posons
$\mathfrak J_m=\sum_{i=1}^n x_i^mB$ pour tout $m>0$~; il est clair que
$\mathfrak n=\mathfrak J_1+\mathfrak mB$~; comme il existe un $h>0$ tel que
$\mathfrak m^h=0$, on a
$\mathfrak n^m\subset\mathfrak J_1^{m-h}\subset\mathfrak n^{m-h}$ pour
$m>h$, et comme on a vu que $\mathfrak J_m\supset\mathfrak J_1^{mn}$, on voit
que les $\mathfrak J_m$ forment un système fondamental de voisinages de 0 dans
$B$. Tout revient par suite à prouver que les $B/\mathfrak J_m$ sont des
$A$-modules \emph{libres}, et il revient au même de voir que ce sont des
$A$-modules \emph{plats} ($\mathbf 0_{\mathrm{III}}$, 10.1.3). Or,
l'hypothèse que $(x_i^0)$ est une suite $B_0$-régulière d'éléments de l'idéal
maximal de $B_0$ entraîne la même propriété pour la suite des $(x_i^0)^m$
($1\leq i\leq n$) pour tout $m>0$ (15.1.20)~; la conclusion résulte donc de
(15.1.16, $b)$ et $c)$).
\end{proof}

\begin{lemma}[19.7.1.2]
\label{0.19.7.1.2-fr}
Soient $A$ un anneau topologique, $B$, $C$ deux $A$-algèbres topologiques qui
sont des anneaux locaux noethériens. On suppose en outre que $C$ est complet
et que le corps résiduel $B/\mathfrak m$ de $B$ est un $A$-module de type fini.
Soit $E$ le produit tensoriel complété $B\widehat\otimes_A C$. Alors~:
\begin{enumerate}
  \item[(i)] $E$ est un anneau semi-local noethérien complet.
  \item[(ii)] L'idéal $\mathfrak mE$ est contenu dans le radical de $E$, et
  pour tout $h>0$, $E/\mathfrak m^hE$ est isomorphe à
  $(B/\mathfrak m^h)\otimes_A C$.
  \item[(iii)] Si $C$ est un $A$-module plat, $E$ est un $B$-module plat.
\end{enumerate}
\end{lemma}

\begin{proof}
Par définition, $E$ est le séparé complété du produit tensoriel $B\otimes_A C$
pour la topologie définie par les idéaux
$\operatorname{Im}(\mathfrak m^i\otimes_A C)+
\operatorname{Im}(B\otimes_A\mathfrak n'{}^j)$
($\mathbf 0_{\mathrm I}$, 7.7.5). Si l'on pose
\[
  \mathfrak r=\operatorname{Im}(\mathfrak m\otimes_A C)+
  \operatorname{Im}(B\otimes_A\mathfrak n'),
\]
on a
$\mathfrak r^{2i}\subset
\operatorname{Im}(\mathfrak m^i\otimes_A C)+
\operatorname{Im}(B\otimes_A\mathfrak n'{}^i)\subset\mathfrak r^i$, donc $E$
est aussi le séparé complété de $B\otimes_A C$ pour la topologie
$\mathfrak r$-préadique. Par hypothèse,
$(B\otimes_A C)/\mathfrak r=(B/\mathfrak m)\otimes_A(C/\mathfrak n')$ est un
$(C/\mathfrak n')$-module de type fini, donc un anneau artinien~; en outre,
$\mathfrak r/\mathfrak r^2$, étant un quotient de
\[
  ((\mathfrak m/\mathfrak m^2)\otimes_A C)
  \oplus(B\otimes_A(\mathfrak n'/\mathfrak n'{}^2)),
\]
est un $(B\otimes_A C)$-module de type fini~; appliquant
($\mathbf 0_{\mathrm I}$, 7.2.11), on voit donc que $E$ est un anneau
\emph{noethérien}~; en outre, $E/\mathfrak rE$, isomorphe à
$(B\otimes_A C)/\mathfrak r$, étant artinien, $E$ est \emph{semi-local}.
Notons maintenant que $E$, qui est isomorphe à
\[
  \varprojlim_{i,j}((B/\mathfrak m^i)\otimes_A(C/\mathfrak n'{}^j)),
\]
est aussi isomorphe, par le théorème de la double limite projective, à
\[
  \varprojlim_i\left(\varprojlim_j
  ((B/\mathfrak m^i)\otimes_A(C/\mathfrak n'{}^j))\right)~;
\]

\oldpage[0\textsubscript{IV-1}]{106}

mais $\varprojlim_j((B/\mathfrak m^i)\otimes_A(C/\mathfrak n'{}^j))$ est le
séparé complété de $(B/\mathfrak m^i)\otimes_A C$, et comme $C$ est complet,
ce n'est autre que $(B/\mathfrak m^i)\otimes_A C$ lui-même,
$B/\mathfrak m^i$ étant un $A$-module de type fini puisque
$\mathfrak m/\mathfrak m^i$ est un $(B/\mathfrak m)$-module de type fini
($\mathbf 0_{\mathrm I}$, 7.3.6). On a donc
\[
  E=\varprojlim_i((B/\mathfrak m^i)\otimes_A C).
\]
Pour tout entier $h>0$, $\mathfrak m^hE$, étant un idéal de $E$, est fermé
dans $E$ ($\mathbf 0_{\mathrm I}$, 7.3.5), donc complet, et d'autre part il est
évidemment dense dans
\[
  \varprojlim_i
  (\operatorname{Im}(\mathfrak m^h/\mathfrak m^{h+i})\otimes_A C),
\]
donc égal à cette dernière limite projective. En outre, tous les systèmes
projectifs considérés sont définis par des homomorphismes \emph{surjectifs}~;
il résulte donc de ($\mathbf 0_{\mathrm{III}}$, 13.2.2) que
$E/\mathfrak m^hE$ est isomorphe à
\[
  (B/\mathfrak m^h)\otimes_A C
  =(B\otimes_A C)/\operatorname{Im}(\mathfrak m^h\otimes_A C).
\]
En particulier, comme
$\operatorname{Im}(\mathfrak m\otimes_A C)\subset\mathfrak r$, cela montre
que $\mathfrak mE$ est contenu dans $\mathfrak rE$, donc dans le radical de
$E$. Enfin, l'hypothèse que $C$ est un $A$-module plat entraîne que
$(B/\mathfrak m^h)\otimes_A C=E/\mathfrak m^hE$ est un
$(B/\mathfrak m^h)$-module plat pour tout $h>0$~: comme $B$ et $E$ sont
noethériens et que $\mathfrak mE$ est contenu dans le radical de $E$, il
résulte de ($\mathbf 0_{\mathrm{III}}$, 10.2.2) que $E$ est un $B$-module
plat.
\end{proof}

\begin{lemma}[19.7.1.3]
\label{0.19.7.1.3-fr}
Soient $A$ un anneau local noethérien, $\mathfrak m$ son idéal maximal, $k$ son
corps résiduel, $B_0$ une $k$-algèbre~; on suppose que $B_0$ est un anneau
local noethérien, complet et régulier. Alors il existe une $A$-algèbre
topologique $B$ qui est un anneau local noethérien complet, un $A$-module plat,
et tel que $B_0$ soit $k$-isomorphe à
$B\otimes_A k=B/\mathfrak mB$.
\end{lemma}

\begin{proof}
Comme $\widehat A$ est plat sur $A$ et a même corps résiduel, on peut se borner
au cas où $A$ est \emph{complet}.

Soit $K$ le corps résiduel de $B_0$, et distinguons deux cas~:

\medskip
\noindent I) $K$ est une extension \emph{séparable} de $k$. En vertu de
(19.6.4), $B_0$ est $k$-isomorphe à un anneau de séries formelles
$K[[T_1,\ldots,T_n]]$. Lorsque $B_0=K$, le lemme a déjà été démontré
($\mathbf 0_{\mathrm{III}}$, 10.3.1)~; soit $C$ un anneau local noethérien
complet qui est un $A$-module plat et tel que $C\otimes_A k$ soit isomorphe à
$K$. Pour $n\geq1$, il suffit de prendre (avec la notation précédente)
$B=C[[T_1,\ldots,T_n]]$~; on sait en effet (Bourbaki,
\emph{Alg. comm.}, chap.~III, §~3, n\textsuperscript{o}~4, cor.~3 du th.~1)
que $B$ est un $C$-module plat, donc aussi un $A$-module plat, et d'autre part,
il est immédiat que $C[[T_1,\ldots,T_n]]\otimes_A k$ est isomorphe à
$(C/\mathfrak mC)[[T_1,\ldots,T_n]]=B_0$.

\medskip
\noindent II) $K$ est de \emph{caractéristique} $p>0$, et par suite il en est
de même de $k$. Notons $P$ le corps premier $\mathbf F_p$, et $W(P)$ l'anneau
local complet des nombres $p$-adiques $\mathbf Z_p$, qui est un anneau de
valuation discrète (donc \emph{régulier}), et a $P$ pour corps résiduel.
Montrons d'abord qu'il existe un homomorphisme continu d'anneaux $W(P)\to A$
faisant donc de $A$ une $W(P)$-\emph{algèbre topologique}. En effet, si
$j:\mathbf Z\to A$ est l'homomorphisme canonique, on a
$j(p\mathbf Z)\subset\mathfrak m$ par hypothèse, d'où
$j^{-1}(\mathfrak m)=p\mathbf Z$, et par suite $j$ se factorise en
$\mathbf Z\to\mathbf Z_{p\mathbf Z}\xrightarrow{f}A$, où $f$ est un
homomorphisme local, donc continu, qui (puisque $A$ est complet) se prolonge
par continuité en l'homomorphisme $W(P)\to A$ cherché.

Comme $k$ est extension séparable de $P$, le cas I) montre qu'il y a un
homomorphisme local $W(P)\to W(k)$, où $W(k)$ est un anneau local noethérien
complet et un $W(P)$-module plat, tel que $W(k)\otimes_{W(P)}P$ soit isomorphe
à $k$. D'ailleurs, comme l'uniformisante $p$ de $W(P)$ est un élément
$W(k)$-régulier par platitude ($\mathbf 0_{\mathrm I}$, 6.3.4) et comme

\oldpage[0\textsubscript{IV-1}]{107}

$W(k)/pW(k)=k$, $pW(k)$ est l'idéal maximal de $W(k)$, ce qui entraîne que ce
dernier anneau est un \emph{anneau de valuation discrète complet} (Bourbaki,
\emph{Alg. comm.}, chap.~VI, §~3, n\textsuperscript{o}~5, prop.~9). Par
(19.7.1.1) on voit en outre (puisque $k$ est séparable sur $P$, donc une
$P$-algèbre formellement lisse (19.6.1)) que $W(k)$ est une $W(P)$-algèbre
\emph{formellement lisse}. Le $W(P)$-homomorphisme continu $W(k)\to k$ se
factorise donc en $W(k)\to A\to k$ (19.3.11), ce qui permet de considérer $A$
comme une $W(k)$-algèbre topologique. Appliquant maintenant le cas I) à $B_0$
considéré comme $P$-algèbre et à $W(P)$, on voit qu'il existe une
$W(P)$-algèbre $B_P$ qui est un anneau local noethérien complet, un
$W(P)$-module plat, et telle que $B_P\otimes_{W(P)}P$ soit $P$-isomorphe à
$B_0$. Utilisant de nouveau le fait que $W(k)$ est une $W(P)$-algèbre
formellement lisse, on voit par (19.3.11) que l'homomorphisme composé
$W(k)\to k\to B_0$ se factorise en $W(k)\to B_P\to B_0$~; en outre, comme
$k=W(k)/pW(k)$, on a
$B_P\otimes_{W(k)}k=B_P/pB_P=B_P\otimes_{W(P)}P=B_0$. Montrons que $B_P$ est
un $W(k)$-module \emph{plat}~; comme $W(k)$ est un anneau de valuation discrète
dont $p$ est l'uniformisante, il suffit de vérifier que $B_P$ est un
$W(k)$-module \emph{sans torsion} ($\mathbf 0_{\mathrm I}$, 6.3.4), ou encore
que $p$ est un élément $B_P$-régulier, ce qui résulte de ce que $B_P$ est un
$W(P)$-module plat ($\mathbf 0_{\mathrm I}$, 6.3.4).

Posons maintenant
\[
  B=B_P\widehat\otimes_{W(k)}A
\]
et notons que le corps résiduel de $A$ étant égal à celui de $W(k)$, est
\emph{a fortiori} un $W(k)$-module de type fini. Il résulte donc tout d'abord
de (19.7.1.2) que $B$ est un anneau semi-local noethérien complet,
$\mathfrak mB$ étant contenu dans le radical de $B$~; en outre
$B/\mathfrak mB$ est $k$-isomorphe à
$B_P\otimes_{W(k)}k=B_0$, donc $B$ est en fait un anneau local. Comme $B_P$
est un $W(k)$-module plat, (19.7.1.2) montre enfin que $B$ est un $A$-module
plat. C.Q.F.D.
\end{proof}

\begin{lemma}[19.7.1.4]
\label{0.19.7.1.4-fr}
Soient $A$ un anneau, $\mathfrak J$ un idéal de $A$, $M$, $N$ deux $A$-modules
séparés pour la topologie $\mathfrak J$-préadique. On suppose en outre que $N$
est complet pour la topologie $\mathfrak J$-préadique et que $M$ soit un
$A$-module plat. Soit $u:N\to M$ un $A$-homomorphisme~; si
$u\otimes1:N\otimes_A(A/\mathfrak J)\to M\otimes_A(A/\mathfrak J)$ est
bijectif, alors $u$ est bijectif.
\end{lemma}

\begin{proof}
Les modules gradués associés étant pris relatifs aux filtrations
$\mathfrak J$-préadiques, il résulte des hypothèses sur $M$ et $N$ relativement
aux topologies $\mathfrak J$-préadiques qu'il suffit de prouver que
$\gr(u):\gr_\bullet(N)\to\gr_\bullet(M)$ est \emph{bijectif} (Bourbaki,
\emph{Alg. comm.}, chap.~III, §~2, n\textsuperscript{o}~8, cor.~3 du th.~1).
Or, on a un diagramme commutatif
\[
\xymatrix{
  \gr_0(N)\otimes_{A_0}\gr_\bullet(A)
    \ar[r]^-{\gr_0(u)\otimes1}\ar[d]_{\varphi_N} &
  \gr_0(M)\otimes_{A_0}\gr_\bullet(A)\ar[d]^{\varphi_M} \\
  \gr_\bullet(N)\ar[r]_-{\gr(u)} & \gr_\bullet(M)
}
\]

\oldpage[0\textsubscript{IV-1}]{108}

où $A_0=A/\mathfrak J=\gr_0(A)$, et $\varphi_M$ et $\varphi_N$ sont les
applications canoniques ($\mathbf 0_{\mathrm{III}}$, 10.1.1.2). Par
hypothèse, $\gr_0(u)$ est bijectif ainsi que $\varphi_M$
($\mathbf 0_{\mathrm{III}}$, 10.2.1), et $\varphi_N$ est surjectif~; on en
déduit d'abord que $\gr_0(u)\otimes1$ est bijectif, puis que $\varphi_N$ est
injectif, donc bijectif, et enfin que $\gr(u)$ est bijectif.
\end{proof}

\begin{lemma}[19.7.1.5]
\label{0.19.7.1.5-fr}
Soient $A$ un anneau noethérien, $\mathfrak J$ un idéal de $A$, $B$, $B'$ deux
$A$-algèbres qui sont des anneaux locaux noethériens, les homomorphismes
$A\to B$, $A\to B'$ étant continus pour la topologie $\mathfrak J$-préadique
sur $A$. On suppose que~:
\begin{enumerate}
  \item[$1^{\mathrm o}$] $B$ et $B'$ sont complets pour les topologies
  $\mathfrak J$-préadiques~;
  \item[$2^{\mathrm o}$] $B$ est une $A$-algèbre formellement lisse~;
  \item[$3^{\mathrm o}$] $B'$ est un $A$-module plat.
\end{enumerate}
Posons $A_0=A/\mathfrak J$, et soit
$u_0:B\otimes_A A_0\to B'\otimes_A A_0$ un $A_0$-isomorphisme~; alors il
existe un $A$-isomorphisme $u:B\to B'$ tel que $u_0=u\otimes1$ (ce qui
entraîne que $B'$ est une $A$-algèbre formellement lisse et $B$ un $A$-module
plat).
\end{lemma}

\begin{proof}
Posons $B_0=B\otimes_A A_0$, $B'_0=B'\otimes_A A_0$. Notons que si
$\mathfrak m$ et $\mathfrak m'$ sont les idéaux maximaux de $B$ et $B'$, les
topologies $\mathfrak J$-préadiques sur $B$ et $B'$ sont séparées puisque
$\mathfrak JB\subset\mathfrak m$ et $\mathfrak JB'\subset\mathfrak m'$~; en
outre, comme $\mathfrak JB'$ est fermé dans $B'$ pour la topologie
$\mathfrak J$-préadique, l'homomorphisme composé
$B\to B_0\xrightarrow{u_0}B'_0$, qui est continu pour les topologies
$\mathfrak J$-préadiques, se factorise en $B\xrightarrow{u}B'\to B'_0$, où
$u$ est un $A$-homomorphisme continu (19.3.10). On a évidemment
$u_0=u\otimes1$, et l'hypothèse que $u_0$ est bijectif entraîne qu'il en est de
même de $u$ en vertu de (19.7.1.4).
\end{proof}

\medskip
\noindent\textbf{(19.7.1.6)}\enspace
\emph{Fin de la démonstration.} --- Pour achever de prouver (19.7.1), il faut
montrer que $a)$ entraîne $b)$~; on sait déjà que $a)$ entraîne que $B_0$ est
une $k$-algèbre formellement lisse (19.3.5, (iii)), donc tout revient à prouver
que $B$ est un $A$-module plat. Il revient au même d'établir que $\widehat B$
est un $\widehat A$-module plat (Bourbaki, \emph{Alg. comm.}, chap.~III, §~5,
n\textsuperscript{o}~4, prop.~4), et l'on sait que $\widehat B$ est une
$\widehat A$-algèbre formellement lisse (19.3.6)~; on peut donc se borner au
cas où $A$ et $B$ sont complets. Comme $B_0$ est une $k$-algèbre formellement
lisse, c'est un anneau régulier (19.6.5) et complet
($\mathbf 0_{\mathrm I}$, 6.3.5)~; appliquant (19.7.1.3), on voit qu'il existe
une $A$-algèbre $B'$ qui est un anneau local noethérien complet et un
$A$-module plat, un homomorphisme local $A\to B'$ et un $k$-isomorphisme
$B_0\simeq B'_0=B'\otimes_A k$. Il suffit alors d'appliquer (19.7.1.5) en
prenant pour $\mathfrak J$ l'idéal maximal de $A$, pour obtenir que $B$ est
$A$-isomorphe à $B'$, donc est un $A$-module plat. C.Q.F.D.

\begin{theorem}[19.7.2]
\label{0.19.7.2-fr}
Soient $A$ un anneau local noethérien, $\mathfrak J$ un idéal contenu dans
l'idéal maximal de $A$, $A_0=A/\mathfrak J$, $B_0$ un anneau local noethérien
complet, $A_0\to B_0$ un homomorphisme local faisant de $B_0$ une
$A_0$-algèbre formellement lisse. Alors il existe un anneau local noethérien
complet $B$, un homomorphisme local $A\to B$ faisant de $B$ un $A$-module
plat, et un $A_0$-isomorphisme
$u:B\otimes_A A_0\xrightarrow{\sim}B_0$. Si $(B',u')$ est un couple
satisfaisant aux mêmes conditions que $(B,u)$, il existe un $A$-isomorphisme
$v:B\xrightarrow{\sim}B'$ rendant commutatif le diagramme
\[
\xymatrix{
  B\otimes_A A_0\ar[rr]^{v\otimes1}_{\sim}\ar[dr]^{\sim}_u &&
  B'\otimes_A A_0\ar[dl]^{u'}_{\sim} \\
  & B_0
}
\]
\end{theorem}

\oldpage[0\textsubscript{IV-1}]{109}

\begin{proof}
Soit $\mathfrak m$ l'idéal maximal de $A$, de sorte que
$\mathfrak m_0=\mathfrak m/\mathfrak J$ est l'idéal maximal de $A_0$, $A$ et
$A_0$ ayant le même corps résiduel $k$. Posons
$B_{00}=B_0\otimes_{A_0}k$~; comme $B_{00}$ est une $k$-algèbre formellement
lisse (19.3.5, (iii)), c'est un anneau local régulier (19.6.5)~; en appliquant
(19.7.1.3), on voit qu'il existe une $A$-algèbre topologique $B$ qui est un
anneau local noethérien complet, un $A$-module plat, et pour laquelle on a un
$k$-isomorphisme
$u_{00}:B\otimes_A k\xrightarrow{\sim}B_{00}$. Notons qu'en vertu de (19.7.1),
$B$ est une $A$-algèbre formellement lisse, donc
$B\otimes_A A_0=B/\mathfrak JB$ est une $A_0$-algèbre formellement lisse
(19.3.5, (iii)) et un anneau local noethérien complet~; en outre, $B_0$ est un
$A_0$-module plat en vertu de l'hypothèse et de (19.7.1)~; comme on a un
$k$-isomorphisme
\[
  u_{00}:B\otimes_A k=(B\otimes_A A_0)\otimes_{A_0}k
  \xrightarrow{\sim}B_0\otimes_{A_0}k=B_{00},
\]
on déduit de (19.7.1.5), appliqué en y remplaçant $A$ par $A_0$ et
$\mathfrak J$ par $\mathfrak m_0$, qu'il existe un $A_0$-isomorphisme
$u:B\otimes_A A_0\xrightarrow{\sim}B_0$ tel que $u_{00}=u\otimes1$. Quant à
l'assertion d'unicité, notons que les idéaux $\mathfrak J^hB$ (resp.
$\mathfrak J^hB'$) sont fermés dans $B$ (resp. $B'$)
($\mathbf 0_{\mathrm I}$, 7.3.5), donc $B$ et $B'$ sont séparés et complets
pour les topologies $\mathfrak J$-préadiques (Bourbaki, \emph{Top. gén.},
chap.~III, 3\textsuperscript{e} éd., §~3, n\textsuperscript{o}~5, cor.~2 de la
prop.~9)~; on a par hypothèse un $A_0$-isomorphisme
$v_0:B\otimes_A A_0\xrightarrow{\sim}B'\otimes_A A_0$ tel que
$u'\circ v_0=u$~; comme $B$ est une $A$-algèbre formellement lisse et $B'$ un
$A$-module plat, on peut appliquer (19.7.1.5), d'où l'existence du
$A$-isomorphisme $v$ répondant à la question.
\end{proof}

\begin{remarks}[19.7.3]
\label{0.19.7.3-fr}
\begin{enumerate}
  \item[(i)] On notera que l'assertion d'unicité dans (19.7.2) est encore
  valable si l'on suppose seulement que $B$ et $B'$ sont complets \emph{pour
  les topologies $\mathfrak J$-préadiques}. Nous ignorons si l'on peut
  améliorer de même l'assertion d'existence, autrement dit si l'on peut se
  dispenser de supposer l'anneau local $B_0$ \emph{complet} (pour sa topologie
  $\mathfrak n_0$-préadique, en désignant par $\mathfrak n_0$ son idéal
  maximal) en exigeant seulement que $B$ soit complet \emph{pour la topologie
  $\mathfrak J$-préadique}. Lorsque $A_0$ est complet pour la topologie
  $\mathfrak m$-préadique, on peut voir que ce problème revient au suivant~:
  si $B_0$ est un anneau local noethérien régulier (non nécessairement
  complet) contenant le corps premier $\mathbf F_p=\mathbf Z/p\mathbf Z$,
  existe-t-il pour tout $n>1$ une $(\mathbf Z/p^n\mathbf Z)$-algèbre plate $B$
  telle que $B/pB$ soit isomorphe à $B_0$~?
  \item[(ii)] On notera qu'en général, l'isomorphisme $v$ dont l'existence est
  affirmée dans (19.7.2) \emph{n'est pas unique} (cf.~(19.8.7)).
\end{enumerate}
\end{remarks}

\subsection{Algèbres de Cohen et $p$-anneaux de Cohen; application à la
structure des anneaux locaux complets}
\label{subsection:0.19.8-fr}

Les résultats de cette section sont des applications immédiates des théorèmes
de (19.7), mais méritent d'être explicités en raison de leur importance
pratique.

\begin{definition}[19.8.1]
\label{0.19.8.1-fr}
Soient $A$, $B$ deux anneaux locaux noethériens, $\mathfrak m$ l'idéal maximal
de $A$, $k=A/\mathfrak m$ son corps résiduel,
$\varphi:A\to B$ un homomorphisme local, faisant de $B$ une $A$-algèbre. On
dit que $B$ est une \emph{$A$-algèbre de Cohen} si elle vérifie les conditions
suivantes~:
\begin{enumerate}
  \item[(i)] $B$ est un anneau \emph{complet}.
  \item[(ii)] $B$ est un $A$-module \emph{plat}.
  \item[(iii)] $B\otimes_A k$ est un \emph{corps} (autrement dit,
  $\mathfrak mB$ est l'idéal maximal de $B$) qui est une \emph{extension
  séparable} de $k$.
\end{enumerate}
\end{definition}

\oldpage[0\textsubscript{IV-1}]{110}

\begin{theorem}[19.8.2]
\label{0.19.8.2-fr}
Soient $A$ un anneau local noethérien, $k$ son corps résiduel.
\begin{enumerate}
  \item[(i)] Si $B$ est une $A$-algèbre de Cohen, $B$ est une $A$-algèbre
  formellement lisse. Pour tout anneau local noethérien complet $C$, tout
  homomorphisme local $A\to C$ et tout idéal $\mathfrak J\neq C$ dans $C$,
  tout $A$-homomorphisme $B\to C/\mathfrak J$ se factorise donc en
  $B\xrightarrow{v}C\to C/\mathfrak J$, où $v$ est un $A$-homomorphisme
  (nécessairement local).
  \item[(ii)] Pour tout corps $K$, extension séparable de $k$, il existe une
  $A$-algèbre de Cohen $B$ telle que $B\otimes_A k$ soit $k$-isomorphe à $K$,
  et une telle $A$-algèbre est unique à isomorphisme près.
\end{enumerate}
\end{theorem}

\begin{proof}
Comme $K$ est une $k$-algèbre formellement lisse (19.6.1), l'assertion (i)
résulte de (19.7.1). Pour prouver (ii), on peut se borner au cas où $A$ est
\emph{complet}, car il revient au même de dire que $B$ est un $A$-module plat
ou un $\widehat A$-module plat ($\mathbf 0_{\mathrm{III}}$, 10.2.3), on a
$\mathfrak mB=\mathfrak m\widehat AB$ et $k$ est le corps résiduel de
$\widehat A$. Il suffit alors d'appliquer (19.7.2) en prenant
$\mathfrak J=\mathfrak m$ et $B_0=K$ (et utilisant (19.6.1)).
\end{proof}

\begin{definition}[19.8.3]
\label{0.19.8.3-fr}
On appelle \emph{anneau local premier} un anneau local de la forme
$\mathbf Z_{p\mathbf Z}$, où $p\mathbf Z$ est un idéal premier de $\mathbf Z$.
On appelle \emph{anneau local complet premier} le complété d'un anneau local
premier.
\end{definition}

Les anneaux locaux premiers sont donc de deux sortes~:
\begin{enumerate}
  \item[$1^{\mathrm o}$] Ceux qui correspondent aux idéaux maximaux
  $p\mathbf Z$ où $p\neq0$ est un nombre premier~; $\mathbf Z_{p\mathbf Z}$
  est un anneau de valuation discrète, dont le complété est \emph{l'anneau des
  entiers $p$-adiques} noté d'ordinaire $\mathbf Z_p$\footnote{Cette notation,
  actuellement universellement utilisée, est dans ce cas en conflit avec la
  notation $A_f$ adoptée dans ($\mathbf 0_{\mathrm I}$, 1.2.3)~: avec
  $A=\mathbf Z$ et $f=p$, $A_f$ signifie en effet l'anneau des nombres
  rationnels de la forme $k/p^n$ ($k\in\mathbf Z$, $n$ entier $\geq0$)~; nous
  éviterons toujours d'utiliser la notation $\mathbf Z_p$ pour désigner ce
  dernier anneau.}.
  \item[$2^{\mathrm o}$] Pour l'idéal premier $p\mathbf Z=(0)$,
  $\mathbf Z_{p\mathbf Z}$ est le corps des nombres rationnels $\mathbf Q$,
  identique à son complété (la topologie étant naturellement la topologie
  d'anneau local noethérien, donc ici la topologie discrète).
\end{enumerate}

La terminologie de (19.8.3), analogue à celle des «~corps premiers~», se
justifie de la même manière~: pour tout anneau local $A$, considérons
l'homomorphisme canonique $\varphi:\mathbf Z\to A$, et soit $p\mathbf Z$
l'image réciproque par cet homomorphisme de l'idéal maximal $\mathfrak m$ de
$A$~; $p\mathbf Z$ est un idéal premier de $\mathbf Z$ et l'homomorphisme
précédent se factorise donc en
$\mathbf Z\to\mathbf Z_{p\mathbf Z}\xrightarrow{\psi}A$~; d'ailleurs, comme
$\varphi$ est l'\emph{unique} homomorphisme de $\mathbf Z$ dans $A$, $p$ et
$\psi$ sont déterminés de façon \emph{unique}. Autrement dit, pour tout anneau
local $A$, il y a un \emph{unique homomorphisme} $\psi:P\to A$, où $P$ est un
anneau local premier~; si en outre $A$ est séparé et complet, on peut prolonger
par complétion cet homomorphisme, et il y a donc un \emph{unique
homomorphisme} $\widehat\psi:\widehat P\to A$, où $\widehat P$ est un anneau
local complet premier. D'ailleurs, par passage aux quotients, $\psi$ donne un
homomorphisme du corps résiduel $\mathbf F_p$ si $p>0$ (resp. $\mathbf Q$ si
$p=0$) dans le corps résiduel $k$ de $A$, et $p$ est donc la
\emph{caractéristique} de $k$.

Si l'on prend en particulier pour $A$ un anneau local premier (resp. complet
premier), on voit qu'il n'existe dans un tel anneau qu'\emph{un seul
endomorphisme}, savoir l'identité.
